% \documentclass[lettersize,10pt,journal,compsoc]{IEEEtran}
\documentclass[lettersize,journal]{IEEEtran}
\usepackage{cite}
\usepackage{graphicx}
\usepackage{amsmath,amssymb,amsfonts}
\usepackage{array}          % 扩展列格式
\usepackage{booktabs}       % \toprule \midrule
\usepackage{multirow}       % 跨行单元格（如需）

% % subfig 若表格需要子表，可按需载入
% \usepackage[caption=false,font=footnotesize,labelfont=sf,textfont=sf]{subfig}
% \usepackage{caption} % put this after \usepackage{subfig}
% \captionsetup[figure]{justification=centering}
% \usepackage{tabularx}
% \usepackage{amssymb}
% \usepackage{pifont}
% \usepackage{adjustbox}
% \usepackage{enumitem}
% \usepackage{soul}

\usepackage[caption=false,font=footnotesize,labelfont=sf,textfont=sf]{subfig}
\usepackage{caption} % 必须在 subfig 后
\captionsetup[figure]{justification=centering}
\usepackage{tabularx}
\usepackage{pifont}
\usepackage{adjustbox}
\usepackage{enumitem}
\usepackage{soul}

\usepackage{amsthm}
\theoremstyle{definition}
\newtheorem{definition}{Definition}[section]

\usepackage{tikz}
\usepackage{url} 
\usetikzlibrary{positioning, arrows.meta, shapes, fit, calc}
\usepackage{hyperref}

\title{Mid-Training of Large Language Models: A Survey}

% \author{Mo \\ 2025-09}
\author{Kaixiang Mo, Yuxin Shi, Weiwei Weng, Zhiqiang Zhou, Shuman Liu, Haibo Zhang$^\dagger$, Anxiang Zeng$^\dagger$
        % <-this % stops a space
\thanks{All authors are with Shopee. E-mails: \{kaixiang.mo, yuxin.shi, weiwei.weng, zhiqiang.zhou, shuman.liu, peter.wu\}@shopee.com, zeng0118@ntu.edu.sg}% <-this % stops a space
\thanks{$^\dagger$Corresponding author: peter.wu@shopee.com, zeng0118@ntu.edu.sg}
\thanks{This work has been submitted to the IEEE for possible publication.}}


% The paper headers
\markboth{Journal of \LaTeX\ Class Files,~Vol.~14, No.~8, August~2021}%
{Shell \MakeLowercase{\textit{et al.}}: A Sample Article Using IEEEtran.cls for IEEE Journals}

% \IEEEpubid{0000--0000/00\$00.00~\copyright~2021 IEEE}
% Remember, if you use this you must call \IEEEpubidadjcol in the second
% column for its text to clear the IEEEpubid mark.

\begin{document}

\maketitle
% —— IEEEtran BST control (must come **before** the first \cite) ——
% See corresponding @IEEEtranBSTCTL entry in main.bib
\bstctlcite{BSTcontrol}



\begin{abstract}
Large language models (LLMs) are typically developed through large-scale pre-training followed by task-specific fine-tuning. Recent advances highlight the importance of an intermediate mid-training stage, where models undergo multiple annealing-style phases that refine data quality, adapts optimization schedules, and extend context length. This stage mitigates diminishing returns from noisy tokens, stabilizes convergence, and expands model capability in late training. Its effectiveness can be explained through gradient noise scale, the information bottleneck, and curriculum learning, which together promote generalization and abstraction. Despite widespread use in state-of-the-art systems, there has been no prior survey of mid-training as a unified paradigm. 
% This paper provides the first comprehensive review of strategies across three domains—data refinement, learning rate scheduling, and context extension—while discussing open challenges and future directions for both research and industry. 
We introduce the first taxonomy of LLM mid-training spanning data distribution, learning-rate scheduling, and long-context extension. We distill practical insights, compile evaluation benchmarks, and report gains to enable structured comparisons across models. We also identify open challenges and propose avenues for future research and practice.

%Mid-Training is an important procedure to improve the performance and extend the context length of the pre-trained model. However, there is no comprehensive updated survey about mid-training. In this paper, we present a comprehensive survey on recent advances in mid-training for large language models. This survey is organized into three key aspects of annealing, data distribution, learning rate schedule \& theory for annealing, and long context extension. This survey provides an organized view of the literature on mid-training, which serves as a valuable resource for industry and academia. 
\end{abstract}
% \clearpage
% \tableofcontents
% \clearpage
\begin{figure*}
    \centering
    \captionsetup{justification=raggedright,singlelinecheck=false}
    \includegraphics[width=1\linewidth]{1_overview.pdf}
    \caption{Mid Training of recent models. Only selected models with detailed data mix are shown. In the long context part, Hybrid means a combination of different attention mechanisms, e.g., self-attention, SSM, lightning-attention, etc. INT is the abbreviation for interleaving global and local attentions. }
    \label{fig:data_lr_longcontext}
\end{figure*}

\section{Introduction}
\label{sec:intro}


In recent large language model (LLM) training, mid-training has emerged as a crucial continuation stage between general pre-training and task-specific supervised fine-tuning (SFT). Rather than relying solely on web-scale corpora, which provide broad but noisy supervision, state-of-the-art models increasingly adopt an annealing-style phase to refine optimization dynamics and sharpen data quality~\cite{dubey2024llama, hu2024minicpm, glorioso2024zamba}. This stage is motivated by both practical and theoretical considerations: lowering the learning rate mitigates gradient variance and stabilizes convergence near favorable minima, while shifting toward more curated or synthetic corpora~\cite{abdin2024phi3} enhances the marginal utility of each additional token. Without such adjustments, models often stagnate in capability development, showing limited gains in reasoning, coding, and long-context understanding despite substantial additional compute~\cite{bi2024deepseek, allal2025smollm2}. 

The effectiveness of mid-training can be theoretically motivated, as it shifts models from memorization toward abstraction by emphasizing structured reasoning, factuality, and instruction following. We summarize this phenomenon from three perspectives: (i) gradient noise scale, where high-quality data improves optimization by enhancing signal variance and mitigating overfitting plateaus~\cite{mccandlish2018empirical, wang2021eliminating, meng2023per}; (ii) the information bottleneck, which compresses noisy features while preserving predictive structures~\cite{izmailov2022feature}; and (iii) curriculum learning, where data distributions are gradually refined to reinforce complex reasoning~\cite{wang2025dump, jia2022sample}. Together, these perspectives provide a principled rationale for mid-training as a strategy to improve generalization and efficiency in large-scale training.

Despite its growing adoption and effectiveness, there has been no comprehensive survey that conceptualizes mid-training as a coherent paradigm. Existing studies have explored isolated aspects (e.g., curriculum-style data annealing~\cite{du2024survey,wu2025data}, adaptive learning-rate schedules~\cite{jin2023rethinking}, or context-length extension~\cite{huang2023advancing, liu2025comprehensive, pawar2024and}) but lack a comprehensive framework. To bridge this gap, we present an integrated review of mid-training strategies across three interconnected domains: data distribution, optimization scheduling, and context extension. By framing these elements as mutually reinforcing, we highlight mid-training as a distinct and coherent stage in the LLM development pipeline, rather than an ad-hoc collection of heuristics.

A central aspect of mid-training is the refinement of the data mixture after large-scale pre-training. At this point, models have acquired broad linguistic and semantic features from trillions of web-scale tokens, making the marginal value of additional noisy data limited. To improve efficiency, the distribution is shifted toward curated or synthetic corpora emphasizing reasoning, coding, STEM, multilinguality, and instruction following. Empirical evidence shows that such targeted data can foster compositional skills (e.g., chain-of-thought reasoning, math problem solving, and multi-step planning) that are underrepresented in general crawls~\cite{abdin2024phi, allal2025smollm2}. Common practice includes down-sampling low-quality tokens and up-sampling knowledge-dense corpora, often through iterative learning rate annealing passes~\cite{NEURIPS2024_19e4ea30, penedo2024fineweb, Su2024NemotronCCTC}. 

Another key component of mid-training is adapting the learning rate schedule. After aggressive pre-training, models approach promising optima where large step sizes risk instability or divergence. Reducing the learning rate with smoother decay stabilizes convergence and suppresses gradient noise~\cite{you1908does, gilmer2022loss, kalra2024warmup}. Existing strategies include linear or cosine decay, multi-stage schedulers, and adaptive schemes tailored to long runs~\cite{goyal2017accurate, loshchilov2016sgdr, bi2024deepseek, hu2024minicpm, shen2024power, ibrahim2024simple}. These methods improve sample efficiency by enabling finer assimilation of high-quality tokens. Yet scheduler design remains largely empirical: the optimal duration and shape of decay vary across model sizes, architectures, and optimizers~\cite{kaplan2020scaling, hoffmann2022training, defazio2023optimal, bergsma2025straight, li2025predictable}. Overly steep decay may halt learning prematurely, while overly cautious schedules waste computation with limited gains.

Finally, mid-training is also a natural stage for extending context length beyond the 4K–8K limits of early pre-training. This extension supports downstream tasks that demand coherence across long documents, multi-file reasoning, or dialogue histories spanning tens of thousands of tokens. Current methods combine positional encoding remapping, such as Position Interpolation~\cite{chen2023extending}, NTK-aware interpolation~\cite{bloc97_ntk_rope_2023}, YaRN~\cite{peng2023yarn} or ABF~\cite{xiong2023effective} with curricula that gradually introduce long-form inputs. Studies show that extending context during mid-training enables models to capture long-range dependencies more effectively without restarting pre-training. 

Together, these elements form the backbone of mid-training: without refined data, the model lacks specialization; without proper scheduling, it risks instability or underutilization of valuable tokens; and without context extension, it remains constrained in scope for downstream applications. Their interdependence suggests that effective mid-training cannot be understood by analyzing each factor in isolation, but rather through their combined impact on model performance and efficiency. This holistic view motivates the need for systematic discussion of mid-training as a distinct stage in the LLM development pipeline.

\textbf{Our contributions}. This paper contributes to academia and industry in the following ways:
\begin{itemize}[leftmargin=10pt]
    \item We propose a taxonomy of LLM mid-training based on three domains (data distribution, learning rate scheduling and long-context extension), and summarize various approaches in each domain. To the best of our knowledge, it is the first such taxonomy on LLM mid-training.
    \item We summarize valuable insights in each domain, providing readers with a convenient reference for future use.
    \item We discuss the common evaluation benchmarks and collect the reported gains from different models, thereby, providing a structured comparison of how mid-training contribute to LLM training.
    \item We outline promising future research direction in LLM mid-training and propose potential ways forward.
\end{itemize}

\section{Possible Theory Behind Mid-Training}
The core rationale for mid-training is to shift the model's learning dynamics from memorization to abstraction. As the model approaches the capacity limits of broad generalization from large-scale general or medium-quality data, high-quality samples enable further gains by sharpening attention to structured reasoning, factuality, and instruction following. Below, we outline several theoretical foundations for this approach.


% \begin{table*}[h]
% \centering
% \scriptsize
% \renewcommand{\arraystretch}{1.2}
% \newcolumntype{L}[1]{>{\raggedright\arraybackslash}p{#1}}
% \caption{Consolidated overview of state-of-the-art models that disclose their mid-training datasets.}
% \label{tab:open_sourced_ds}

% \begin{adjustbox}{width=\textwidth}
% \begin{tabularx}{\textwidth}{@{} l L{3cm} L{2.5cm} L{3cm} L{3cm} L{3cm} @{}}
% \toprule
% \textbf{Model} & \textbf{HQ Web Data} & \textbf{Code} & \textbf{Math} & \textbf{Instruction\& QA Style} & \textbf{Synthetic Data} \\
% \midrule
% \textbf{MiniCPM}\mbox{\cite{hu2024minicpm}} & CommonCrawl\_Chn; Dolma~\cite{soldaini2024dolma}; C4~\cite{C4dataset}; Pile~\cite{gao2020pile}; Wikipedia~\cite{wikidump}; Arxiv~\cite{clement2019usearxivdataset}; BaiduBaike & Stack~\cite{Kocetkov2022TheStack} & OpenWebMath~\cite{paster2023openwebmath} & SlimOrca~\cite{SlimOrca}; OssInstruct~\cite{wei2024magicoder}; UltraChat~\cite{ding2023enhancing}; EvolInstruct~\cite{surge2024openbezoar}; Stack Exchange QA & \textemdash \\
% % \midrule
% \textbf{Zamba-v1}\mbox{\cite{glorioso2024zamba}} & Pile~\cite{gao2020pile}; C4~\cite{C4dataset}; PeS2o~\cite{peS2o}; Arxiv; RefinedWeb~\cite{penedo2023refinedweb} & EvolInstructCode\cite{luo2023wizardcoder} & StackMathQA~\cite{stackmathqa2024} & OpenOrca~\cite{OpenOrca} & Cosmopedia~\cite{benallal2024cosmopedia} \\
% % \midrule
% \textbf{MAP-Neo}\mbox{\cite{Zhang2024MAPNeoHC}} & Matrix~\cite{Zhang2024MAPNeoHC}; US-PD-Books~\cite{US-PD-Books} & Matrix~\cite{Zhang2024MAPNeoHC}; Stack V2~\cite{lozhkov2024starcoder} & Matrix~\cite{Zhang2024MAPNeoHC}; AutoMathText~\cite{zhang2025autonomous}; OpenWebMath~\cite{paster2023openwebmath} & Matrix~\cite{Zhang2024MAPNeoHC}; BioInstrutQA~\cite{BioInstructQA}; SMolInstruct~\cite{yu2024llasmol}; Agent-FLAN~\cite{chen2024agent} & Matrix~\cite{Zhang2024MAPNeoHC}; Cosmopedia~\cite{benallal2024cosmopedia} \\
% % \midrule
% \textbf{OLMo-2}\mbox{\cite{olmo20242}} & 
% Dolmino-Mix-1124~\cite{olmo20242} (DCLM-Baseline~\cite{NEURIPS2024_19e4ea30}; PeS2o-Dolma1.7~\cite{soldaini2024dolma}; Wiki-Dolma1.7~\cite{soldaini2024dolma}) & Dolmino-Mix-1124~\cite{olmo20242} (CodeSearchNet-OWMfilter~\cite{husain2019codesearchnet}) & Dolmino-Mix-1124~\cite{olmo20242} (Metamath-OWMfilter~\cite{yu2023metamath}; GSM8K-trainsplit~\cite{cobbe2021gsm8k}) & Dolmino-Mix-1124~\cite{olmo20242} (Flan-Dolma1.7~\cite{soldaini2024dolma}; StackExchange-RedPajamaV1~\cite{together2023redpajama}) & Dolmino-Mix-1124~\cite{olmo20242} (TuleMath; Dolmino SynthMath; TinyGSM-MIND; MathCoder2-synthetic)\\
% % \midrule
% \textbf{SmolLM2}\mbox{\cite{allal2025smollm2}} & DCLM-Baseline~\cite{NEURIPS2024_19e4ea30}; FineWeb-Edu~\cite{penedo2024fineweb} & Stack-Edu~\cite{allal2025smollm2} & InfiWebMath-3+~\cite{han2024infimmwebmath40badvancingmultimodalpretraining}; FineMath 4+~\cite{allal2025smollm2}; OpenWebMath~\cite{paster2023openwebmath}; AugGSM8K~\cite{li2023mugglemath} & \textemdash & Cosmopedia v2~\cite{benallal2024smollmcorpus} \\
% \bottomrule
% \end{tabularx}
% \end{adjustbox}
% \end{table*}

\begin{table*}[h]
\centering
\scriptsize
\renewcommand{\arraystretch}{1.2}
\newcolumntype{L}[1]{>{\raggedright\arraybackslash}p{#1}}
\caption{Consolidated overview of state-of-the-art models that disclose their mid-training datasets.}
\label{tab:open_sourced_ds}

\begin{adjustbox}{width=\textwidth}
\begin{tabularx}{\textwidth}{@{} l L{3cm} L{2.5cm} L{3cm} L{3cm} L{3cm} @{}}
\toprule
\textbf{Model} & \textbf{HQ Web Data} & \textbf{Code} & \textbf{Math} & \textbf{Instruction\& QA Style} & \textbf{Synthetic Data} \\
\midrule
\textbf{MiniCPM}\mbox{\cite{hu2024minicpm}} & CommonCrawl\_Chn; Dolma~\mbox{\cite{soldaini2024dolma}}; C4~\mbox{\cite{C4dataset}}; Pile~\mbox{\cite{gao2020pile}}; Wikipedia~\mbox{\cite{wikidump}}; Arxiv~\mbox{\cite{clement2019usearxivdataset}}; BaiduBaike & Stack~\mbox{\cite{Kocetkov2022TheStack}} & OpenWebMath~\mbox{\cite{paster2023openwebmath}} & SlimOrca~\mbox{\cite{SlimOrca}}; OssInstruct~\mbox{\cite{wei2024magicoder}}; UltraChat~\mbox{\cite{ding2023enhancing}}; EvolInstruct~\mbox{\cite{surge2024openbezoar}}; Stack Exchange QA & \textemdash \\
% \midrule
\textbf{Zamba-v1}\mbox{\cite{glorioso2024zamba}} & Pile~\mbox{\cite{gao2020pile}}; C4~\mbox{\cite{C4dataset}}; PeS2o~\mbox{\cite{peS2o}}; Arxiv; RefinedWeb~\mbox{\cite{penedo2023refinedweb}} & EvolInstructCode\mbox{\cite{luo2023wizardcoder}} & StackMathQA~\mbox{\cite{stackmathqa2024}} & OpenOrca~\mbox{\cite{OpenOrca}} & Cosmopedia~\mbox{\cite{benallal2024cosmopedia}} \\
% \midrule
\textbf{MAP-Neo}\mbox{\cite{Zhang2024MAPNeoHC}} & Matrix~\mbox{\cite{Zhang2024MAPNeoHC}}; US-PD-Books~\mbox{\cite{US-PD-Books}} & Matrix~\mbox{\cite{Zhang2024MAPNeoHC}}; Stack V2~\mbox{\cite{lozhkov2024starcoder}} & Matrix~\mbox{\cite{Zhang2024MAPNeoHC}}; AutoMathText~\mbox{\cite{zhang2025autonomous}}; OpenWebMath~\mbox{\cite{paster2023openwebmath}} & Matrix~\mbox{\cite{Zhang2024MAPNeoHC}}; BioInstrutQA~\mbox{\cite{BioInstructQA}}; SMolInstruct~\mbox{\cite{yu2024llasmol}}; Agent-FLAN~\mbox{\cite{chen2024agent}} & Matrix~\mbox{\cite{Zhang2024MAPNeoHC}}; Cosmopedia~\mbox{\cite{benallal2024cosmopedia}} \\
% \midrule
\textbf{OLMo-2}\mbox{\cite{olmo20242}} & 
Dolmino-Mix-1124~\mbox{\cite{olmo20242}} (DCLM-Baseline~\mbox{\cite{NEURIPS2024_19e4ea30}}; PeS2o-Dolma1.7~\mbox{\cite{soldaini2024dolma}}; Wiki-Dolma1.7~\mbox{\cite{soldaini2024dolma}}) & Dolmino-Mix-1124~\mbox{\cite{olmo20242}} (CodeSearchNet-OWMfilter~\mbox{\cite{husain2019codesearchnet}}) & Dolmino-Mix-1124~\mbox{\cite{olmo20242}} (Metamath-OWMfilter~\mbox{\cite{yu2023metamath}}; GSM8K-trainsplit~\mbox{\cite{cobbe2021gsm8k}}) & Dolmino-Mix-1124~\mbox{\cite{olmo20242}} (Flan-Dolma1.7~\mbox{\cite{soldaini2024dolma}}; StackExchange-RedPajamaV1~\mbox{\cite{together2023redpajama}}) & Dolmino-Mix-1124~\mbox{\cite{olmo20242}} (TuleMath; Dolmino SynthMath; TinyGSM-MIND; MathCoder2-synthetic)\\
% \midrule
\textbf{SmolLM2}\mbox{\cite{allal2025smollm2}} & DCLM-Baseline~\mbox{\cite{NEURIPS2024_19e4ea30}}; FineWeb-Edu~\mbox{\cite{penedo2024fineweb}} & Stack-Edu~\mbox{\cite{allal2025smollm2}} & InfiWebMath-3+~\mbox{\cite{han2024infimmwebmath40badvancingmultimodalpretraining}}; FineMath 4+~\mbox{\cite{allal2025smollm2}}; OpenWebMath~\mbox{\cite{paster2023openwebmath}}; AugGSM8K~\mbox{\cite{li2023mugglemath}} & \textemdash & Cosmopedia v2~\mbox{\cite{benallal2024smollmcorpus}} \\
\bottomrule
\end{tabularx}
\end{adjustbox}
\end{table*}


\noindent\textbf{Gradient Noise Scale:} 
% Prior work shows that the gradient noise scale (GNS) reflects the amount of useful signal per update step~\cite{mccandlish2018empirical}. Intuitively, higher-quality data is expected to induce larger gradient variance across examples, which can manifest as a higher GNS, whereas redundant or noisy data may reduce gradient diversity. Higher GNS enables models to traverse sharper minima in the loss landscape and avoid overfitting to low-quality plateaus, consistent with findings that gradient noise can help models escape sharp minima and improve generalization~\cite{neelakantan2015adding, wang2021eliminating}. Empirically, this has been observed to increase the effectiveness of optimization during late training stages, where signal sparsity can otherwise stall convergence, and recent work further shows that improving signal-to-noise characteristics in the gradient can mitigate overfitting on noisy or low-quality data~\cite{meng2023per}.
Prior work shows that the gradient noise scale (GNS) reflects the amount of useful signal per update step~\cite{mccandlish2018empirical}. Higher-quality data tends to induce greater gradient variance, yielding a higher GNS, whereas redundant or noisy data reduces diversity. A larger GNS helps models escape sharp minima and avoid overfitting to low-quality plateaus~\cite{neelakantan2015adding, wang2021eliminating}. Empirically, this improves optimization in late training stages, where signal sparsity may stall convergence, and recent work shows that enhancing gradient signal-to-noise can further mitigate overfitting on noisy data~\cite{meng2023per}.

\noindent\textbf{Information Bottleneck:} 
% From the perspective of the information bottleneck theory, the model during mid-training compresses its internal representation by retaining only the most relevant features for downstream tasks. High-quality data provides clearer, less entropic supervision signals, which help the model identify semantically rich structures and discard spurious correlations learned earlier in training. This mechanism helps transform large-scale memorization into usable abstraction.
In the information bottleneck (IB) framework, representation learning in neural networks can be interpreted as a process of compressing internal states while retaining task-relevant information~\cite{tishby2015deep}. During the learning rate annealing phase of training, the model progressively reduces reliance on noisy or redundant features. The model's representations may be viewed as increasingly emphasizing those features most predictive for downstream objectives~\cite{alemi2016deep}. In this view, high-quality supervision signals provide clearer, lower-entropy guidance that facilitates the identification of semantically meaningful structures and the attenuation of spurious correlations~\cite{izmailov2022feature}. Recent extensions of the IB principle to supervised settings further highlight how such compression can transform large-scale memorization into more abstract, generalizable representations~\cite{weingarten2025supervised}.

\noindent\textbf{Curriculum Learning:} 
% Learning rate annealing naturally aligns with the principles of curriculum learning. Early pretraining exposes the model to diverse and noisy corpora for generalization. As training progresses, the data distribution can be gradually shifted toward more challenging and informative examples—exactly the role of learning rate annealing. This structured progression helps optimize learning efficiency and reinforces complex skills such as multi-step reasoning and code generation~\cite{nair2024curriculum}.
Learning rate annealing naturally aligns with the principles of curriculum learning. Early pretraining exposes the model to diverse and noisy corpora for generalization, after which the data distribution can be gradually shifted toward more challenging and informative examples—exactly the role of curriculum-style scheduling~\cite{nair2024curriculum, chaudhry2024data}. This structured progression helps optimize learning efficiency and has been shown to reinforce complex skills such as multi-step reasoning and code generation~\cite{wang2025dump, jia2022sample}.



\section{Data Distributions}
\label{sec:data}
% In this section, we study the question "why use high quality data" from three different theory perspectives, and then we will survey datasets and distributions used in the annealing of the state-of-the-art models. 
In this section, we will study the mid-training dataset distribution and quality in the state-of-the-art LLM approaches.


\subsection{Type of Commonly Used Mid-Training Data}
The mid-training stage in LLM pretraining represents a critical phase where data quality becomes increasingly prioritized to refine model capabilities, improve alignment, and optimize performance on downstream tasks. Across state-of-the-art LLMs, several distinct categories of data are commonly employed during mid-training. These can be broadly grouped into the following types:



\noindent\textbf{High-Quality Filtered Web Data:}
% Models such as MiniCPM, Hunyuan-Large, Llama3, SmolLM2, OLMo-2, and Phi4 extensively utilize curated subsets of web-scale corpora. 
Models~\cite{hu2024minicpm, sun2024hunyuan, Dubey2024TheL3, allal2025smollm2, olmo20242, abdin2024phi} extensively utilize curated subsets of web-scale corpora. 
These subsets are meticulously filtered based on rigorous criteria including content quality, topic diversity, low toxicity, educational value, and minimal redundancy. In contrast to datasets used during initial pretraining stages, these refined datasets often undergo comprehensive filtering processes involving heuristic-based rules or learned scoring mechanisms to ensure selection of the most informative and high-quality samples.
% Representative datasets include general web-scale corpora such as CommonCrawl, C4~\cite{C4dataset}, Wikipedia~\cite{wikidump}, Dolma~\cite{soldaini2024dolma}, RedPajama-Data-V2, Cultrax, RefinedWeb~\cite{penedo2023refinedweb} (600B tokens), SlimPajam and Matrix~\cite{Zhang2024MAPNeoHC} (4.69T tokens). Further specialized filtering has yielded higher-quality subsets such as FineWeb-Edu~\cite{penedo2024fineweb} (1.3T tokens), Ultra-FineWeb (1.12T tokens), DCLM-baseline~\cite{NEURIPS2024_19e4ea30} (3.8T tokens), peS2o~\cite{peS2o} ($\sim$45B tokens).
Representative datasets include CommonCrawl, C4~\cite{C4dataset}, Wikipedia~\cite{wikidump}, Dolma~\cite{soldaini2024dolma}, RedPajama-Data-V2~\cite{together2023redpajama}, Culturax~\cite{nguyen2023culturax}, RefinedWeb~\cite{penedo2023refinedweb}, SlimPajama~\cite{shen2023slimpajama} and Matrix~\cite{Zhang2024MAPNeoHC}. Further specialized filtering has yielded higher-quality subsets such as FineWeb-Edu~\cite{penedo2024fineweb}, Ultra-FineWeb, DCLM-baseline~\cite{NEURIPS2024_19e4ea30}, peS2o~\cite{peS2o}.

% Many models such as (MiniCPM, Hunyuan-Large, Llama3, SmolLM2, OLMo-2, Phi4) utilize a curated subset of web-scale corpora, where documents are filtered based on criteria such as quality, topical diversity, toxicity, educational value and low redundancy. Compared to early-stage pretraining data, these datasets are significantly more refined, often passing multiple heuristics or learned scoring mechanisms to identify high-quality samples.

% Related datasets: CommonCrawl, C4, Wikipedia, Dolma, RedPajama-Data-V2, Cultrax, Refined-Web, SlimPajam.

% Related filtered datasets: FineWeb-Edu, Ultra-FineWeb (1.12T), DCLM-baseline (3.71T tokens), peS2o (from Dolma 1.7, 58.6B), Wikipedia\&Wikibooks (from Dolma 1.7, 3.7B), 

\noindent\textbf{Code and Mathematical Content:}
% Code and mathematics data significantly enhance symbolic reasoning and program synthesis capabilities in models such as AFM, SmolLM2, MAP-Neo, MiniCPM, Minimax-m1, SmolLM2, and Qwen3.
Code and math data significantly enhance symbolic reasoning and program synthesis capabilities in mid-training stages~\cite{gunter2024apple, allal2025smollm2, Zhang2024MAPNeoHC, hu2024minicpm, chen2025minimax, yang2025qwen3, yang2025qwen3}.
Incorporating structured content from open-source repositories, curated coding benchmarks, and mathematical question-answer pairs or derivations enriches the model's ability to handle complex logic and technical reasoning.
% Representative datasets include Stack v1~\cite{Kocetkov2022TheStack} ($\sim$200B tokens), Stack v2~\cite{lozhkov2024starcoder} ($\sim$900B tokens). More specifically filtered datasets include StackMathQA, OpenWebMath~\cite{paster2023openwebmath} (12B tokens), InfiMM-WebMath (40B tokens), FineMath~\cite{hu2024minicpm} (54B tokens focusing on math and reasoning through classifer-based filtering), FineMath4+~\cite{hu2024minicpm} (10B tokens), FineMath3+~\cite{hu2024minicpm}
% (34B tokens), InfiWebMath~\cite{han2024infimmwebmath40badvancingmultimodalpretraining} (40B tokens), Infi-WebMath4+~\cite{hu2024minicpm} (8.5B tokens), Infi-WebMath3+~\cite{hu2024minicpm} (20.5B tokens), Stack-Edu~\cite{allal2025smollm2} (125B tokens, a filtered variant of StarCoder2Data), RefineCode~\cite{huang2024opencoder} (960B tokens), and EvolInstructCode~\cite{luo2023wizardcoder}.
Representative datasets include Stack v1~\cite{Kocetkov2022TheStack}, Stack v2~\cite{lozhkov2024starcoder}. More specifically filtered datasets include StackMathQA, OpenWebMath~\cite{paster2023openwebmath}, InfiMM-WebMath~\cite{han2024infimm}, FineMath~\cite{hu2024minicpm}, FineMath4+~\cite{hu2024minicpm}, FineMath3+~\cite{hu2024minicpm}, InfiWebMath~\cite{han2024infimmwebmath40badvancingmultimodalpretraining} , Infi-WebMath4+~\cite{hu2024minicpm}, Infi-WebMath3+~\cite{hu2024minicpm}, Stack-Edu~\cite{allal2025smollm2}, RefineCode~\cite{huang2024opencoder}, and EvolInstructCode~\cite{luo2023wizardcoder}.
% Code and math-heavy data play a prominent role in the mid-training stage for models like AFM, SmolLM2, MAP-Neo and X(MiniCPM, Minimax-m1, SmolLM2, Qwen3). The inclusion of such structured content improves the model's symbolic reasoning and program synthesis capabilities. Datasets often include open-source repositories, curated coding benchmarks, and mathematical question-answer pairs or derivations.

% Related datasets: Stack v1, Stack v2, StarCoderData, StarCoder2Data, 
% Related filtered dataset: OpenWebMath (12B tokens), InfiMM-WebMath (40B tokens), FineMath (using Llama-3.1-70B-Instruct), Stack-Edu (125B token, a filtered variant of StarCoder2Data filtered by Llama3-70B-Instruct), RefineCode (960B tokens, proposed by opencoder)

\noindent\textbf{Instruction-Tuned and QA-Style Data:}
% Instruction-following data is increasingly prevalent during the annealing stage, prominently utilized by models including Nemotron-4, MAP-Neo, MiniCPM, Hunyuan-Large, Minimax-m1, Pangu Pro MoE, OLMo-2, and Phi4. 
Instruction-following and QA-style data is increasingly prevalent during the mid-training~\cite{adler2024nemotron, hu2024minicpm, abdin2024phi, Zhang2024MAPNeoHC, sun2024hunyuan, olmo20242, chen2025minimax, tang2025pangu}. 
Such datasets typically comprise synthetic or curated question-answer pairs, instruction-response prompts, and alignment-tuning corpora. This type of data aims to enhance the model’s understanding of human intent, improve the consistency of the response, and strengthen the reasoning ability of the models~\cite{abdin2024phi}.
% Representative datasets include Ultrachat~\cite{ding2023enhancing} (1.5M multi-turn samples), EvolInstruct~\cite{surge2024openbezoar} (2K samples), OssInstruct~\cite{wei2024magicoder} (75K samples), ShareGPT4, StackExchangeQA, and FLAN~\cite{longpre2023flan} (15M samples), OpenOrca~\cite{OpenOrca} (4.2M samples), SMolInstruct~\cite{yu2024llasmol} (3M samples).
Representative datasets include Ultrachat~\cite{ding2023enhancing}, EvolInstruct~\cite{surge2024openbezoar}, OssInstruct~\cite{wei2024magicoder}, StackExchangeQA~\cite{olmo20242}, and FLAN~\cite{longpre2023flan}, OpenOrca~\cite{OpenOrca}, SMolInstruct~\cite{yu2024llasmol}.
% Instruction-following data is increasingly prevalent during mid-training, as seen in Nemotron-4, MAP-Neo and (MiniCPM, Hunyuan-Large, Minimax-m1, Pangu Pro MoE, OLMo-2, Phi4 ). These examples typically include synthetic or curated question-answer pairs, instruction-response prompts, and alignment-tuning corpora. The goal is to nudge the model closer to human intent understanding and response consistency.

% Related datasets: Ultrachat, EvolInstruct, OssInstruct, ShareGPT4, StackExchangeQA, FLAN (from Dolma 1.7, 17B)

\noindent\textbf{Synthetic Textbooks and Knowledge-Dense Data:}
Synthetic textbook-style and knowledge-dense data significantly enhance LLMs by providing high-quality educational content~\cite{abdin2024phi, allal2025smollm2, glorioso2024zamba, hu2024minicpm, sun2024hunyuan, gunter2024apple, olmo20242, yang2025qwen3}, particularly in low-resource domains such as math, coding, and multilingual contexts. These synthetic datasets are typically derived from knowledge graphs or generated using advanced LLMs by utilizing high-quality seeds sourced from multiple domains~\cite{abdin2024phi}. The synthetic data enrich the models with factual, explanatory, and pedagogically structured information. Such data strengthens the factual grounding and expands the world knowledge of models, improving their general reasoning and explanatory capabilities. 
% Models leveraging this data include Phi, SmolLM2, Zamba-v1, MiniCPM, Hunyuan-Large, AFM, OLMo-2, and Qwen3. 
% Models leveraging this data include~\cite{abdin2024phi, allal2025smollm2, glorioso2024zamba, hu2024minicpm, sun2024hunyuan, gunter2024apple, olmo20242, yang2025qwen3}. 
% Representative datasets include TuluMath (230M tokens), Dolmino SynthMath (28.7M tokens), TinyGSM-MIND (6.48B tokens), MathCoder2 Synthetic (3.87B tokens), Cosmopedia~\cite{benallal2024cosmopedia} (25B tokens), Cosmopedia v2~\cite{benallal2024smollmcorpus} (28B tokens), OpenHermes-2.5 (1M samples).
Representative datasets include TuluMath~\cite{olmo20242}, Dolmino SynthMath~\cite{olmo20242}, TinyGSM-MIND~\cite{olmo20242}, MathCoder2 Synthetic~\cite{olmo20242}, Cosmopedia~\cite{benallal2024cosmopedia}, Cosmopedia v2~\cite{benallal2024smollmcorpus}, OpenHermes-2.5.

\noindent\textbf{Long-context Data:}
% For models that support extended context lengths, such as Qwen3, MiniMax-M1, Hunyuan-Large and X (Qwen3), the annealing stage may incorporate long-context documents or long-form Q\&A designed to span tens or hundreds of thousands of tokens. 
For models that support extended context lengths~\cite{yang2025qwen3, chen2025minimax, sun2024hunyuan, abdin2024phi}, the mid-training stage may incorporate long-context documents or long-form Q\&A designed to span thousands of tokens. 
These long-context data are either filtered from existing training corpus~\cite{sun2024hunyuan, abdin2024phi, allal2025smollm2, gao2024train} or obtained from data synthesis~\cite{abdin2024phi}. These samples improve the model’s memory, coherence, and reasoning across document-scale inputs.
Related datasets used to extract long-context data: DCLM~\cite{NEURIPS2024_19e4ea30}, FineWeb-Edu~\cite{penedo2024fineweb}, Dolma~\cite{soldaini2024dolma}



\noindent\textbf{Reasoning and CoT Data:}
Using reasoning data during mid-training has become a trend~\cite{tang2025pangu, yang2025qwen3}. Reasoning data, such as CoT annotations, provides explicit demonstrations of how to decompose complex problems into smaller steps, helping the model learn structured and interpretable reasoning patterns rather than relying only on surface-level cues. CoT reasoning data gives strong performance on math or logic tasks~\cite{sprague2024cot}. Such data is generally obtained from synthetic generation~\cite{abdin2024phi3, abdin2024phi} or through web data filtering. 
% Introducing a proportion of reasoning data in the mid-training phase can help the model gradually strengthen its multi-step inference ability, improve robustness on reasoning-intensive tasks, and enhance the interpretability of its outputs, while still maintaining balanced performance on general instruction-following benchmarks.


% \textbf{Fill-in-Middle (FIM) Data:} Deepseek~\cite{DeepSeekAI2024DeepSeekV3TR, zhu2024deepseek} observes that the FIM strategy does not compromise the next-token prediction capability while enabling the model to accurately predict middle text based on contextual cues. 
\noindent\textbf{Fill-in-Middle (FIM) Data:} Deepseek~\cite{DeepSeekAI2024DeepSeekV3TR, zhu2024deepseek} observes that the FIM strategy does not compromise the next-token prediction capability while enabling the model to accurately predict middle text based on contextual cues. Within the mid-training stage, FIM data is particularly beneficial: it enforces bi-directional contextual reasoning, strengthens long-range dependency modeling, and improves robustness under fragmented inputs. At the same time, FIM increases data efficiency by allowing multiple span-prediction opportunities from a single sample, making it well-suited for late-stage training where high-quality data is scarce. Similar observations have been made in code-pretrained models such as StarCoder~\cite{li2023starcoder} and CodeT5+~\cite{wang2023codet5+}, where infilling objectives significantly enhance both data utilization and downstream generalization.

Table.~\ref{tab:open_sourced_ds} offers a consolidated overview of state-of-the-art models that disclose their mid-training datasets, serving as a resource for the research community.

% 6. Curriculum and Capability-Guided Data:





% \begin{table*}[h]
% \centering
% \scriptsize
% \renewcommand{\arraystretch}{1.2}
% \newcolumntype{L}[1]{>{\raggedright\arraybackslash}p{#1}}
% \caption{Summary of LLM pretraining strategies, token amounts, and annealing data configurations.}
% \begin{adjustbox}{width=\textwidth}
% \begin{tabularx}{\textwidth}{@{} l c L{6cm} L{8cm} @{}}
% \toprule
% \textbf{Model} & \textbf{\# Stages} & \textbf{Stages with Token Amounts} & \textbf{Annealing Data Description} \\
% \midrule
% Nemotron-4~\cite{adler2024nemotron} & 2 & General (8T) $\rightarrow$ Annealing (1T) & Reweighted pre-training tokens emphasizing high-quality sources; added QA-style alignment and upweighted weak domains \\
% \addlinespace
% MiniCPM~\cite{hu2024minicpm} & 2 & General (1T) $\rightarrow$ Annealing (20B) & Diverse and high-quality knowledge, SFT-style data \\
% \addlinespace
% Phi-4~\cite{abdin2024phi} & 2 & General (10T) $\rightarrow$ Annealing (?) $\rightarrow$ Long-context (250B) & Mix of synthetic data and highly filtered reasoning web content; Long-context phase uses 30\% curated long-context data and 70\% recall from pretraining. \\
% \addlinespace
% Deepseek-V3~\cite{DeepSeekAI2024DeepSeekV3TR} & 3 & General (14.8T) $\rightarrow$ Long-context 32k (60B) $\rightarrow$ Long-context 128k (60B) & Context extension only; no specific annealing \\
% \addlinespace
% Zamba-v1~\cite{glorioso2024zamba} & 2 & General (950B) $\rightarrow$ Annealing (1T) & 60\% original pretraining data and 40\% curated math/code/instruction/synthetic datasets \\
% \addlinespace
% MAP-Neo~\cite{Zhang2024MAPNeoHC} & 2 & General (4.5T) $\rightarrow$ Annealing (778B) & 74.5\% English, 17\% code, 8.5\% Chinese; high-quality instructional data \\
% \addlinespace
% AFM~\cite{gunter2024apple} & 3 & General (6.3T) $\rightarrow$ Annealing (1T) $\rightarrow$ Long-context (100B) & Downweight low-quality webcrawl and upweight code, math; Long-context phase uses some annealing data and synthetic long-context Q\&A \\
% \addlinespace
% LLaMA3-405B~\cite{Dubey2024TheL3} & 3 & General (15T) $\rightarrow$ Long-context (0.8T) $\rightarrow$ Annealing (40B) & Upsample very high-quality sources during annealing \\
% \addlinespace
% Granite 3.0~\cite{granite2024granite} & 2 & General (8-10T) $\rightarrow$ Annealing (2T) & Mixed prior and synthetic high-quality data: 45\% Web, 11\% Domain, 10\% Code, 10\% Math, and Instruction, 5\% Multilingual, 5\% Academic, 4\% Technical \\
% \addlinespace
% Hunyuan-Large~\cite{sun2024hunyuan} & 3 & General (7T) $\rightarrow$ Annealing ($\sim$0.35T) $\rightarrow$ Long-context (10B) & Top 5\% highest-quality tokens; Long-context phase uses 25\% long-context data from books/code, 75\% previous pretraining data \\
% \addlinespace
% Yi-Lightning~\cite{wake2024yi} & 3 & Initial (?) $\rightarrow$ Annealing 1 (?) $\rightarrow$ Annealing 2 (Fast-decay) ($\sim$12.5\%) $\rightarrow$ Long-context (20B) & Phase 1 upsamples high-quality, multilingual reasoning data; Phase 2 further intensifies upsampling and introduces early instruction tuning \\
% \addlinespace
% OLMo-2~\cite{olmo20242} & 2 & General (5T) $\rightarrow$ Annealing ($\sim$500B) & Domain-specific high quality data and synthetic data (Dolmino Mix 1124); multiple annealing runs averaged \\
% \addlinespace
% MiniMax-M1~\cite{chen2025minimax} & 2 & Continued training on Minimax-Text-01 (7.5T) $\rightarrow$ Long-context (?) & Diverse natural QA pairs, 70\% STEM, code, book, reasoning data\\
% \addlinespace
% SmolLM2~\cite{allal2025smollm2} & 5 & General 1 (6T) $\rightarrow$ General 2 (2T) $\rightarrow$ General 3 (2T) $\rightarrow$ Annealing (1T) $\rightarrow$ Long-context (75B) & 4\% Cosmopedia v2, 14\% math, 24\% code, 58\% English web; Long-context phase uses 40\% long-context documents, 60\% previous annealing data \\
% \addlinespace
% Pangu Pro MoE~\cite{tang2025pangu} & 3 & General (9.6T) $\rightarrow$ Reasoning (3T) $\rightarrow$ Annealing (0.4T) & Reasoning phase emphasizes complex STEM, coding, and synthetic CoT data; annealing phase includes 20\% instruction-style and 18\% STEM data.\\
% \addlinespace
% Qwen3~\cite{yang2025qwen3} & 3 & General (30T) $\rightarrow$ Annealing (5T) $\rightarrow$ Long-context ($\sim$0.5T) & High-quality tokens increasing STEM, coding, reasoning, synthetic content; Long-context phases uses 75\% 16k--32k, 25\% 4k--16k long-context data \\
% \addlinespace
% Hunyuan-A13B~\cite{yang2025qwen3} & 4 & General (20T) $\rightarrow$ Annealing (300B) $\rightarrow$ Long-context 32k (?) $\rightarrow$ Long-context 256k (?) & No specific annealing dataset details \\
% \bottomrule
% \end{tabularx}
% \label{tab:llm-pretraining-annealing}
% \end{adjustbox}
% \end{table*}


\begin{table*}[h]
\centering
\scriptsize
\caption{Structured summary of stages and data distributions employed during mid-training. For each stage, we indicate the number of phases and the total tokens consumed. In the case of the long-context stage, the number of phase reflects how many steps is used to progressively extend the sequence length.}
\label{fig:data_distribution_in_various_model}
\begin{adjustbox}{width=\textwidth}
\begin{tabularx}{\textwidth}{l| lll |*{8}{>{\centering\arraybackslash}X}}
\toprule
\textbf{Model} & \multicolumn{3}{c|}{\textbf{Stage Configuration (\# Phases, \# Tokens)}} & \multicolumn{8}{c}{\textbf{Data types (used in annealing \& long-context stages)}} \\
\cmidrule(lr){2-4} \cmidrule(lr){5-12}
& \textbf{General} & \textbf{Annealing} & \textbf{Long-context} & \textbf{HQ Web} & \textbf{Code} & \textbf{Math/ STEM} & \textbf{Instruct \& QA} & \textbf{Synthetic} & \textbf{Reasoning \& CoT} & \textbf{FIM} & \textbf{Long-context} \\

% \textbf{Model} & \textbf{General} & \textbf{Annealing} & \textbf{Long-context} & \textbf{HQ Web} & \textbf{Code} & \textbf{Math} & \textbf{Instruct \& QA} & \textbf{Synthetic} & \textbf{Reasoning \& CoT} & \textbf{FIM} & \textbf{Long-context} \\
\midrule
Nemotron-4~\cite{adler2024nemotron} & 1 (8T) & 1 (1T) & - & \ding{51} & \ding{51} & \ding{51}  & \ding{51} &   &   &   &   \\
MiniCPM~\cite{hu2024minicpm} & 1 (1T) & 1 (20B) & - & \ding{51} & \ding{51} & \ding{51} & \ding{51} & \ding{51} &   &   &   \\
Phi-4~\cite{abdin2024phi} & 1 (10T) & 1 & 1 (250B) & \ding{51} & \ding{51} & \ding{51} & \ding{51} & \ding{51} & \ding{51} &  & \ding{51} \\
Deepseek-V3~\cite{DeepSeekAI2024DeepSeekV3TR} & 1 (14.8T) & - & 2 (120B) & \ding{51} & \ding{51} & \ding{51} &   &  &   & \ding{51} & \ding{51} \\
Zamba-v1~\cite{glorioso2024zamba} & 1 (950B) & 1 (1T) & - & \ding{51} & \ding{51} & \ding{51} & \ding{51} & \ding{51} &   &   &   \\
MAP-Neo~\cite{Zhang2024MAPNeoHC} & 1 (4.5T) & 1 (778B) & - & \ding{51} & \ding{51} & \ding{51} & \ding{51} & \ding{51} &   &   &   \\
AFM~\cite{gunter2024apple} & 1 (6.3T) & 1 (1T) & 1 (100B) & \ding{51} & \ding{51} & \ding{51} & \ding{51} & \ding{51} &   &   & \ding{51} \\
LLaMA3-405B~\cite{Dubey2024TheL3} & 1 (15T) & 1 (40B) & 6 (800B) & \ding{51} & \ding{51} & \ding{51} &   &   & \ding{51} &   & \ding{51} \\
Granite 3.0~\cite{granite2024granite} & 1 (8-10T) & 1 (2T) & - & \ding{51} & \ding{51} & \ding{51} & \ding{51} & \ding{51} &   &   &   \\
Hunyuan-Large~\cite{sun2024hunyuan} & 1 (7T) & 1 (350B) & 2 (10B) & \ding{51} & \ding{51} & \ding{51} & \ding{51} & \ding{51} &   &   & \ding{51} \\
Yi-Lightning~\cite{wake2024yi} & 1 & 2 & 3 (20B) & \ding{51} & \ding{51} & \ding{51} & \ding{51} &   & \ding{51} &   & \ding{51} \\
OLMo-2 13B~\cite{olmo20242} & 1 (5T) & 1 ($\sim$600B) & - & \ding{51} & \ding{51} & \ding{51} & \ding{51} & \ding{51} & \ding{51} &   &   \\
MiniMax-M1~\cite{chen2025minimax} & - & 1 (7.5T)\textsuperscript{*} & 3 (358B) & \ding{51} & \ding{51} & \ding{51} & \ding{51} &   & \ding{51} &   & \ding{51} \\
SmolLM2~\cite{allal2025smollm2} & 3 (10T) & 1 (1T) & 1 (75B) & \ding{51} & \ding{51} & \ding{51} &   & \ding{51} & \ding{51} &   & \ding{51} \\
Pangu Pro MoE~\cite{tang2025pangu} & 1 (9.6T) & 2 (3.4T, 32K)$^\dagger$ & - & \ding{51} & \ding{51} & \ding{51} & \ding{51} & \ding{51} & \ding{51} &   &   \\
Qwen3~\cite{yang2025qwen3} & 1 (30T) & 1 (5T) & 1 ($\sim$500B) & \ding{51} & \ding{51} & \ding{51} & \ding{51} & \ding{51} & \ding{51} &   & \ding{51}  \\
Mimo-7B~\cite{xiaomi2025mimo} & 1 (18T) & 1 (4T)  & 2 (2T) & \ding{51} & \ding{51} & \ding{51} &   & \ding{51} & \ding{51} &   & \ding{51} \\
\bottomrule
\end{tabularx}
\end{adjustbox}

\vspace{1mm}
\begin{minipage}{0.95\linewidth}
\scriptsize
\textsuperscript{*} Continued training on Minimax-Text-01. \\
$^\dagger$ Pangu Pro MoE does not have a specific long-context stage. Instead, it is trained with 32K sequence length during the two annealing stages.
\end{minipage}
\end{table*}


% \subsection{Impact of high quality dataset}
% The annealing stage has emerged as a critical opportunity to refine large language models (LLMs) through selective exposure to high-quality data. Unlike earlier phases of pretraining which emphasize coverage and scale, annealing focuses on precision: aligning the model’s learned representation with high-value examples to improve downstream generalization, reasoning, and alignment. In this section, we explore both the theoretical motivations and empirical evidence supporting this design.

% \subsubsection{Experimental Validation}
% Empirical results across multiple open-source LLMs support the efficacy of annealing with high-quality data:

% \textbf{LLaMA3} adopts an explicit annealing phase in which the data mixture is modified to upsample only the most reliable and instruction-aligned content~\cite{Dubey2024TheL3}. This stage was shown to improve reasoning and coherence, especially under long-context inference (see LLaMA3 data curation pipeline).

% \textbf{MiniCPM} finds that introducing high-quality data at the beginning of the decay stage yields significantly greater benefits than reserving it for supervised fine-tuning (SFT)~\cite{hu2024minicpm}. This insight supports initiating capability specialization directly within pretraining.

% \textbf{SmolLM2} performs ablation studies on different English web subsets, which inform its data curation strategy for the annealing stage~\cite{allal2025smollm2}. The model demonstrates measurable performance gains by filtering low-signal web content in favor of educational, math, and coding sources. Training ablation models on 350B tokens each from FineWeb-Edu~\cite{penedo2024fineweb} and DCLM attained performance shown in Table 1. We find that FineWeb-Edu achieves higher scores on the educational benchmarks MMLU, ARC, and OpenBookQA, while DCLM performs better on HellaSwag and CommonsenseQA. These results align with the datasets’ content: FineWeb-Edu prioritizes educational material, while DCLM captures more diverse, conversational styles.

% \textbf{MiniMax-M1} emphasizes the importance of reducing redundancy and repetition during annealing~\cite{chen2025minimax}. The model applies semantic deduplication and uses heuristic-quality scoring to prioritize question-answer and STEM-style data, achieving improved long-context and reasoning robustness.

% \textbf{Deepseek} performs ablation studies on using instruction data during the annealing stage versus incorporating it only during the SFT stage~\cite{bi2024deepseek}. They experimented with adding 5 million instruction samples into the final 10\% of the pretraining phase and observed improved benchmark performance. However, the gains were nearly identical to those achieved by adding the same data during SFT. This suggests that while annealing-stage instruction data can strengthen the base model, it offers limited additional benefit compared to SFT-only usage. DeepSeek ultimately chose not to include instruction data in the pretraining phase.

% \textbf{Phi4} indicates that their experience from previous models showed that question-answer data contributed significantly to various capabilities, such
% as mathematical reasoning and academic performance. Our ablation studies showed that organic questions are substantially more effective than synthetic questions. 

% For long context data, phi4 also conducts the ablation studies on data sources that are inherently long context and artificially created long context data where samples are padded together to fill the sequence. The results show that the former to perform better in longer context tasks.


% \textbf{OpenCoder}~\cite{huang2024opencoder} indicates that it is necessary to ensure that the overall data distribution remains similar to the pretraining phase. A significant distribution shift can lead to catastrophic forgetting in the model’s knowledge. Hence, OpenCoder ensures that 84\% of the annealing data comes from the original distribution of its pretraining corpus.

% \textbf{ProLong-8B}~\cite{gao2024train} indicates that 1) code repositories and books are good sources of long-context data, 2) Training only on long data hurts long-context performance, 3) Choosing a high-quality short-context mix is important.


% ~\cite{springer2025overtrained} proposes a phenomenon called catastrophic overtraining, indicating that extended pre-training can make models harder to fine-tune. Catastrophic overtraining is the phenomenon where extending pre-training beyond a certain token budget results in a decrease in the model’s performance after subsequent modifications.

% OLMo 2~\cite{olmo20242} conducts experiments to show that 1) domain specific data is helpful even in small proportions. This suggests that it is not strictly necessary to have a large proportion of domain-specific data in the annealing mixture, just that domain-specific data is present.

% OLMo 2~\cite{olmo20242} conducts experiments to show that 2) some duplication is beneficial. We see that one copy of the math yields a GSM* score of 61, two copies yields a score of 66, and four copies yields a score of 65. This suggests that even if there is a scarcity of high-quality domain-specific data, duplicating it a small number of times can still provide some gains.

% OLMo 2~\cite{olmo20242} conducts experiments to show that 3) rewriting can help dramatically. We run a microannealing run on a mixture using an inline-annotated form of TinyGSM and compare it to just the ‘Problem Solving’ MIND rewritten variant of TinyGSM. Relative to the baseline, the code version of TinyGSM degrades GSM* performance, while the rewritten version dramatically improves the performance. This suggests the power of rewriting as a tool to cheaply convert data to a more
% amenable form for training.

% \cite{charton2024emergent} On three problems of mathematics: the greatest common divisor, modular multiplication, and matrix eigenvalues, we show that for a fixed number of training steps, models trained on smaller sets of repeated examples outperform models trained on larger sets of single-use examples. This high-
% lights that the benefits of repetition can outweigh those of data diversity. 

% These findings collectively underscore the importance of annealing not merely as a continuation of pretraining, but as a strategically designed specialization phase. It is during this phase that the model's general capabilities are polished and directed, often setting the foundation for successful alignment and instruction tuning in later stages.





% "Why use high quality dataset?" 

% \subsubsection{Gradient-Noise-Scale}
% Because high-quality data has higher gradient noise scale. 

% \subsubsection{Information Bottleneck}
% According to information bottleneck theory, 

% \subsubsection{Curriculum Learning}
% In the theory of curriculum learning, we 

% \subsubsection{experimental validation}
% Llama3 model uses annealing stage to assesses data quality.
% Llama3 data processing process.

% the effect of repetition: minimax.

% SmolLM2 has conducted english web ablations to curate pre-training data.

% MiniCPM: the results of MiniCPM indicate that the benefits of introducing high-quality data at the beginning of the decay stage are much higher than simply adding it during the SFT phase. Therefore, we recommend that specialization and enhancement of model capabilities should start from the decay phase. (contains the data mixture of annealing stage, see section 6.3)



\subsection{Mid-training Data Used In Popular Models}
In this section, we review and compare the mid-training datasets adopted in popular models. To facilitate clarity, Table.\ref{fig:data_distribution_in_various_model} provides a structured summary of the data distributions used during mid-training, along with the number of pre-training stages and total training tokens.

% The Qwen3 models~\cite{yang2025qwen3} are pre-trained through a three-stage process. It is first trained over 30 trillion tokens covering 119 languages and dialects, then further optimized on 5T high-quality tokens at a sequence length of 4K tokens which increases the proportion of STEM, coding, reasoning and synthetic data. For the third stages of long context, Qwen3 is trained on hundreds of billions of tokens, consisting of 75\% 16k–32k and 25\% 4k–16k long-context data.
The Qwen3 models~\cite{yang2025qwen3} are pre-trained in three stages: first on 30T tokens across 119 languages and dialects, then on 5T high-quality 4K-sequence tokens with increased STEM, coding, reasoning, and synthetic data, and finally on hundreds of billions of tokens for long-context training (75\% 16k–32k, 25\% 4k–16k).


% The Phi4 models~\cite{abdin2024phi} are pre-trained on $sim$10T tokens through a two-stage process. Most training tokens (mostly filtered web data) were used in stage 1 of pre-training. Stage 2 was trained on a mixture of synthetic tokens and a much smaller allocation for ultra-filtered and reasoning-heavy web data. Then, the Phi4 models conduct a mid-training of long context training on 75B curated long-context token and 175B recall tokens from the previous stages.
The Phi4 models~\cite{abdin2024phi} are pre-trained on $\sim$10T tokens in two stages: stage 1 with mostly filtered web data, and stage 2 with a mix of synthetic tokens and a smaller portion of ultra-filtered, reasoning-focused web data. For mid training, the models are further trained on 75B curated long-context tokens and 175B recall tokens from earlier stages.

% The Llama3 405B models~\cite{Dubey2024TheL3} are pre-trained on 15T multilingual tokens. The whole pre-training process consists of three stages: (1) initial, (2) long-context (800B tokens), and (3) annealing (final 40M tokens). The pre-training data mix contains roughly 50\% of general knowledge tokens, 25\% of math and reasoning tokens, 17\% code tokens, 8\% multilingual tokens. During the annealing, the data mixture is adjusted to upsample data sources of very high quality. 
The Llama3 405B models~\cite{Dubey2024TheL3} are pre-trained on 15T multilingual tokens across three stages: 1) initial, 2) long-context (800B tokens), and 3) annealing (final 40M tokens). The data mix comprises 50\% general knowledge, 17\% code, 25\% math/reasoning, and 8\% multilingual tokens, with the annealing stage upsampling the highest-quality sources.

The Apple Foundation Model (AFM)~\cite{gunter2024apple} consists of three stages: (1) pre-training stage on 6.3T tokens with a sequence length of 4k, (2) annealing stage on 1T tokens with a sequence length of 8k, which downweights low-quality webcrawl and upweight code and math data, (3) long context stage on 100B tokens with a sequence length of 32k, using the data mixture from the annealing stage, augmented with synthetic long-context Q\&A data.

The Deepseek-V3 models~\cite{DeepSeekAI2024DeepSeekV3TR} are first pre-trained on 14.8T diverse and high-quality tokens. After that, Deepseek-V3 performs two context extension stages, first trained on 60B tokens with a sequence length of 32K and then trained on 60B tokens with a sequence length of 128K.

% The SmolLM2 models~\cite{allal2025smollm2} are trained on 11T tokens, employing three pre-training stages, one annealing stage and one context length extension stage. For the first pre-training stage, the model is trained on 6T high-quality and curated tokens. They adopted 90\% web data and 10\% code data. The web data consists of 60\% FineWeb-Edu~\cite{penedo2024fineweb} and 40\% DCLM ratio for web data. For the second pre-training stage, SmolLM2 is trained on 2T tokens with a mixture of 75\% English web data, 20\% code data, and 5\% math data. For the third pre-training stage, the model is trained on the 2T tokens with a mixture of 74\% English web data, 16\% code data, and 10\% math data. For the fourth annealing stage, the model is trained on 1T high quality tokens with a mixture of 4\% synthetic textbook data (Cosmopedia v2), 14\% math data, 24\% code data, and 58\% english web data. To support long-context extension, SmolLM2 extended the context length from 2k to 8k tokens for 75B tokens. The data mixture contains 40\% long-context documents, while the remaining 60\% followed the stage 4 mixture.
The SmolLM2 models~\cite{allal2025smollm2} are trained on 11T tokens across three pre-training stages, one annealing stage, and one context extension stage. Stage 1 uses 6T curated tokens (90\% web, 10\% code; web split into 60\% FineWeb-Edu~\cite{penedo2024fineweb} and 40\% DCLM). Stage 2 trains on 2T tokens (75\% web, 20\% code, 5\% math), and Stage 3 on 2T tokens (74\% web, 16\% code, 10\% math). The annealing stage adds 1T high-quality tokens (58\% web, 24\% code, 14\% math, 4\% synthetic textbooks). Finally, long-context training extends context length from 2k to 8k tokens over 75B tokens, with 40\% long documents and 60\% from the Stage 4 mixture.

% The Nemotron-4 models~\cite{adler2024nemotron} are trained on 9T tokens for pre-training, with 8T tokens used in general pretraining and the last 1T in annealing. The pre-train data mixture consists of 70\% English, 15\% multilingual, and 15\% code data. 
The Nemotron-4 models~\cite{adler2024nemotron} are pre-trained on 9T tokens, with 8T for general pretraining and 1T for mid-training. The data mix includes 70\% English, 15\% multilingual, and 15\% code.
In the mid-training phase, two distinct data distributions are employed~\cite{parmar2024nemotron}. The first distribution, which makes up the majority of this phase, consists of tokens already seen during pretraining but reweighted to emphasize higher-quality sources. 
% The second distribution introduces a small number of QA-style alignment examples to help the model improve its ability to answer such queries in downstream tasks. Additionally, data sources from areas where the model previously showed lower accuracy are up-weighted.
The second distribution introduces a smaller portion QA-style alignment examples and up-weights domains where the model underperforms.

% The OLMo 2 models~\cite{olmo20242} adopt a 2-stage process for pre-training. For the first stage of general pre-training, the model is trained on $\sim$5T tokens (primarily web-based data). For the second stage of annealing, the model is trained on the Dolmino Mix 1124 dataset, which is smaller but contains higher-quality data and synthetic data targeted at improving weak model capabilities. To better leverage this high-quality data and improve generalization, they repeat annealing stage multiple times with different random data orders and average the resulting models. For the 7B model, it performs three runs of 50B tokens each and average the results. For the 13B model, it runs three times for 100B tokens and a fourth time for 300B tokens, then averages all four models.
The OLMo 2 models~\cite{olmo20242} adopt a two-stage pre-training process. The first stage uses $\sim$5T tokens of primarily web data. The second mid-training stage employs the Dolmino Mix 1124 dataset, which is smaller but higher-quality and includes synthetic data for strengthening weak capabilities. To maximize generalization, OLMo 2 repeats mid-training with different random orders and averages the resulting models. The 7B model trains three runs of 50B tokens each, while the 13B model performs three runs of 100B tokens plus one of 300B tokens, averaging across all.

% The Yi-Lightning model~\cite{wake2024yi} adopts a three-stage approach (initial pre-training, annealing, and fast-decay training stage). The initial pre-training stage emphasizes data diversity to establish robust foundational capabilities across diverse domains. The annealing phase shifts data distribution and gradually upsamples high-quality data, with a focus on complex reasoning and low-resource multilingual support. The final fast-decay stage, accounting for 12.5\% of total tokens, boosts high-quality data emphasis and incorporates early instruction-tuning adaptation. After the three pre-training stage, the model is further trained on 20B training tokens to develop robust long-context capabilities.
The Yi-Lightning model~\cite{wake2024yi} follows a three-stage training paradigm: an initial pre-training stage emphasizing data diversity for broad foundational capabilities; an annealing stage that gradually upsamples high-quality data with emphasis on complex reasoning and low-resource multilingual support; and a fast-decay stage (12.5\% of total tokens) that further strengthens high-quality data usage and introduces early instruction-tuning adaptation. After these three stages, the model undergoes an additional 20B-token training phase to enhance long-context performance.

% The Zamba-v1~\cite{glorioso2024zamba} is pre-trained in two phases. It is first trained on 950B tokens from open web datasets. During the annealing stage, Zamba-v1 is trained over a dataset composed of 60\% original pretraining data and 40\% high-quality annealing data. The annealing set included over 100 curated datasets spanning math (e.g., StackMathQA), code (e.g., EvolInstructCode), instruction tuning (e.g., OpenOrca), and synthetic examples from stronger language models. Most datasets were trained for one epoch, with selected high-quality subsets receiving two. 
The Zamba-v1 model~\cite{glorioso2024zamba} is pre-trained in two phases: an initial stage on 950B open web tokens, followed by a mid-training stage with a mixture of 60\% pretraining data and 40\% high-quality datasets. The mid-training set spans over 100 curated sources, including math (StackMathQA), code (EvolInstructCode), instruction tuning (OpenOrca), and synthetic data from stronger models. Most datasets were trained for one epoch, while select high-quality subsets were upsampled for two.

% For MiniMax-M1 model~\cite{chen2025minimax}, they continue training on Minimax-Text-01 model with an additional 7.5T tokens using optimized data quality and mixture to improve reasoning and long-context capabilities while maintaining diversity. They refine web and PDF parsing and apply enhanced cleaning and semantic deduplication to prioritize natural QA pairs from diverse sources while avoiding synthetic data. STEM, code, book, and reasoning-related content make up 70\% of the dataset.  Training proceeds with a constant learning rate of 8e-5 for 2.5T tokens, followed by decay to 8e-6 over 5T tokens. For long-context extension, we adopt a four-stage approach, gradually increasing the context window from 32K to 1M tokens to avoid instability caused by mismatched optimization between early and late layers in lightning attention architectures.
For the MiniMax-M1 model~\cite{chen2025minimax}, training continues from Minimax-Text-01 with an additional 7.5T tokens using optimized mixtures to enhance reasoning and long-context capabilities while preserving diversity. Data quality is improved through refined web/PDF parsing, enhanced cleaning, and semantic de-duplication, prioritizing natural QA pairs over synthetic data. STEM, code, books, and reasoning-related content constitute 70\% of the corpus. Training uses a constant learning rate of 8e-5 for 2.5T tokens, then decays to 8e-6 over 5T. Long-context extension follows a four-stage schedule, expanding the window from 32K to 1M tokens to ensure stability in lightning attention.

% The pre-training of Pangu Pro MoE~\cite{tang2025pangu} follows a three-phase structure inspired by cognitive development: a general phase (9.6T tokens) to build foundational knowledge and language skills, a reasoning phase (3T) to enhance complex reasoning through high-quality STEM, code, and synthetic CoT data, and an annealing phase (0.4T) to refine behavior and transition into instruction tuning. During the reasoning phase, they designed synthetic short and long CoT difficult samples and extended sequence lengths (32K) to improve alignment with long CoT tasks, while leveraging more reasoning-focused data than prior models. The annealing phase increases instruction-style data to 20\% and advanced STEM content to 18\%, applying curriculum-based sampling and ablation studies with a 7B proxy model to optimize data strategies. 
The pre-training of Pangu Pro MoE~\cite{tang2025pangu} follows a three-phase process inspired by cognitive development: a general phase (9.6T tokens) for foundational knowledge, a reasoning phase (3T) to strengthen complex reasoning with high-quality STEM, code, and synthetic CoT data, and an annealing phase (0.4T) to refine behavior and transition into instruction tuning. In the reasoning phase, synthetic short- and long-form CoT samples and extended contexts (32K) are introduced to align with long reasoning tasks. The annealing stage increases instruction-style data (20\%) and advanced STEM content (18\%), using curriculum-based sampling and ablation with a 7B proxy to optimize data strategies.

The Hunyuan-Large model~\cite{sun2024hunyuan} is pre-trained on 7T tokens, which contains nearly 1.5T tokens of high-quality and diverse synthetic data. During the learning rate annealing phase, the model is trained on 5\% of the highest-quality pre-training tokens, which plays a pivotal role in augmenting the model's performance. After the learning rate annealing phase, Hunyuan-Large is trained on longer sequences to enable its longer-context capability (up to 256K tokens). The training corpus during long-text phase is consists of 25\% natural long-context data obtained from books and codes and 75\% normal length pre-training data, inspired by~\cite{gao2024train}.

The Hunyuan-A13B model~\cite{yang2025qwen3} is pre-trained on 20T tokens. Then in midt-traininig, it implemented a fast annealing stage on 300B tokens. Following the annealing phase, Hunyuan-A13B progressed through two long-context stages to expand its context length to 32k tokens, then to 256K tokens.

The MiniCPM model~\cite{hu2024minicpm} is trained by a two-stage pre-training strategy. During the general pre-training phase, MiniCPM uses 1T coarse-quality pre-training data, which is abundant and can support continuous training when provided with more computational resources. During the mid-training phase, the model is trained on 20B tokens using a mixture of the pre-training data and high-quality knowledge \& ability-oriented SFT data.

% The MAP-Neo model~\cite{Zhang2024MAPNeoHC} is pre-trained in two phases. In the general pre-training stage, the model is trained on 4.5T generic text tokens (60\% English, 15\% code, 25\% Chinese) to develop its general text generation capability. In the annealing phase, the model is trained on 778B high-quality tokens (74.5\% English, 17\% code, 8.5\% Chinese), which primarily consists of high-quality instruction data. The proportion of code increases from 15\% to 17\%. The deliberate enrichment of the dataset with a higher ratio of code, coupled with instructional inputs, ensures a more robust and versatile model, adept at tackling complex coding tasks as well as understanding and generating professional responses in different fields.
The MAP-Neo model~\cite{Zhang2024MAPNeoHC} is pre-trained in two phases. In the general stage, it uses 4.5T tokens (60\% English, 25\% Chinese, 15\% code) to build general text generation ability. The mid-training phase adds 778B high-quality tokens (74.5\% English, 8.5\% Chinese, 17\% code), with a greater emphasis on instructional and coding data. This enrichment improves robustness and equips the model for complex coding tasks as well as professional domain-specific responses.

% The Granite 3.0 models~\cite{granite2024granite} are trained with a two-stage pre-training stage. In the first general pre-training stage, the dense and MoE models are trained on 10T and 8T tokens respectively. Stage 1 training data consists of unstructured multilingual language data from 5\% Web,
% 11\% Domain, 10\% Code, 10\% Math, 10\% Instruction, 5\% Multilingual, 5\% Academic and 4\% Technical. In stage 2, we train on a mixture of 2T tokens of data. Some of the data sources for stage 2 are the same as the stage 1 data sources, mixed with a small amount of high-quality open-source and synthetic corpora with permissive licenses. 
The Granite 3.0 models~\cite{granite2024granite} follow a two-stage pre-training setup. In Stage 1, dense and MoE variants are trained on 10T and 8T tokens, respectively, using a mixture of 5\% Web, 11\% Domain, 10\% Code, 10\% Math, 10\% Instruction, 5\% Multilingual, 5\% Academic and 4\% Technical. Stage 2 adds 2T tokens drawn partly from Stage 1 sources, supplemented with high-quality open-source and synthetic corpora under permissive licenses.

MiMo-7B~\cite{xiaomi2025mimo} uses a 3-stage pre-training strategy trained on 25T tokens: Stage 1 (General) focuses on high-quality natural data by removing low-value content and upsampling professional domains. Stage 2 (Annealing) increases mathematics and code data to about 70\% to strengthen specialized skills. Stage 3 (Long-context) incorporates around 10\% synthetic data for tasks like math, code, and creative writing, while extending the context length from 8K to 32K tokens, then to 64K tokens, to improve complex task performance.


\subsection{Insights For Mid-Training Data}
% \textbf{Upsampling high-quality, instruction-aligned, and educational content consistently improves reasoning and coherence.}
% Several models, including LLaMA3 and MiniCPM, emphasize the importance of emphasizing high-quality, instruction-following data during annealing. LLaMA3 explicitly adjusts its data mixture to upsample only the most reliable, instruction-aligned content, significantly improving reasoning and coherence, especially in long-context inference. Annealing with a high emphasis on mathematics and coding leaves general-domain perplexity unchanged, while simultaneously enhancing performance on targeted tasks; hence, the modified data mixture is retained in LLaMA3’s final checkpoints. Similarly, MiniCPM demonstrates that introducing high-quality data early in the decay phase yields greater performance gains than reserving it for supervised fine-tuning, suggesting that capability specialization should begin during pretraining.
% \textbf{Selecting high-quality content improves model performance.}
% In contrast to the early stages of pretraining, which prioritize scale and coverage, the annealing phase emphasizes data quality. By aligning the model’s learned representations with carefully curated high-quality data, annealing enhances downstream generalization, reasoning, and alignment. Ablation studies in \cite{Su2024NemotronCCTC, wang2025ultra, penedo2024fineweb} consistently demonstrate that models trained on filtered, high-quality corpora outperform those trained on unfiltered data. Moreover, although there are many available high-quality public dataset, each dataset has distinct strengths. For example, FineWeb-Edu~\cite{penedo2024fineweb} excels on academic benchmarks such as MMLU and ARC, whereas DCLM~\cite{NEURIPS2024_19e4ea30} performs better on commonsense and reasoning tasks. Hence, strategically combining complementary datasets is crucial for maximizing model performance across diverse domains.
\textbf{Selecting high-quality content improves model performance.}
Unlike early pre-training, which prioritizes scale and coverage, the mid-training phase emphasizes quality. Curated datasets align model representations with reasoning, generalization, and alignment objectives. Ablation studies confirm that filtered corpora consistently outperform unfiltered ones~\cite{Su2024NemotronCCTC, wang2025ultra, penedo2024fineweb}. Different datasets, however, have complementary strengths—for instance, FineWeb-Edu excels on academic benchmarks such as MMLU and ARC, while DCLM performs better on commonsense and reasoning tasks~\cite{penedo2024fineweb, NEURIPS2024_19e4ea30}. Thus, combining diverse high-quality sources is critical for broad performance gains.

\noindent\textbf{Prioritizing educational, coding, and math data boosts performance on reasoning and STEM benchmarks.} SmolLM2~\cite{allal2025smollm2} has conducted extensive ablation studies to refine their mid-training mixtures. SmolLM2 demonstrates consistent gains by replacing low-signal web text with structured educational, math, and coding data, indicating that high-quality, domain-focused corpora provide stronger supervision signals during late-stage training. MiniMax-M1 complements this approach by introducing semantic deduplication and heuristic scoring to prioritize question-answer pairs and STEM-oriented data, thereby enhancing reasoning accuracy and improving long-context generalization. Together, these findings highlight the value of systematically curating high-utility data sources in mid-training to maximize both efficiency and downstream task performance.

% \noindent\textbf{The role of instruction-style data during annealing remains uncertain, with models reporting both benefits and limitations.}
% Upsampling high-quality, instruction-aligned, and educational content is shown to consistently improve reasoning and coherence, as evidenced by models like LLaMA3~\cite{Dubey2024TheL3} and MiniCPM~\cite{hu2024minicpm}. LLaMA3 selectively amplifies reliable, instruction-following data during annealing, enhancing performance on reasoning-intensive tasks, particularly under long-context conditions, without harming general-domain perplexity. MiniCPM further demonstrates that incorporating high-quality data early in the annealing phase yields greater returns than postponing it to supervised fine-tuning, indicating that capability specialization can begin during pretraining. However, other studies raise doubts about the marginal utility of instruction data at this stage. Deepseek~\cite{bi2024deepseek} finds that adding 5 million instruction samples near the end of pretraining yields similar benefits to applying them during SFT, suggesting diminishing returns for instruction-style data beyond fine-tuning. This divergence highlights an ongoing debate: while instruction-aligned data can strengthen core reasoning abilities, its optimal placement within the training pipeline remains unclear.

\noindent\textbf{The role of instruction-style data during mid-training remains uncertain, with models reporting both benefits and limitations.}
Upsampling high-quality instruction and educational content often improves reasoning and coherence~\cite{Dubey2024TheL3, hu2024minicpm}. LLaMA3~\cite{Dubey2024TheL3} amplifies reliable instruction-following data to boost long-context reasoning without harming general perplexity, while MiniCPM~\cite{hu2024minicpm} shows that introducing such data earlier in mid-training is more effective than deferring it to fine-tuning. However, Deepseek~\cite{bi2024deepseek} reports that adding 5M instruction samples late in pretraining provides gains comparable to SFT, suggesting diminishing returns. This divergence underscores an open question: instruction data can aid reasoning, but its optimal placement in the training pipeline remains unresolved.


\noindent\textbf{Long-context capabilities benefit from careful data source selection, not just sequence length.}
ProLong-8B~\cite{gao2024train} and Phi-4~\cite{abdin2024phi} investigate strategies for long-context optimization. Both find that training solely on long data can hurt performance, and a balanced mixture with high-quality short-context data is critical. Phi-4 shows that naturally long-form content outperforms artificially concatenated data in long-context reasoning. Books and code repositories are highlighted as effective long-context sources.


\noindent\textbf{Domain-specific data is impactful even in small proportions.}
OLMo 2~\cite{olmo20242} finds that even a small fraction of domain-specific data (e.g., math and coding) within the mid-training mix can produce disproportionate performance gains. This underscores the quality-over-quantity principle at this stage, where carefully selected high-utility domains provide strong supervision signals, enhance reasoning ability, and improve transfer to specialized downstream tasks without requiring large-scale rebalancing of the corpus.


% \textbf{Synthetic data}
% The broad usage of synthetic data improves the quality and diversity
% of training data, which enables the model to learn richer representations effectively and generalize better to unseen data~\cite{sun2024hunyuan}.

\noindent\textbf{Synthetic data substantially enhances training efficiency and generalization.}
The broad usage of synthetic data improves both the quality and diversity of the training corpus, enabling models to acquire richer representations and generalize more effectively to unseen data~\cite{sun2024hunyuan}. Beyond simple data augmentation, synthetic corpora can be tailored to emphasize underrepresented domains, complex reasoning, or structured formats, thereby compensating for limitations in naturally collected text. This makes synthetic data particularly valuable in the mid-training stage, where curated, high-signal samples amplify learning efficiency and reinforce specialized capabilities without requiring large-scale raw corpus expansion.

% \textbf{Reasoning data}
% FineMath4+ achieves a 2x improvement on GSM8K and a 6x improvement on MATH compared to InfiMM-WebMath, demonstrating the importance of retaining high-quality mathematical content with reasoning~\cite{allal2025smollm2}.
\noindent\textbf{Reasoning-oriented data is critical for strengthening mathematical and logical capabilities.}
FineMath4+ achieves a 2x improvement on GSM8K and a 6x improvement on MATH compared to InfiMM-WebMath, underscoring the importance of preserving high-quality math corpora with step-by-step reasoning~\cite{allal2025smollm2}. Such data provides explicit demonstrations of problem decomposition and logical inference, allowing models to internalize structured reasoning patterns rather than relying solely on surface heuristics. Incorporating reasoning data during mid-training thus amplifies gains in mathematical problem-solving and enhances generalization to broader reasoning benchmarks.

\noindent\textbf{Moderate repetition and rewriting of high-quality data can substitute for scale.}
OLMo 2~\cite{olmo20242} and related studies explore the role of repetition and rewriting. Repeating high-value data (e.g., math examples) a few times yields performance gains, with diminishing returns beyond a certain point. Additionally, rewriting tasks into more amenable formats (e.g., inline annotations or simplified prompts) dramatically improves performance over structurally similar but unrewritten variants. Charton et al.\ further validate that models trained on smaller repeated datasets can outperform those trained on larger, unrepeated datasets, especially in mathematical reasoning.

\noindent\textbf{Maintaining distributional continuity prevents catastrophic forgetting.}
OpenCoder~\cite{huang2024opencoder} highlights the risk of distributional shift during mid-training, where a mismatch between mid-training and pretraining corpora can lead to catastrophic forgetting and degraded generalization. To address this, OpenCoder retains 84\% of its mid-training data from the original pretraining distribution, ensuring continuity while gradually introducing higher-quality subsets. This strategy stabilizes knowledge retention and preserves broad capabilities, showing the importance of balancing innovation in data selection with consistency in distributional coverage.

\noindent\textbf{Overtraining beyond a critical token budget may degrade downstream tunability.}
A recent study~\cite{springer2025overtrained} on catastrophic overtraining warns that excessive pretraining reduces a model’s adaptability during subsequent fine-tuning, as representations may become overspecialized and less flexible. This highlights the importance of defining careful token budgets and stopping criteria during mid-training to balance knowledge consolidation with downstream tunability.

Collectively, these insights position mid-training not as a mere continuation of pretraining, but as a targeted specialization phase.
It is during mid-training that models are strategically refined, leveraging focused, high-quality, and instruction-rich data to enhance core reasoning, coherence, and long-context understanding. These refinements often lay the groundwork for successful instruction tuning and alignment in downstream stages.

\section{Learning Rate Scheduler}
\label{sec:lrsched}
% In this section, we review the literature on learning rate schedulers. We begin with an overview of theoretical motivations and hypotheses surrounding learning rate schedulers in the pretraining of LLMs, highlighting their role in shaping optimization dynamics and generalization. We then introduce common types learning rate schedulers, and discuss by surveying the specific scheduling strategies adopted in recent state-of-the-art LLMs.
In this section, we review the literature on learning rate (LR) schedulers. We begin with key factors affected LR scheduling design, then we introduce common types of LR schedulers and the LR schedulers in recent LLMs. We also summarize some valuable key insights for selecting appropriate LR scheduler.

\subsection{Key factors affected LR Scheduling}

{LR schedulers are critical components in optimizing LLMs, as they directly affect optimization stability, convergence speed and generalization performance. 
LR scheduling are largely affected by several factors:
\begin{itemize}[leftmargin=*]
    \item \textit{Warmup setting:} The warmup phase mitigates optimization instability during the early iterations, while the decay phase governs convergence dynamics, typically using cosine decay. Recent analyses \cite{kalra2024warmup} further emphasize the effectiveness of properly setting warmup parameters in stabilizing training and preventing divergence in training Transformer models. Gilmer et al. \cite{gilmer2022loss} empirically show that warm-up reduces loss sharpness and improves optimization conditioning, enabling larger LR.
    \item \textit{Decay strategy:} The decay stage of LR schedulers have also been widely studied. You et al. \cite{you1908does} investigate the role of LR decay in modern neural networks, showing that a large initial LR suppresses the memorization of noisy data, while a gradual decay facilitates the learning of complex patterns.
    \item \textit{Batch size:} Increasing the batch size lowers the variance of stochastic gradient estimates, which in turn supports the use of higher learning rates. Empirical results by Goyal et al.~\cite{goyal2017accurate} demonstrate that, in large mini-batch regimes, employing a cautious initial rate with a properly tuned warm-up phase alleviates optimization difficulties.
\end{itemize}
% Interestingly, according to Kaplan et al.\cite{kaplan2020scaling}, the final performance of language models is largely insensitive to the choice of learning rate scheduler, provided that the total learning rate budget is sufficiently large. In contrast, Hoffmann et al.\cite{hoffmann2022training} provided empirical evidence that learning rate schedulers significantly affect optimization outcomes, demonstrating that the choice of decay strategy can have a substantial impact on final model performance.
While these factors provide practical guidance for designing effective schedulers, the broader question of how much LR scheduling truly matters remains debated. Kaplan et al.~\cite{kaplan2020scaling} argued that the final performance of language models is largely insensitive to the specific scheduler, so long as the total LR budget is sufficiently large. In contrast, Hoffmann et al.~\cite{hoffmann2022training} presented empirical evidence that the choice of decay strategy substantially influences optimization dynamics and final model quality, suggesting that scheduler design can have non-trivial consequences in large-scale training.}

\subsection{Common Learning Rate Scheduler Types}

{The design of LR schedules in LLM training has evolved substantially over time. The original Transformer~\cite{vaswani2017attention} introduced a short warm-up phase followed by an inverse square root decay, whose advantage was that it did not require prior knowledge of the total number of training steps. As models and datasets scaled up, however, cosine decay schedules~\cite{loshchilov2016sgdr} became widely adopted in LLM pretraining due to their empirical effectiveness, as seen in GPT-3~\cite{brown2020language}, Gopher~\cite{rae2021scaling}, and Chinchilla~\cite{hoffmann2022training}.}

{Subsequent work questioned whether cosine annealing represents the optimal choice. Defazio et al.~\cite{defazio2023optimal} demonstrated that linear decay consistently outperforms cosine and other schedulers across optimization methods, while Bergsma et al.~\cite{bergsma2025straight} reported similar findings when training LLMs with AdamW. More recently, Ibrahim et al.~\cite{ibrahim2024simple} proposed an “infinite” LR schedule, which eliminates the need for fixed token budgets and repeated warmups by employing a four-phase strategy, thus supporting continual training with reduced forgetting. Despite these advances, many modern LLMs continue to adopt the 10x cosine variant introduced by~\cite{kalra2024warmup}, given its slight empirical advantage over standard cosine decay.}

{Across these developments, most training pipelines converge on a composite structure: a linear warm-up phase, which gradually increases the LR to mitigate instability, followed by a decay phase whose form—cosine, linear, or multi-stage—varies across models. In multi-phase or annealing-style training paradigms, stage-wise schedulers such as WSD apply distinct decay behaviors at different stages to better align with the evolving data distribution and optimization objectives. By modulating the LR in tandem with training phases, these approaches facilitate a smooth transition from rapid exploration in early stages to fine-grained convergence later, thereby enhancing both optimization efficiency and generalization performance.}


{In the following, we provide a detailed overview of mainstream decay-phase LR schedulers. Figure~\ref{fig:lrs} provides a unified visualization of commonly used LR schedules (linear, cosine, exponential, cyclical, and WSD), to highlight their characteristic dynamics over training steps. To facilitate a clear comparison, all LR schedulers are set up with the same warmup steps, peak LR, total training steps, and initial and final learning rates of zero.
}

\begin{figure}
    \centering
    \includegraphics[width=\linewidth]{2_lrs.png}
    \caption{Comparison of different LR schedulers.}
    \label{fig:lrs}
\end{figure}

\subsubsection{Linear Scheduler}

The linear LR scheduler was originally proposed in \cite{goyal2017accurate} and has since been widely adopted in large-scale model training.

In the decay phase, the LR decreases linearly from a peak value \(\eta_{peak}\) to a final value \(\eta_{end}\) over \(S_{\text{dc}}\) steps:
\begin{equation}
\eta_t = \eta_{peak} - (\eta_{peak} - \eta_{end}) \cdot \frac{t}{S_{\text{dc}}}, \quad 0 \leq t \leq S_{\text{dc}}.
\end{equation}

This linear decay schedule enables a smooth transition from rapid exploration in early training to more stable convergence in later stages. Its simplicity and effectiveness make it a common choice in both academic and industrial settings.



\subsubsection{Cosine Scheduler}
The cosine LR scheduler typically consists of a linear warmup phase followed by a cosine decay, gradually decreases the LR following a cosine curve after it reaches its maximum after warmup phase.
The decay function is given by
\[
\eta_t = \eta_{end} + \frac{1}{2}(\eta_{peak} - \eta_{end}) \left(1 + \cos\left(\frac{\pi t}{S_{\text{dc}}}\right)\right), \quad 0 \leq t \leq S_{\text{dc}},
\]
where \(\eta_{peak}\) and \(\eta_{end}\) denote the peak and final learning rates, respectively.

A key hyperparameter is the decay length \(S_{\text{dc}}\) at which cosine decay decreases to the minimum for the first time, often set equal to the total number of training steps \(S_{\text{total}}\) when training length is predetermined. Prior studies~\cite{hu2024minicpm, kaplan2020scaling, hoffmann2022training} have shown that both \(S_{\text{dc}} < S_{\text{total}}\) and \(S_{\text{dc}} > S_{\text{total}}\) lead to suboptimal performance. In particular, setting \(S_{\text{dc}} = S_{\text{total}}\) improves training efficiency, as it avoids prematurely decaying the LR or keeping it high for too long. Possible explanation for the effectiveness of cosine LR schedulers with \(S_{\text{dc}} = S_{\text{total}}\) lies in the balance between an extended high LR phase due to cosine decay, which may aid in global exploration, and a full decay phase, which could promote stable convergence dynamics. Consequently, the intermediate checkpoints tend to be suboptimal, complicating the continual pretraining of an existing language model.


\subsubsection{Exponential Scheduler}

The exponential LR scheduler~\cite{li2019exponential} gradually decreases the LR by multiplying it with a fixed decay factor at each step, enabling smooth and continuous reduction throughout training. It is defined as:
\[
\eta_t =  \eta_{init} \cdot e^{-kt}
\]
where \(\eta_{init}\) is the initial LR, and \(k\) is a decay constant that controls the rate of exponential decay.

Its simplicity and continuous decay pattern make it suitable for tasks with stable long-term optimization needs. However, it offers less control and interpretability over scheduling stages and may decay too quickly without careful tuning.

\subsubsection{Knee Scheduler}

The Knee scheduler~\cite{iyer2023wide} is an explore–exploit LR strategy, inspired by the wide-minima density hypothesis, which hypothesizes that narrow minima are significantly higher than wide minima in the loss landscape of deep neural networks.  It comprises two distinct phases: an initial exploration phase, during which the model is trained with a constant high LR for a sufficient duration to increase the likelihood of converging near a wide minimum; and a subsequent exploitation phase, wherein the LR linearly decays to zero without requiring any additional hyperparameters.


\subsubsection{Cyclical Scheduler}

The cyclical scheduler~\cite{smith2017cyclical} cyclically vary between predefined boundary values in a cyclic manner rather than decaying monotonically. The cyclical LR at training step \( t \) is defined as
\begin{align}
\eta_t &= \eta_{init} + (\eta_{peak} - \eta_{init}) \cdot \nonumber \\
&\quad \max\left(0,\ 1 - \left| \frac{t}{s} - 2 \left\lfloor 1 + \frac{t}{2s} \right\rfloor + 1 \right| \right)
\end{align}
where \( \eta_{init} \) and \( \eta_{peak} \) are the minimum and maximum LR, \( s \) is the half-cycle length (i.e., the number of steps from \( \eta_{init} \) to \( \eta_{peak} \)), and \( \lfloor \cdot \rfloor \) denotes the floor function, which returns the greatest integer less than or equal to its input. Its adaptive nature eliminates the need for extensive experimentation to identify optimal learning rates and schedules, while often achieving near-optimal accuracy in fewer iterations. 


\subsubsection{WSD Scheduler}
MiniCPM~\cite{hu2024minicpm} proposes Warmup-Stable-Decay (WSD) LR scheduler that divide the training stage into three phases: the warmup stage, the stable training stage, and the decay stage. The function form of WSD is:

\[
WSD(T; s) = 
\begin{cases}
\frac{s}{W} \eta, & s < W \\
\eta, & W \leq s \leq T \\
f(s - T)\eta, & T < s < S
\end{cases}
\]

where $0<f(s-T)\leq1$ is a decreasing function about $s$, $\eta$ is the maximum LR. The WSD scheduler adopts a three-phase structure designed together with two-phase pretraining. It enables efficient convergence by achieving rapid loss reduction in the final decay stage using only 10\% of total tokens, and is particularly effective when high-quality data is introduced mid-training. Compared to cosine schedulers, which require a predetermined total number of steps to achieve optimal decay, WSD achieves comparable performance without this constraint. Its stage-wise design supports resuming from intermediate checkpoints for efficient decay, making it highly flexible for continual pretraining.


\subsubsection{Power Scheduler}

The Power scheduler~\cite{shen2024power} extends from the WSD scheduler and, building on Maximum Update Parameterization ($\mu P$), enables zero-shot LR transfer across different hyperparameters by uncovering a power-law relationship that governs the optimal LR. It is defined as
\begin{equation}
\eta_{\text{power}}(n) = \min\left(\eta_{peak},\; \beta \cdot a n^b\right),
\end{equation}
where $\beta$ is the batch size, $n$ is the number of trained tokens, $a$ and $b$ are the power-law coefficients (amplitude and decay exponent, respectively), and $\eta_{peak}$ is the upper bound on the LR. The final PowerLR scheduler, combined with a warmup and decay stage, is defined as follows:

\begin{equation}
\small
\eta_n = 
\begin{cases}
\frac{n}{N_{\text{}}} \cdot \eta_{\text{power}}(N_{\text{}}) & n < N_{\text{}}, \\
\eta_{\text{power}}(n) & N_{\text{}} \leq n \leq N - N_{\text{decay}}, \\
f(n) \cdot \eta_{\text{power}}(N - N_{\text{decay}}) & n > N - N_{\text{decay}}.
\end{cases}
\label{eq:power}
\end{equation}

The generality of Power scheduler is substantiated by extensive empirical results reported in~\cite{shen2024power}, which confirm its effectiveness across varying hyperparameter scales and model configurations.


\subsubsection{Multi-step Scheduler}
DeepSeek~\cite{bi2024deepseek} replaces the commonly used cosine scheduler with a multi-step LR schedule featuring a linear warmup phase, followed by two discrete decay steps triggered at specific token ratios. Two discrete LR drops are applied at the stage boundaries, enabling efficient continual training by allowing earlier phases to be reused across different training scales.

% \paragraph{Section Summary.} Learning rate schedulers play a pivotal role in shaping training dynamics. Traditional schedulers like linear and exponential offer simplicity, while modern variants such as cosine, WSD, and Power schedulers provide more flexible, phase-aware control. Each scheduler type offers trade-offs between stability, generalization, and adaptability for continual training or large-scale pretraining.

% Widely adopted learning rate schedulers, such as linear, cosine, or exponential schedulers, are primarily designed to smoothly adjust the learning rate over training. More recent LLMs adopt multi-phase strategies that partition training into distinct stages, with each phase governed by its own scheduling policy. Notably, both the WSD scheduler and the multi-step scheduler in DeepSeek follow this principle, differing primarily in how they define and utilize these stages.


\subsection{Learning Rate Schedulers in Recent LLMs}

\begin{table*}[ht]
\centering
\caption{Overview of LR Schedulers in representative LLMs.}
\resizebox{\textwidth}{!}{
\begin{tabular}{lllllc}
\toprule
\textbf{Model} & \textbf{LR Decay Scheduler} & \textbf{LR Path} & \textbf{Optimizer} & \textbf{Staged Training} \\
\midrule
GPT-3 (175B)         & Cosine      & (0) $\xrightarrow{\text{}}$ {6e-4} $\rightarrow$ {6e-5} & Adam & \ding{55} \\
Gopher (280B)    & Cosine      & {1e-7} $\xrightarrow{\text{}}$ {4e-5} $\rightarrow$ {4e-6} & Adam    & \ding{55}  \\
Chinchilla (70B)   & Cosine      & {1e-7} $\xrightarrow{\text{}}$ {1.25e-4} $\rightarrow$ {1.25e-5} & AdamW    & \ding{55}  \\
PaLM (540B) & Inverse square root & {1e-2} (constant) $\rightarrow$ {1.98e-5} & Adafactor &  \ding{55} \\
OPT (175B) & Linear (restart\textsuperscript{\dag}) & 0 $\xrightarrow{\text{}}$ {1.2e-4} $\rightarrow$ {1.2e-5} & AdamW & \ding{55} \\
LLaMA (65.2B)  & Cosine       & (0) $\xrightarrow{\text{}}$ {1.5e-4} $\rightarrow$ {1.5e-5} & AdamW    & \ding{55} \\
LLaMA 2 (70B)  & Cosine       & (0) $\xrightarrow{\text{}}$ {1.5e-4} $\rightarrow$ {1.5e-5} & AdamW    & \ding{55}   \\
Jais (13B)  & Linear       & 0 $\xrightarrow{\text{}}$ {1.2e-2} $\rightarrow$ {1.2e-3} & AdamW  &  \ding{55} \\
Qwen (14B)  & Cosine       & 0 $\xrightarrow{\text{}}$ {3e-4} $\rightarrow$ {3e-5} & AdamW    & \ding{55} \\
Falcon (40B)   & Cosine       & (0) $\xrightarrow{\text{}}$ {1.85e-4} $\rightarrow$ {1.85e-5} & AdamW    & \ding{55}   \\
DeepSeek (67B)  & Multi-stage       & (0) $\xrightarrow{\text{}}$ {3.2e-4} $\rightarrow$ {1e-4} $\rightarrow$ {3.2e-5} & AdamW    & \ding{51} \\
OLMo 1.7 (7B)   & Cosine  $\rightarrow$ Linear       & (0) $\xrightarrow{\text{}}$ {3e-4} $\rightarrow$ {3e-5} & AdamW & \ding{51} \\
InternLM2 (20B)  & Cosine & (0) $\xrightarrow{\text{}}$ {3e-4} $\rightarrow$ {3e-5}  & AdamW    & \ding{51} \\
MiniCPM (2.4B)  & WSD & (0) $\xrightarrow{\text{}}$ 0.01 $\rightarrow$ $0.01  \cdot f(s - T)$ & Adam    & \ding{51} \\
DeepSeek V2 (67B)  & Multi-stage & 0 $\xrightarrow{\text{}}$ {2.4e-4} $\rightarrow$ {7.6e-5} $\rightarrow$ {2.4e-5}  & AdamW & \ding{51} \\
Zamba (7B) & Cosine $\rightarrow$ Exponential &  {1.5e-4} $\rightarrow$ {7.5e-5} (cosine); {1.1e-4} $\rightarrow$ {1e-7} (exponential) & Zero-1 distributed &  \ding{51} \\
MAP-Neo (7B) & Two-stage & {2e-5} $\xrightarrow{\text{warmup}}$ {2e-4} $\rightarrow$ {2e-5} (cosine); {2e-4} $\rightarrow$ {1e-4} (exponential) & AdamW & \ding{51}  \\
AFM (3B) & Two-stage & (0) $\xrightarrow{\text{}}$ 0.01 $\rightarrow$ {5e-5} (cosine); (0) $\rightarrow$ {3e-4} $\rightarrow$ {3e-7} & AdamW & \ding{51} \\
LLaMA 3 (405B)   & Cosine $\rightarrow$ Linear & (0) $\xrightarrow{\text{}}$ {8e-5} $\rightarrow$ {8e-7} (cosine) $\rightarrow$ 0 (linear) & AdamW    & \ding{51} \\
Hunyuan-Large (52B out of 389B MoE)   & Three phases & (0) $\xrightarrow{\text{}}$ $\eta_{max}$ $\rightarrow$ $\frac{1}{10}\eta_{max}$ (gradual decay) & AdamW    & \ding{51}   \\
Open-Coder (8B)  & WSD & (0) $\xrightarrow{\text{}}$ {3e-4} $\rightarrow$ {1e-5} & - & \ding{51}  \\
Yi-Lightning   & Three-stage & (0) $\xrightarrow{\text{}}$ $\eta_{\max} \rightarrow 0.5\,\eta_{\max} \rightarrow \eta_{\min}$ & - & \ding{51} \\
Granite 3.0 (8B)  & Power & (0) $\xrightarrow{\text{}}$ 0.02 $\rightarrow$ $4\beta \cdot f(n, N, N_{\text{decay}}) \cdot (N - N_{\text{decay}})^{-0.51}$ & AdamW & \ding{51} \\
Phi-4 (14B)  & Linear & 0 $\xrightarrow{\text{}}$ {3e-4} $\rightarrow$ (0) &  - &  \ding{51}  \\
YuLan-Mini (2.42B)  & WSD & 0 $\xrightarrow{\text{}}$ {1e-2} $\rightarrow$ {5.22e-5} (1-sqrt) &  AdamW & \ding{51}   \\
DeepSeek V3 (37B out of 671B MoE)  & Multi-stage & 0 $\xrightarrow{\text{}}$ {2.2e-4} $\xrightarrow{\text{constant}}$ {2.2e-5} (cosine) $\xrightarrow{\text{constant}}$ {7.3e-6} (constant) &  AdamW & \ding{51}   \\
OLMo 2 (13B)  & Cosine $\rightarrow$ Linear & 0 $\xrightarrow{\text{}}$ {9e-4} $\rightarrow$ {9e-5} (cosine) $\rightarrow$ 0 (linear) & AdamW & \ding{51}  \\
MiniMax-Text-01 (45.9B out of 456B MoE)  & Multi-stage & (0) $\xrightarrow{\text{}}$ {2e-4} $\xrightarrow{\text{constant}}$ {1.3e-4} $\xrightarrow{\text{constant}}$ {3e-5} (exponential) $\rightarrow$ 0 (linear) & AdamW & \ding{51}  \\
SmolLM2 (1.7B)  & WSD & 0 $\xrightarrow{\text{}}$ {5e-4} $\rightarrow$ 0 & AdamW & \ding{51} & \\
xGen-small (9B)  & Revised power~\cite{shen2024power} & - & (Optax lib) & \ding{51}  \\
MiMo (7B) & Multi-stage & 0 $\xrightarrow{\text{}}$ {1.07e-4} $\xrightarrow{\text{constant}}$ {3e-5} (cosine) $\xrightarrow{\text{constant}}$ {1e-5} (cosine) & AdamW & \ding{51}  \\
Pangu Pro MoE (3B)  & Cosine & {3e-4} $\rightarrow$ {1e-7} & AdamW & \ding{51}  \\
MiniCPM4 (8B)  & WSD & N.A. & - & \ding{51}  \\
Hunyuan A13B  & Multi-stage & 0 $\xrightarrow{\text{}}$ {3e-4} $\rightarrow$ {3e-5} (cosine) $\rightarrow$ {8e-6} (rapid cosine) & - & \ding{51}  \\
SmolLM3   & WSD & (0) $\xrightarrow{\text{}}$ {2e-4} $\rightarrow$ 0 & AdamW & \ding{51}  \\
\bottomrule
\end{tabular}
}

\vspace{1mm}  % 表格和注释之间留空
\noindent
\begin{minipage}{\linewidth}
\textbf{Note:}
The “Staged Training” column indicates distinct pre-training phases prior to supervised fine-tuning, such as data distribution shifts or long-context adaptation. Restart\textsuperscript{\dag} means that LR restarts resume from the checkpointed value. Unless otherwise specified, both the initial and final learning rates are assumed to be zero and denoted as $(0)$. The decay function is defined as \( f(s - T) = 0.5^{\frac{s - S}{5000}} \), where \( S \) denotes the total number of training steps and \( s \) is the current step.
\end{minipage}

\label{tab:lr-schedulers}

\end{table*}


The emergence of LLMs marks a major milestone in the development of NLP. This shift was initiated by the introduction of the Transformer architecture~\cite{vaswani2017attention}, which enabled the development of early large-scale pre-trained models such as BERT~\cite{devlin2019bert} and GPT~\cite{brown2020language}. These early works laid the foundation for scaling both model size and training tokens. The release of GPT-3~\cite{brown2020language} further demonstrated the power of large-scale pre-training, revealing emergent capabilities that arise from the substantial  increase in model parameters and training tokens. Since then, numerous LLMs, including Gopher~\cite{rae2021scaling}, Megatron-Turing~\cite{adler2024nemotron}, Chinchilla~\cite{hoffmann2022training}, PaLM~\cite{chowdhery2023palm}, OPT~\cite{zhang2022opt}, and BLOOM~\cite{workshop2022bloom}, have been released, advancing the frontiers of model scaling, training efficiency, and downstream task performance. In this section, we provide a comprehensive review of the LR schedulers adopted in pre-training phase of the representative LLMs. Furthermore, the mid-training stage plays a critical role in smoothing the transition of learning dynamics and maintaining training stability between pre-training and supervised fine-tuning. We analyze the LR scheduling strategies of LLMs that explicitly incorporate an annealing phase.

With the rapid increase in model size, careful tuning of hyperparameters, including the LR scheduler and batch size, is essential for maintaining training stability and achieving convergence under a limited computational budget. Several works~\cite{mccandlish2018empirical}\cite{kreps2022all}\cite{brown2020language} found that training larger models benefits from the use of larger batch sizes combined with smaller learning rates, contributing to more stable and efficient optimization. The LR scheduler in \textit{GPT-3}~\cite{brown2020language} employs a linear warm-up and followed by a cosine decay, reducing to 10\% of the initial rate, followed by continued training at this reduced level. Consequently, the paradigm shift to focusing heavily on model scaling and pre-training strategies.
\textit{Gopher}~\cite{rae2021scaling} scales transformer-based models up to 280 billion parameters, achieving strong performance across a range of benchmarks. It utilizes the Adam optimizer with a warm-up that increases the LR from $10^{-7}$ to a peak value, followed by cosine decay by a factor of 10. 
Building on this, \textit{Chinchilla}~\cite{hoffmann2022training} investigates the optimal scaling between model size and training tokens under a fixed FLOPs budget, concluding that model size and tokens should scale proportionally. The hypothesis that aligning the cosine decay schedule length with the total number of training tokens improves performance is empirically validated, and is thus adopted in \textit{Chinchilla}. Besides, they empirically show that AdamW outperforms Adam in large-scale training, and therefore adopt AdamW as the optimizer for \textit{Chinchilla}.
PaLM~\cite{chowdhery2023palm} adopts Adafactor~\cite{shazeer2018adafactor} with a fixed LR in the early phase, followed by an inverse square root decay with respect to the training step, and incorporates dynamic weight decay along with empirically tuned hyperparameters to ensure stability in large-scale training.
\textit{LLaMA}~\cite{touvron2023llama, touvron2023llama2, dubey2024llama} form a family of open-weight transformer-based language models developed by Meta, with parameter scales ranging from 7B to 405B across versions. LLaMA 1 and 2 use the AdamW optimizer and a LR schedule consisting of a linear warm-up followed by cosine decay to 10\% of the peak LR. LLaMA 3 extends this recipe into a three-stage process, using a longer linear warm-up and cosine decay to 1\% of the peak LR during the initial pre-training phase, followed by a final annealing phase that linearly decays the LR to zero over the final 40 million tokens. % llama4 moe [TODO] add?
Similarly, Qwen~\cite{bai2023qwen} uses a cosine LR schedule with a specified peak LR followed by a decay to 10\% of the peak LR. In Qwen 2.5~\cite{qwen2025qwen25technicalreport}, the authors further explore scaling laws to optimize hyperparameters across both dense and Mixture-of-Experts (MoE) models, enabling efficient training and performance parity. The pretraining process is restructured into three stages, general stage, reasoning stage and long context stage, with accelerated LR decay during the reasoning stage to facilitate more effective optimization in Qwen 3~\cite{yang2025qwen3}.
\textit{Nemotron-4}~\cite{parmar2024nemotron, adler2024nemotron} introduces an additional continued training phase, employing two distinct data distributions and a steeper LR decay slope over absolute magnitude to help the model transition smoothly from pre-training corpus and better learn newly emphasized corpus.
\textit{DeepSeek}~\cite{bi2024deepseek} and \textit{DeepSeek-V2}~\cite{liu2024deepseek} employ a warmup-and-step-decay LR schedule, consisting of a 2K-step linear warm-up followed by discrete drops at fixed training milestones, a design that facilitates continual training reuse while achieving performance comparable to that of cosine decay. \textit{DeepSeek-V3}~\cite{liu2024deepseekv3} extends this scheduling strategy by introducing a plateau phase followed by cosine decay after the warm-up, and subsequently transitions into a long-context extension stage.
Meanwhile, \cite{hu2024minicpm} put forward the Warmup-Stable-Decay (WSD) LR scheduler, which closely resembles that of \textit{DeepSeek} and is likewise tailored for continual training, enabling effective reuse of intermediate model checkpoints. Their experiments show that, during the decay stage, the loss rapidly drops as the LR decreases, reaching or even falling below that of the Cosine schedule at step $T = S$. 
The Power scheduler~\cite{shen2024power}, developed through further research based on WSD and $\mu$P, introduces a new LR schedule that combines a linear warmup phase, a slow power-law decay, and a fast exponential decay. Granite 3.0~\cite{granite2024granite} adopts the Power scheduler to train its lightweight foundation models. % TODO

Beyond the mainstream LLMs discussed above, a variety of LR schedulers have been explored in LLM. Among them, the \textit{cosine LR scheduler} is the most widely adopted~\cite{kaplan2020scaling, hoffmann2022training, rae2021scaling, gunter2024apple, cai2024internlm2, touvron2023llama, touvron2023llama2, dubey2024llama, brown2020language, olmo20242, tang2025pangu, bai2023qwen, almazrouei2023falcon}, where the LR increases linearly during the warm-up phase and then decays following a cosine curve.
Another commonly used approach is the linear LR scheduler, where the LR increases linearly during warm-up and then decreases linearly to zero or a minimal value over the which linearly increases the LR during warm-up and then decays it linearly to zero or a small constant, is another widely adopted strategy~\cite{zhang2022opt, abdin2024phi, bergsma2025straight, defazio2023optimal}.
\cite{glorioso2024zamba, ibrahim2024simple, hu2024minicpm} find it better to rewarm the LR followed by a rapid exponential (instead of linear) decay.
However, some studies suggest that linear decay may not always be optimal. \cite{glorioso2024zamba, ibrahim2024simple, hu2024minicpm} argue that exponential decay leads to better performance than linear decay.
In addition, the inverse square root LR scheduler, typically combined with the Adafactor optimizer, has been adopted by several models, including \textit{T5}~\cite{raffel2020exploring}, \textit{PaLM}~\cite{chowdhery2023palm}, and \textit{OpenMoE}~\cite{xue2024openmoe}.

In addition to those continuous decay strategies, some models have adopted variants of the multi-step LR scheduler. \textit{DeepSeek} employs discrete LR drops at fixed milestones, while \textit{MiniCPM} integrates the WSD scheduler; both strategies facilitates checkpoint reuse and support continual training. Empirical results from both models show that the multi-step LR scheduler performs comparably to cosine decay during pretraining. With the introduction of the annealing stage, the WSD strategy has been adopted in several recent works~\cite{hu2024yulan, allal2025smollm2, huang2024opencoder}, owing to its flexibility in accommodating multi-stage pretraining.
\textit{MAP-Neo}, extending from \textit{MiniCPM}, employs a two-stage scheduler combining warm-up and cosine decay, followed by a subsequent exponential decay phase. Similarly, an increasing number of models~\cite{sun2024hunyuan,li2025minimax,wake2024yi,nijkamp2025xgen}, incorporate multiple pretraining stages, typically involving distribution shifts and/or sequence length extensions, and accordingly adopt multi-phase LR strategies with distinct designs. A recent trend emphasizes the use of an explicit annealing phase, which has shown effectiveness in improving optimization. Models that explicitly adopt such an annealing stage~\cite{} typically employ multi-stage LR schedules, such as those used in \cite{hu2024minicpm} and \cite{bi2024deepseek}.
Table~\ref{tab:lr-schedulers} provides a comparative summary of LR schedulers and optimizer settings, as explicitly reported for representative LLMs in official sources.
% gemini 1 2 没写lrs optimizer或者pertaining stage的训练细节



% \begin{table*}[t]
% \centering
% \footnotesize
% \begin{tabularx}{\textwidth}{@{} l X c c @{}}
% \toprule
% \textbf{Model} & \textbf{Learning Rate Scheduling} & \textbf{Annealing} & \textbf{Scale} \\
% \midrule
% Qwen & Cosine learning rate schedule with a specified peak learning rate for each model size. The learning rate is decayed to 10\% of the peak. AdamW optimizer. & no & LLM \\
% % Qwen 2.5 & --- & no & LLM \\
% Qwen 3 & Three-stage process. Accelerated LR decay during Reasoning Stage (S2). & no & LLM \\
% MiniCPM & Warmup-Stable-Decay (WSD). After more training ($S'$), can anneal to match Cosine LRS at Cosine($S'$). & yes & SLM \\
% Gemma & Similar to ZeRO-3. & no & SLM \\
% % Gemma 3 & --- & --- & --- \\
% LLaMA 3 405B & AdamW with peak LR $8 \times 10^{-5}$, linear warm-up (8k steps), cosine decay to $8 \times 10^{-7}$ over 1.2M steps. & yes & LLM \\
% LLaMA 4 & MetaP to set per-layer learning rates and initialization scales. & mid-training & LLM \\
% AFM & Linear warm-up (5k steps), cosine decay to 0.005 of peak. & no & SLM \\
% xGen-small & Revised power schedule with extended flat phase, followed by slow decay and fast annealing decay on high-quality data. & yes & SLM \\
% DeepSeek V3 & Linear ramp to $2.2\times10^{-4}$ (2k steps), constant till 10T, cosine decay to $2.2\times10^{-5}$ over 4.3T, final 500B split into $2.2\times10^{-5}$ (333B) and $7.3\times10^{-6}$ (167B). & no & LLM \\
% SmolLM2 & AdamW with $(\beta,\beta_2) = (0.9,0.95)$, WSD: 2k-step warmup, stable at $5.0\times10^{-4}$, optional decay over 10\% of total steps. & yes & SLM \\
% MiniMax-Text-01 & Linear warm-up (500 steps) to $2\times10^{-4}$, constant training for 7.2T, then exponential decay to $3\times10^{-5}$ over 1T tokens. & no & LLM \\
% % GRANITE 3.0 & --- & --- & --- \\
% Phi-4 & Linear warm-up and decay. Peak LR 0.0003, weight decay 0.1, batch size 5760. & no (midtraining stage) & LLM \\
% Nemotron-4 15B & LR decay schedule at training end improves quality. & no & LLM \\
% Nemotron-4 340B & LR decay slope prioritized over magnitude to aid transition from pretraining to final stage. & no & LLM \\
% Zamba & Cosine decay from $1.5\times10^{-4}$ to $7.5\times10^{-5}$, then exponential decay. Annealing follows infinite LR schedules. & yes & LLM \\
% \bottomrule
% \end{tabularx}
% \caption{Learning rate scheduling and annealing strategies across language models. Scale indicates model size class, e.g., LLM = Large Language Model, SLM = Small Language Model.}
% \label{tab:lrs-table}
% \end{table*}


\subsection{Key Insights}

\noindent\textbf{Linear warm-up with cosine/linear decay is widely used but largely heuristic.}
This scheme is the most common choice in LLM training due to its empirical stability, even though its design is guided more by practice than theory. While several alternatives, such as the infinite LR schedule~\cite{ibrahim2024simple}, have been proposed to relax assumptions like fixed token budgets and repeated warm-ups, they often lack robustness across architectures and tasks. Consequently, cosine and linear decay remain the prevailing strategies, reflecting both their demonstrated empirical reliability and the broader reliance on heuristic tuning in the field.

\noindent\textbf{Warm-up helps stability.}
Warm-up phases are known to alleviate gradient instability and loss sharpness in early training~\cite{goyal2017accurate, gilmer2022loss, kalra2024warmup}, thereby facilitating the use of larger peak learning rates. However, the optimal duration and scaling of warm-up still lack theoretical justification and vary significantly across model scales and optimizers.

\noindent\textbf{Decay strategies remain unsettled and system-dependent.}
The design of the decay phase in LR schedules is far from resolved. Although cosine annealing has been widely adopted in LLM training~\cite{brown2020language, rae2021scaling, hoffmann2022training}, recent empirical analyses suggest that linear decay can yield superior convergence and generalization in certain settings~\cite{defazio2023optimal, bergsma2025straight}. Beyond these comparisons, the effectiveness of a decay strategy interacts intricately with factors such as batch size, model scale, and optimizer dynamics, leading to system-dependent behaviors that are not yet theoretically characterized.

\noindent\textbf{Scaling laws help LR selection.}
Empirical scaling laws provide practical guidance for choosing learning rates across model sizes and training regimes, though they remain incomplete and sensitive to model-specific factors. Kaplan et al.~\cite{kaplan2020scaling} uses scaling law to show that the optimal LR depends on the target loss: smaller values are needed near convergence for stability, whereas larger rates can be effective in short, compute-limited runs. Their analysis also indicated that larger models require smaller learning rates, while smaller models tolerate more aggressive ones. More recently, Li et al.~\cite{li2025predictable} proposed a universal scaling law for hyperparameter selection in LLM pretraining, finding that a fixed final LR leads to more stable convergence.

\noindent\textbf{Current evidence highlights both the critical role and the unresolved uncertainties of LR scheduling.}
Conflicting empirical results, for example, the claim that final performance is insensitive to schedule shape under sufficient LR budget~\cite{kaplan2020scaling}, versus evidence of significant variance across decay types~\cite{hoffmann2022training, defazio2023optimal}, highlight the lack of consensus and reproducible benchmarks in this space.
\section{Long Context Extension}
\label{sec:longcontext}
In this section, we focus on the long context extension phases in popular LLMs. 

% \subsection{Background and Formulation}
% \subsubsection{Self‑attention with Positional Encoding}
% \label{sec:prelim-sa}

% Let $\mathbb S_N=\{w_i\}_{i=1}^{N}$ be an $N$‑token sequence and
% $\mathbb E_N=\{\mathbf x_i\}_{i=1}^{N}$ its word‑embedding vectors,
% $\mathbf x_i\in\mathbb R^{d}$.  
% Positional information is injected via task‑specific functions
% $f_q,f_k,f_v$ that transform the \emph{content} embedding $\mathbf x_i$
% and the \emph{position} index $i$ into query, key and value vectors:
% \begin{equation}
% \label{eq:qkv}
% \mathbf q_{m}=f_q(\mathbf x_{m},m),\quad
% \mathbf k_{n}=f_k(\mathbf x_{n},n),\quad
% \mathbf v_{n}=f_v(\mathbf x_{n},n).
% \end{equation}

% Given \eqref{eq:qkv}, single‑head scaled dot‑product self‑attention
% computes the weight of position $n$ when attending from position $m$:
% \begin{equation}
% \label{eq:attn-weights}
% a_{m,n}\;=\;
% \frac{\exp\bigl(\mathbf q_{m}^{\top}\mathbf k_{n}/\sqrt d\bigr)}
%      {\displaystyle\sum_{j=1}^{N}\exp\bigl(\mathbf q_{m}^{\top}\mathbf k_{j}/\sqrt d\bigr)},
% \qquad
% \mathbf o_{m}\;=\;\sum_{n=1}^{N}a_{m,n}\,\mathbf v_{n},
% \end{equation}
% where $\mathbf o_{m}$ is the output representation for token $w_{m}$.

% \subsubsection{Rotary Position Embedding (RoPE)}
% \label{sec:rope}

% Let $W_q,W_k\in\mathbb R^{d\times d}$ be the linear maps that turn a
% content embedding $\mathbf x_m$ into \emph{base} query/key vectors
% $\tilde{\mathbf q}_m=W_q\mathbf x_m$ and
% $\tilde{\mathbf k}_m=W_k\mathbf x_m$.
% RoPE applies a position‑dependent \emph{rotation} in each 2‑D
% sub‑plane \emph{after} these linear projections:
% \begin{align}\label{eq:rope-rotate-qk}
% \mathbf q_m
%     = f_q(\mathbf x_{m},m)
%     = R^d_{\Theta,m}\,\tilde{\mathbf q}_m
%     = R^d_{\Theta,m}\,W_q\mathbf x_m, 
% \\
% \mathbf k_n
%     = f_k(\mathbf x_{n},n)
%     = R^d_{\Theta,n}\,\tilde{\mathbf k}_n
%     = R^d_{\Theta,n}\,W_k\mathbf x_n .
% \end{align}

% The rotation matrix $R^d_{\Theta,m}$ is defined as follows. 
% Write the $d$‑dimensional rotation as the block‑diagonal matrix
% \begin{equation}
% R^d_{\Theta,m}
% =\bigoplus_{j=1}^{d/2}
% \begin{pmatrix}
% \cos m\theta_j & -\sin m\theta_j\\
% \sin m\theta_j &  \cos m\theta_j
% \end{pmatrix},
% \label{eq:R-block}
% \end{equation}
% where $\bigoplus$ concatenates the $2{\times}2$ rotation blocks along the embedding dimension, 
% $\Theta=\bigl\{ \theta_{j}, j\in[1,,2,\dots,d/2] \bigr\}$, 
% and each block rotates a pair
% $\bigl(\tilde q_{m,2j-1},\tilde q_{m,2j}\bigr)$ (and likewise for
% $\tilde{\mathbf k}_n$) by the phase
% \[
% \theta_j = b^{-\frac{2(j-1)}{d}}, \qquad b=10000, 
% \qquad j=1,\dots,d/2 .
% \]

% The relative‑position property is as follows. 
% Because $R^d_{\Theta,m}$ is orthogonal and
% $R^d_{\Theta,m}^{\top} R^d_{\Theta,n}=R^d_{\Theta,{m-n}}$,
% the dot product inside scaled self‑attention becomes
% \begin{equation}\label{eq:rope-relative}
% \mathbf q_m^{\top}\mathbf k_n
%   = \tilde{\mathbf q}_m^{\top} R^d_{\Theta_{m-n}} \tilde{\mathbf k}_n ,
%   = \mathbf x_{m}^{\top} W_q^{\top} R^d_{\Theta_{m-n}} W_k {\mathbf x}_n
% \end{equation}
% which depends only on the \emph{relative} position $m-n$.
% This invariance enables extrapolation to context lengths well beyond
% those seen during pre‑training.

\subsection{Background and Formulation}
\subsubsection{Self‑attention with Positional Encoding}
\label{sec:prelim-sa}

We briefly fix notation and recall scaled dot–product self‑attention.
Let $\mathbb S_N=\{w_i\}_{i=1}^{N}$ be an $N$‑token sequence and
$\mathbb E_N=\{\mathbf x_i\}_{i=1}^{N}$ the corresponding token embeddings with
$\mathbf x_i\in\mathbb R^{d}$.
To inject order information, we use task‑specific maps
$f_q,f_k,f_v$ that take the \emph{content} embedding $\mathbf x_i$ together with its
\emph{position} index $i$ and produce the query, key and value vectors:
\begin{equation}
\label{eq:qkv}
\mathbf q_{m}=f_q(\mathbf x_{m},m),\quad
\mathbf k_{n}=f_k(\mathbf x_{n},n),\quad
\mathbf v_{n}=f_v(\mathbf x_{n},n).
\end{equation}
Given \eqref{eq:qkv}, single‑head scaled dot–product attention from position $m$ to all
positions $n\in\{1,\dots,N\}$ computes
\begin{equation}
\label{eq:attn-weights}
a_{m,n}\;=\;
\frac{\exp\!\bigl(\mathbf q_{m}^{\top}\mathbf k_{n}/\sqrt d\bigr)}
     {\displaystyle\sum_{j=1}^{N}\exp\!\bigl(\mathbf q_{m}^{\top}\mathbf k_{j}/\sqrt d\bigr)},
\qquad
\mathbf o_{m}\;=\;\sum_{n=1}^{N}a_{m,n}\,\mathbf v_{n},
\end{equation}
where $\mathbf o_m$ denotes the output at token $w_m$.
For clarity, we omit the (standard) causal mask and multi‑head split; all subsequent
derivations apply per head identically.

\subsubsection{Rotary Position Embedding (RoPE)}
\label{sec:rope}

RoPE encodes positions by \emph{rotating} the query/key vectors in a set of
two‑dimensional sub‑planes of the embedding space.
Let $W_q,W_k\in\mathbb R^{d\times d}$ denote the linear maps that first produce the
\emph{base} query/key vectors
$\tilde{\mathbf q}_m=W_q\mathbf x_m$ and $\tilde{\mathbf k}_n=W_k\mathbf x_n$.
A position‑dependent rotation is applied after the projections:
\begin{align}\label{eq:rope-rotate-qk}
\mathbf q_m
    = f_q(\mathbf x_{m},m)
    = R^d_{\Theta,m}\,\tilde{\mathbf q}_m
    = R^d_{\Theta,m}\,W_q\mathbf x_m, 
\\
\mathbf k_n
    = f_k(\mathbf x_{n},n)
    = R^d_{\Theta,n}\,\tilde{\mathbf k}_n
    = R^d_{\Theta,n}\,W_k\mathbf x_n .
\end{align}
Here $R^d_{\Theta,m}$ is an orthogonal block‑diagonal rotation that acts independently
on $d/2$ disjoint 2‑D sub‑planes:
\begin{equation}
\label{eq:R-block}
R^d_{\Theta,m}
=\bigoplus_{j=1}^{d/2}
\begin{pmatrix}
\cos(m\theta_j) & -\sin(m\theta_j)\\
\sin(m\theta_j) &  \cos(m\theta_j)
\end{pmatrix},
\end{equation}
where $\bigoplus$ concatenates the $2{\times}2$ blocks along the embedding dimension,
$\Theta=\{\theta_{j}\mid j=1,2,\dots,d/2\}$ is the set of base angular frequencies, and
each block rotates the coordinate pair
$\bigl(\tilde q_{m,2j-1},\tilde q_{m,2j}\bigr)$ (and analogously for
$\tilde{\mathbf k}_n$).
The original RoPE schedule sets
\[
\theta_j = b^{-\frac{2(j-1)}{d}}, \qquad b=10000, \qquad j=1,\dots,d/2 .
\]
Intuitively, the $j$‑th 2‑D sub‑plane behaves like a complex plane whose phase advances
linearly with the position index $m$ at rate $\theta_j$.

A crucial property of RoPE is that inner products of rotated vectors depend only on
\emph{relative} positions. Because $R^d_{\Theta,m}$ is orthogonal and
${R^d_{\Theta,m}}^{\top} R^d_{\Theta,n}=R^d_{\Theta,{m-n}}$, 
the attention dot‑product becomes
\begin{equation}\label{eq:rope-relative}
\mathbf q_m^{\top}\mathbf k_n
  = \tilde{\mathbf q}_m^{\top} R^d_{\Theta,{m-n}} \tilde{\mathbf k}_n
  = \mathbf x_{m}^{\top} W_q^{\top} R^d_{\Theta,{m-n}} W_k\,\mathbf x_n ,
\end{equation}
which depends only on the displacement $m-n$.
This relative‑position invariance is at the heart of many long‑context strategies:
with suitable frequency remapping, models trained at a certain window length can
generalize to much longer contexts during fine‑tuning or inference.

% ---------------------------------------------------------------
%  Generic RoPE extrapolation: only the frequency map h(θd) varies
% ---------------------------------------------------------------
\subsection{A Unified Overview of Context Extension Methods}
% \subsubsection{Generic frequency–remapping template}
\label{sec:rope-h-only}

% Unified long-context frequency-scaling table (drop-in snippet)
% Requires: \usepackage{booktabs}
\begin{table*}[t]
\centering
\small
\caption{Unified frequency-scaling view of long-context families (fix absolute index $g(m)=m$). Let $s=L_{\text{new}}/L_{\text{train}}$. All methods differ only in how they remap RoPE base frequencies $\theta \mapsto h(\theta)$; attention uses $\cos(m\,h(\theta)),\sin(m\,h(\theta))$.}
\label{tab:unified-h-theta}
\resizebox{\textwidth}{!}{%
\begin{tabular}{p{2cm} p{4.6cm} p{2.7cm} p{4.8cm} p{4.8cm} p{4.8cm}}
\toprule
\textbf{Family} & \textbf{Unified map $h(\theta)$} & \textbf{Extra stabilizer} & \textbf{Fixes vs previous} & \textbf{Pros} & \textbf{Cons} \\
\midrule
PI (Positional Interpolation)
& $h(\theta)=\theta/s$ \, (uniform) 
& None
& Avoids naive extrapolation by compressing positions into the seen range
& Conceptually simple; light tuning; robust up to moderate $s$
& Uniform compression hurts high-frequency (local) signals; short-context quality may drop at large $s$ \\
\addlinespace
NTK-aware (ABF-style)
& $h(\theta)=\theta/s^{\rho(\theta)}$, with $\rho(\theta)\!\in[0,1]$ increasing toward low frequencies
& None
& Reduces long-range attention decay vs.\ PI while keeping more locality
& Better locality/short-length retention than PI; works with minimal tuning
& Hyperparameter coupling (e.g., target length vs.\ curvature); mid-band still mixes interp/extrap \\
\addlinespace
NTK-by-part (piecewise)
& $\displaystyle
h(\theta)=\begin{cases}
\theta & \text{(high freq)}\\[2pt]
\theta/s & \text{(low freq)}\\[2pt]
\theta/s^{\rho(\theta)} & \text{(mid freq)}
\end{cases}$
& None
& Preserves the shortest wavelengths (local detail) while still extending range
& Strong short-context preservation; cleaner trade-off across bands
& More boundaries/knobs to set; modest tuning often needed \\
\addlinespace
YaRN
& Same piecewise $h(\theta)$ as NTK-by-part
& Length-dependent logit temperature / scaling $t(s)$
& Stabilizes loss/perplexity for fast extension; reduces data/steps vs.\ prior methods
& Efficient route to 32k–128k; robust convergence with few steps
& Extra temperature and schedule to choose; still handcrafted frequency policy \\
\addlinespace
LongRoPE
& Non-uniform $h(\theta)$ discovered by search; applied progressively (stages)
& Progressive extension + short-length readjustment
& Surpasses handcrafted rules by optimizing $h$; avoids catastrophic angles at ultra-long lengths
& Reaches $\sim$\,2M context with limited extra steps; recovers short-length performance
& Pipeline complexity (search + stages + readjustment); fewer closed-form guarantees \\
\bottomrule
\end{tabular}%
}
\vspace{-4pt}
\end{table*}

\noindent\textbf{Goal.} We present a \emph{single} formulation that covers five representative long-context families---\textbf{PI}, \textbf{NTK}, \textbf{NTK-by-part}, \textbf{YaRN}, and \textbf{LongRoPE}---while keeping the \emph{absolute position} fixed. This isolates what changes in all methods: the \emph{frequency spectrum} used by RoPE.

\noindent\textbf{Unified formulation for context extension.}
We introduce a single scalar function
\[
  h:\theta_j \;\mapsto\; h(\theta_j)
\]
that rescales each base frequency $\theta_j$.  
Any of these extended‑context position embedding $f'$ can be expressed with the original rope embedding $f$ as 
\begin{equation}
  f^{\prime}_W(\mathbf x_m,m,\theta_j)
  \;=\;
  f_W\!\bigl(\mathbf x_m,\,m,\,h(\theta_j)\bigr)
  \;=\;
  \mathbf R_{h(\theta_j),m}\,W\,\mathbf x_m, 
  \label{eq:rope-extrapol-h}
\end{equation}
where $W\in \{W_q, W_k\}$. These methods are summarize in Table~\ref{tab:unified-h-theta}.

% The \emph{scale factor} $s$ is defined as 
% $
%     s = L_{\text{target}} / L_{\text{train}}, 
%     \label{eq:scale-h}
% $
% where $L_{\text{train}}$ is the pre‑training window and
% $L_{\text{target}}>L_{\text{train}}$ is the desired length. 

% The \emph{wavelength} $\lambda_j$ is defined as the number of tokens
% required to complete a full $2\pi$ rotation:
% $
%   \lambda_j
%   \;=\;
%   2\pi/\theta_j
%   \;=\;
%   2\pi\,b^{2(j-1)/d},
%   \label{eq:wavelength}
% $
% for every two‑dimensional rotary sub‑plane indexed by $j$
% ($j=1,\dots,\tfrac{d}{2}$, and $d$ is the embedding dimension),
% and $\theta_j$ is its base angular frequency
% (cf.\ Eq.~\eqref{eq:rope-extrapol-h}; typically
% $\theta_j=b^{-2(j-1)/d}$ with $b=10\,000$).  
% Smaller $\theta_j$ (i.e.\ larger $\lambda_j$) correspond to
% lower‑frequency dimensions, while larger $\theta_j$ yield
% high‑frequency dimensions that rotate more rapidly across positions.

% \begin{figure*}[t]
%   \centering
%   \subfloat[$h(\theta_j)$]{\includegraphics[width=.32\linewidth]{CompassLLM/figures/4_long_context/Effective_Frequency.png}\label{fig:freq-map}}
%   \hfill
%   \subfloat[$\lambda_j$]{\includegraphics[width=.32\linewidth]{CompassLLM/figures/4_long_context/Effective_wavelength.png}\label{fig:wavelengths}}
%   \hfill
%   \subfloat[$h(\theta_j)/\theta_j$]{\includegraphics[width=.32\linewidth]{CompassLLM/figures/4_long_context/Relative_Frequency_Scaling.png}\label{fig:scaling-ratio}}
%   \caption{\centering Comparison of frequency remapping schemes.}
%   \label{fig:context-extension-suite}
% \end{figure*}

\begin{figure}
    \centering
    \includegraphics[width=1\linewidth]{4_Relative_Frequency_Scaling.png}
    \caption{Comparison of frequency remapping schemes.}
    \label{fig:scaling-ratio}
\end{figure}

As shown in Fig.~\ref{fig:scaling-ratio}, all concrete methods now differ \emph{only} by their choice of
$h(\theta_j)$ as a (possibly $j$‑dependent) function of $s$.
% Typical examples to be instantiated later are
% \[
%   \begin{aligned}
%     h_{\text{PI}}(\theta_j)   &= \theta_j/s,\\
%     h_{\text{NTK}}(\theta_j)  &= \theta_j\,s^{-\gamma_j},\\
%     h_{\text{NTP}}(\theta_j)  &= (1-\gamma)\,\theta_j/s+\gamma\,\theta_j,
%   \end{aligned}
% \]
% where $\gamma_j$ or $\gamma$ encode dimension‑specific heuristics.
% Plugging any $h(\theta_j)$ into
% Eq.~\eqref{eq:rope-extrapol-h} yields the full forward rule for
% $\mathbf q_m$ and $\mathbf k_n$ at the extended context length.

\subsubsection{Position Interpolation (PI)}
Position Interpolation rescales absolute positions when evaluating sequences longer than the pre‑training window~\cite{chen2023extending}.  Let $L_{\text{train}}$ and $L_{\text{target}}>L_{\text{train}}$ be the maximum lengths during training and inference. 

Each position $m$ is mapped to
$
m'=\alpha m,\qquad \alpha=(L_{\text{train}}-1)/(L_{\text{target}}-1),
\label{eq:pi_scale}
$ and $h_{\text{PI}}(\theta_j),m$ is defined as 
$$h_{\text{PI}}(\theta_j)=\alpha \theta_j.$$ 
\begin{definition}[Position Interpolation]\label{def:pi}
\begin{equation}
  f_W^{\text{PI}}(\mathbf x_m,m,\theta_j)
  \;=\;
  \mathbf R_{h_{\text{PI}}(\theta_j),m}\,
  W\,\mathbf x_m 
  % \;=\;
  % \mathbf R_{\alpha h(\theta_j),m}\
  % W\,\mathbf x_m
  \label{eq:pi-final}
\end{equation}
 This equals to setting the rotary angle $h(\theta_j)$ to $\alpha \theta_j$ in Eq.~\eqref{eq:rope-extrapol-h}.  This preserves the overall frequency range while re‑using trained low‑frequency dimensions. PI is easy to implement, but it compresses frequency bands uniformly and might affect the model performance in short problems.
\end{definition}

\subsubsection{Neural Tangent Kernel (NTK) -aware Interpolation}
Inspired by NTK theory~\cite{tancik2020fourier}, the neural networks have trouble learning high‑frequency information when the input dimension is low. This might be the reason why the perplexity of PI deteriorates a bit.

To solve the problem of losing high-frequency information, the NTK-aware~\cite{bloc97_ntk_rope_2023} interpolation is developed. The idea is to extend the low-frequency dimensions more and the high-frequency dimensions less.

With the scale factor
$
  s=L_{\text{target}}/L_{\text{train}}\;>1,
$
define
$
  b' \;=\; b \, s^{d/(d-2)},
  \label{eq:ntk-base}
$
and set the frequency‑remapping function to
\begin{equation}
  h_{\text{NTK}}(\theta_j)
  \;=\;
  (b')^{-\tfrac{2j}{d}},
  \qquad
  j=0,\dots,d/2 .
  \label{eq:ntk-h}
\end{equation}

\begin{definition}[NTK‑aware interpolation]\label{def:ntk-aware}
Using $g(m)=m$ and the map $h_{\text{NTK}}$ in
Eq.~\eqref{eq:rope-extrapol-h} yields
\begin{equation}
  f_W^{\text{NTK}}(\mathbf x_m,m,\theta_j)
  \;=\;
  \mathbf R_{h_{\text{NTK}}(\theta_j),m}\,
  W\,\mathbf x_m .
  \label{eq:ntk-final}
\end{equation}
\end{definition}

Equation~\eqref{eq:ntk-h} spreads the interpolation pressure across
dimensions as shown in Figure~\ref{fig:scaling-ratio}: low‑frequency dimensions (small $j$) receive a larger
scaling $h(\theta_j)/\theta_j$, while high‑frequency ones are
affected less, thereby preserving fine‑grained relative information
and reducing perplexity degradation at long context lengths. But it still blends interpolation/extrapolation in mid frequencies and affects performance under short context. 


\subsubsection{NTK‑by‑parts Interpolation}
\label{sec:ntk-by-parts}

Recall the wavelength of the $j$‑th rotary pair:
\[
  \lambda_j = 2\pi\,b^{\tfrac{2(j-1)}{d}},
  \qquad j = 1,\dots,\tfrac{d}{2}.
\]

Given a context size, some dimensions $j$ have a wavelength $\lambda_j$ longer than the maximum context length seen during the pretraining, so these dimensions' embeddings might not be trained on some part of the rotational domains. In such cases, these dimensions all have unique position pairs, making them behave like absolute position embeddings. 

For other dimensions, all tokens become closer to each other, so the dot product of two vectors rotated by a lesser amount is bigger. This has a negative impact on the LLM's ability to understand small local relationships between nearby embeddings. 

So for dimensions with wavelength smaller than context size, we do not interpolate; for wavelengths bigger than the context size, we only interpolate and avoid extrapolations, unlike the NTK-aware method; for dimensions in-between, we extrapolate like the NTK-ware method. 

\noindent\textbf{Relative ratio:}
Let $L$ be the pre‑training context length.  
Define the dimension‑specific ratio
\begin{equation}
  r(j)
  \;=\;
  \frac{L}{\lambda_j}
  \;=\;
  \frac{L}{2\pi\,b^{\tfrac{2(j-1)}{d}}}.
  \label{eq:ratio-r}
\end{equation}
Intuitively, $r(j)$ measures how many full rotary periods fit inside
the training window for dimension~$j$.

\noindent\textbf{Ramp (mask) function:}
Choose two hyper‑parameters $\alpha<\beta$ that delimit three
frequency regimes and define
\begin{equation}
  \gamma\!\bigl(r(j)\bigr)
  \;=\;
  \begin{cases}
    0, & r(j)<\alpha,\\[6pt]
    1, & r(j)>\beta,\\[6pt]
    \dfrac{r(j)-\alpha}{\beta-\alpha}, & \text{otherwise}.
  \end{cases}
  \label{eq:ramp}
\end{equation}
Thus $\gamma=0$ flags \emph{low‑frequency} dimensions
($\lambda_j\!>\!L/\alpha$), $\gamma=1$ flags
\emph{high‑frequency} dimensions
($\lambda_j\!<\!L/\beta$), and the linear middle section creates a
smooth transition.

\noindent\textbf{Frequency map:}
With scale factor $s\!=\!L_{\text{target}}/L_{\text{train}}>1$, set
\begin{equation}
  h_{\text{NTP}}(\theta_j)
  \;=\;
  \bigl(1-\gamma(r(j))\bigr)\,\frac{\theta_j}{s}
  \;+\;
  \gamma(r(j))\,\theta_j,
  \label{eq:h-ntp}
\end{equation}
which \emph{interpolates} low‑frequency dimensions by the factor
$1/s$ (PI‑like) while \emph{leaving} high‑frequency dimensions
unchanged; mid‑range dimensions receive a weighted mix, as shown in Figure~\ref{fig:scaling-ratio}. 

\begin{definition}[NTK‑by‑parts interpolation]\label{def:ntk-by-parts}
Using $g(m)=m$ and the map $h_{\text{NTP}}$ from
Eq.~\eqref{eq:h-ntp} in the generic template
Eq.~\eqref{eq:rope-extrapol-h} yields
\begin{equation}
  f_W^{\text{NTP}}(\mathbf x_m,m,\theta_j)
  \;=\;
  \mathbf R_{\,m\,h_{\text{NTP}}(\theta_j)}\,W\,\mathbf x_m .
  \label{eq:ntp-final}
\end{equation}
\end{definition}

Recommended defaults for LLaMA‑family models are
$\alpha=1$ and $\beta=32$; these can be tuned per architecture
and target length.

\subsubsection{YaRN}
\label{sec:yarn}

YaRN augments the \textit{NTK‑by‑parts} frequency map
$h_{\text{NTP}}(\theta_j)$ (Sec.~\ref{sec:ntk-by-parts}) with a
\emph{temperature} $t>0$ applied to the attention logits:
\begin{equation}
  \operatorname{softmax}\!\Bigl(
     \tfrac{\mathbf q_m^{\top}\mathbf k_n}{\,t\sqrt{d}\,}
  \Bigr),
  \label{eq:yarn-softmax}
\end{equation}
where $d$ is the model dimension.
Empirically, the following scale‑dependent rule provides near‑optimal
perplexity across LLaMA‑family models:
\begin{equation}
  \sqrt{\smash[b]{\tfrac1t}}
  \;=\;
  \alpha(s)
  \;=\;
  0.1\,\ln s + 1,
  \qquad
  s=\frac{L_{\text{target}}}{L_{\text{train}}}>1 .
  \label{eq:yarn-alpha}
\end{equation}
Hence $t=\alpha(s)^{-2}$.

\noindent\textbf{Implementation via length scaling:}
Because RoPE can be viewed as a bank of $2\times2$ rotation matrices,
Eq.~\eqref{eq:yarn-softmax} is equivalent to scaling the \emph{RoPE
embeddings themselves}:
\[
  \mathbf q_m^{\text{YaRN}}
    \;=\;
    \alpha(s)\,
    \mathbf q_m^{\text{NTP}},
  \qquad
  \mathbf k_n^{\text{YaRN}}
    \;=\;
    \alpha(s)\,
    \mathbf k_n^{\text{NTP}},
\]
where $\mathbf q^{\text{NTP}}$ and $\mathbf k^{\text{NTP}}$ use the
frequency remapping $h_{\text{NTP}}(\theta_j)$ from
Eq.~\eqref{eq:h-ntp}.  
No changes to the attention kernel are required; RoPE embeddings are
scaled once and cached for all forward passes.

\begin{definition}[YaRN]\label{def:yarn}
Combine
\[
  g(m)=m,\quad
  h(\theta_j)=h_{\text{NTP}}(\theta_j),\quad
  \text{and}\quad
  \alpha(s)\text{ from Eq.~\eqref{eq:yarn-alpha}}
\]
to obtain
\begin{equation}
  f_W^{\text{YaRN}}(\mathbf x_m,m,\theta_j)
  \;=\;
  \alpha(s)\,
  \mathbf R_{\,m\,h_{\text{NTP}}(\theta_j)}\,
  W\,\mathbf x_m.
  \label{eq:yarn-final}
\end{equation}
\end{definition}

Recommended values for Eq.~\eqref{eq:yarn-alpha} were obtained by
fitting $\alpha(s)$ on LLaMA‑7B/13B/65B for scale factors
$s\!\in\![1,512]$ without fine‑tuning; the same rule transfers well
to Llama‑2 models (7B, 13B, 70B).\footnote{%
  For alternative kernels such as Flash‑Attention 2 the same
  embedding‑side scaling applies because the softmax temperature in
  Eq.~\eqref{eq:yarn-softmax} is absorbed into $\alpha(s)$.} 
YaRN proves itself in many models, but it still requires extra hyperparameter tuning for $\alpha(s)$. 

\subsubsection{LongRoPE}
\label{sec:longrope}

Unlike PI, NTK-aware or YaRN—which apply a \emph{global}
dimension-wise rule—\textbf{LongRoPE}~\cite{ding2024longrope}  learns a \emph{token-dependent}
scale for every rotary sub-plane.  
Let \(P=\{p_0,p_1,\dots,p_K\}\) be a set of \emph{anchor} positions
with \(0=p_0<\dots<p_K=L_{\text{target}}\).
For each dimension \(j\in[1,\tfrac{d}{2}]\) and every interval
\([p_k,p_{k+1})\) we introduce a learnable scale
\(\sigma_{j,k}>0\).
The frequency map therefore becomes a piece-wise constant function
\begin{equation}
  h_{\text{LR}}\bigl(\theta_j,m\bigr)
  \;=\;
  \sigma_{j,k}\,\theta_j,
  \qquad
  p_k\le m < p_{k+1}.
  \label{eq:h-longrope}
\end{equation}

\noindent\textbf{Optimisation objective:}
Given the pretrained attention logits
\(z_{m,n}=\mathbf q_m^{\top}\mathbf k_n\)
(\(L_{\text{train}}\!\times\!L_{\text{train}}\) region),
LongRoPE selects \(\{\sigma_{j,k}\}\) by minimising the discrepancy
between the original logits and those produced with the scaled RoPE:
\begin{equation}
  \min_{\{\sigma_{j,k}\}}
  \; \mathbb{E}_{(m,n)\sim\mathcal D}
  \bigl\|
    z_{m,n}
    -\tilde z_{m,n}
  \bigr\|_2^2,
  \label{eq:longrope-loss}
\end{equation}
where \(\tilde z_{m,n}\) is obtained from
Eq.~\eqref{eq:rope-extrapol-h} with
\(h(\theta_j)=h_{\text{LR}}(\theta_j,m)\).
The sampling distribution \(\mathcal D\) mixes pairs
\((m,n)\) \emph{within} pre-training window
(\(m,n<L_{\text{train}}\)) and \emph{out-of-window}
pairs (\(m\) or \(n\ge L_{\text{train}}\)); this encourages model
to preserve short-range behaviour while smoothly extending to
\(L_{\text{target}}\).

\noindent\textbf{Search procedure:}
The resulting multidimensional, non-uniform interpolation is cast as a
search problem over the discrete grid \(P\):
LongRoPE performs coordinate-wise optimisation (or Bayesian search)
until convergence, then caches the \(\sigma_{j,k}\) table for use at
inference time.  
In practice \(K\!\ll\!L_{\text{target}}\) (e.g.\ anchors every 512
tokens), so memory overhead is negligible.

LongRoPE is more capable since it learns the parameter \(\sigma_{j,k}>0\) from the data and preserves original short-length quality, but it requires high-quality data for the adjustment process, and has a complex pipeline. 

\begin{table*}[t]
\centering
\small
\caption{Context extension at a glance. Tags: \textit{Context Extension} = frequency remapping \(h(\cdot)\) plus optional length/temperature scaling (e.g., YaRN); Arch.\ = architectural assists.}
\label{tab:ctx-ext-at-a-glance}
\resizebox{\textwidth}{!}{
\begin{tabular}{p{3.5cm}lllp{4.2cm}}
\toprule
Model & Max Len & Context Extension & Arch. & Datasets for long-context \\
\midrule
Nemotron-4 15B~\cite{parmar2024nemotron} & 4k & RoPE (native) & — & N/A \\
Gemma 1 (7B) & 8k & RoPE (native) & — & N/A \\
InternLM2~\cite{cai2024internlm2} & 32k & ABF & — & Book, patent, paper, CC. \\
MiniCPM-128K~\cite{hu2024minicpm} & 128k & ABF + NTK-aware & — & — \\
Phi-3-Mini~\cite{abdin2024phi3} & 128k & LongRoPE~\cite{ding2024longrope} & BlocksparseAttn & — \\
DeepSeek-V2~\cite{liu2024deepseek} & 128k & YaRN & MLA, MoE & — \\
Zamba-v1~\cite{glorioso2024zamba} & 4k & RoPE (native) & Hybrid ssm/global 6:1 & — \\
MAP-Neo~\cite{Zhang2024MAPNeoHC} & 8k & RoPE (native) & — & — \\
Nemotron-4 340B~\cite{adler2024nemotron} & 4k & RoPE (native) & — & N/A \\
AFM~\cite{gunter2024apple} & 32k & ABF & — & — \\
Llama-3.1~\cite{Dubey2024TheL3} & 128k & ABF + YaRN~\cite{peng2023yarn} & — & — \\
Gemma 2~\cite{team2024gemma} & 8k & RoPE (native) & Interleave local/global 1:1 & N/A \\
Granite 3.0~\cite{granite2024granite} & 4k & RoPE (native) & — & N/A \\
Hunyuan-Large~\cite{sun2024hunyuan} & 256k & ABF & MoE & 25\% Book, Code \\
Yi-Lightning~\cite{wake2024yi} & 64k & ABF & Interleave local/global 3:1 & — \\
Qwen-2.5/2.5-Turbo~\cite{Yang2024Qwen25TR} & 128k/1M & ABF +YaRN +DCA  & SparseAttn (Minference) & — \\
DeepSeek-V3~\cite{DeepSeekAI2024DeepSeekV3TR} & 128k & YaRN & MLA, MoE & — \\
OLMo-2~\cite{olmo20242} & 4k & RoPE (native) & — & N/A \\
MiniMax-01~\cite{li2025minimax} & 4M & ABF & Interleave LightningAttn/global 7:1 & +10\% long QA late-stage \\
Qwen-2.5-1M~\cite{yang2025qwen2} & 1M & ABF +YaRN +DCA & SparseAttn (Minference) & 40\% long +synthetic tasks \\
SmolLM2~\cite{allal2025smollm2} & 8k & ABF & — & 40\% long data mix \\
Gemma 3~\cite{team2025gemma} & 128k & ABF + PI & Interleave local/global 5:1 & — \\
Llama-4 Scout~\cite{meta2025llama4} & 10M & NoPE/iRoPE + temp scale & Interleave local/global & — \\
Phi-4~\cite{abdin2024phi} & 16k & ABF & — & +30\% curated long data \\
MiMo-7B~\cite{xiaomi2025mimo} & 32k & ABF & — & — \\
Qwen-3~\cite{yang2025qwen3} & 128k & ABF +YaRN +DCA & — & 75\% long data \\
PANGU PRO MoE~\cite{tang2025pangu} & 32k & — & MoE & STEM/code-heavy later stages \\
MiniCPM-4 8B~\cite{team2025minicpm4} & 128k & LongRoPE +YaRN & SparseAttn (InfLLM v2) & — \\
Hunyuan-A13B & 256k & NTK-aware & MoE & — \\
SmolLM3~\cite{bakouch2025smollm3} & 128k & ABF +YaRN & Interleave NoPE & math, code, reasoning-heavy \\
\bottomrule
\end{tabular}
}
\end{table*}


% ===============================================================
% A Unified View of Context Extension (frequency maps h(θ_j))
% ===============================================================
\subsubsection{A Unified View of Context Extension}
\label{sec:unified-context-ext}

For each rotary sub-plane \(j=1,\dots,\tfrac{d}{2}\) with base
frequency \(\theta_j\) (wavelength \(\lambda_j=2\pi/\theta_j\)),
and scale factor \(s=L_{\text{target}}/L_{\text{train}}>1\),
the following frequency maps \(h(\theta_j)\) instantiate the generic
template in Eq.~\eqref{eq:rope-extrapol-h}, as show in Table~\ref{tab:unified-h-theta}.

\medskip
\noindent\textbf{(i) Position Interpolation (PI).}
\begin{equation}
  h_{\text{PI}}(\theta_j)=\frac{\theta_j}{s}.
  \label{eq:h-pi}
\end{equation}

\medskip
\noindent\textbf{(ii) NTK-aware Scaling.}
\begin{equation}
  h_{\text{NTK}}(\theta_j)
  = \theta_j\,s^{-\tfrac{2(j-1)}{\,d-2\,}}.
  \label{eq:h-ntk}
\end{equation}

\medskip
\noindent\textbf{(iii) NTK-by-parts.}
Define the ratio
\begin{equation}
  r(j)=\frac{L_{\text{train}}}{\lambda_j}
      =\frac{L_{\text{train}}\theta_j}{2\pi},
  \label{eq:rj}
\end{equation}
and a ramp function with hyper-parameters \(\alpha<\beta\):
\begin{equation}
  \gamma(r)=
  \begin{cases}
    0, & r<\alpha,\\[4pt]
    1, & r>\beta,\\[4pt]
    \dfrac{r-\alpha}{\beta-\alpha}, & \text{otherwise}.
  \end{cases}
  \label{eq:gamma}
\end{equation}
Then
\begin{equation}
  h_{\text{NTP}}(\theta_j)
    = \bigl(1-\gamma(r(j))\bigr)\,\frac{\theta_j}{s}
      + \gamma(r(j))\,\theta_j.
  \label{eq:h-ntp-unified}
\end{equation}

\medskip
\noindent\textbf{(iv) YaRN.}
YaRN uses the NTK-by-parts frequency map unchanged:
\begin{equation}
  h_{\text{YaRN}}(\theta_j)=h_{\text{NTP}}(\theta_j).
  \label{eq:h-yarn}
\end{equation}
In addition, YaRN applies a length-dependent scaling to the resulting
queries/keys (or equivalently a softmax temperature):
\begin{equation}
  \alpha(s)=0.1\ln s + 1,
  \qquad
  t=\alpha(s)^{-2}.
  \label{eq:yarn-alpha-unified}
\end{equation}
(See Sec.~\ref{sec:yarn} for details.)

\medskip
\noindent\textbf{(v) LongRoPE.}
LongRoPE introduces non-uniform rescaling across RoPE dimensions and (optionally) across token positions. 
Let \(s_j \in [1,s]\) be a per-dimension rescale factor and let \(\tau\in\mathbb{N}\) denote the number of initial tokens that keep the original RoPE (no interpolation). 
Then the frequency map is
\begin{equation}
  h_{\text{LongRoPE}}(\theta_j;m)
  \;=\;
  \begin{cases}
    \theta_j, & m<\tau,\\[4pt]
    \dfrac{\theta_j}{\,s_j\,}, & m\ge \tau,
  \end{cases}
  \label{eq:h-longrope-unified}
\end{equation}
which recovers PI when \(s_j\equiv s\) and \(\tau=0\), and contains NTK/YaRN-like non-uniformity as special cases when \(s_j\) varies with \(j\). 
In our generic template \eqref{eq:rope-extrapol-h}, this corresponds to applying the rotation with an \(m\)- and \(j\)-dependent frequency,
\(
\mathbf R_{h_{\text{LongRoPE}}(\theta_j;m),\,m}.
\)
(See \cite{ding2024longrope} for the search-based procedure that selects \(\{s_j\}\) and \(\tau\).)

% \subsection{Context Extension in Popular LLMs}
% In this section, we discuss the long-context extension part of the mainstream LLMs. 

% \begin{table*}[h]
%     \centering
%     \caption{Long Context Extension in the state-of-the-art Models}
%     \small            % keeps the wide table on the page
%     \begin{tabular}{lllp{4.5cm}}
%     \hline
%     \textbf{Model} & \textbf{Context} & \textbf{Tech for context extension} & \textbf{Datasets for long-context} \\ \hline
%     Nemotron-4 (15B)    & 4k   & RoPE                            &  \\
%     Gemma 1 (7B)        & 8K   & RoPE                            &  \\
%     InternLM2 (7B)      & 200k & ABF                             & Book, patent, paper, CC. \\
%     Phi-3-Mini (3.8B)   & 128K & LongRoPE two-stage    & Undisclosed \\
%     DeepSeek-V2 (21B/236B MoE) & 128K & MLA, YaRN                       & Undisclosed \\
%     Nemotron-4 (340B)   & 4K   & RoPE                            &  \\
%     Llama-3/3.1 (3B,70B,405B)  & 128K & ABF, YaRN                       & Undisclosed \\
%     Gemma 2 (2B,7B)     & 8K   & Interleave 50\% local, RoPE & \\
%     Granite 3.0 (8B)    & 4K   & RoPE                            & \\
%     Hunyuan-Large (52B/389B MoE) & 256K & ABF two-stage                   & 25\% Book, Code 75\% Normal \\
%     Yi-Lightning        & 64K  & Interleave 75\% local, ABF & Undisclosed \\
%     Qwen 2.5 (0.5B-72B) & 128K & ABF, YaRN, DCA                  & Undisclosed \\
%     DeepSeek-V3 (37B/671B MoE) & 128K & MLA, YaRN two-stage            & Undisclosed \\
%     OLMo 2 (7B,13B)     & 4K   & RoPE                            & \\
%     MiniMax-01 (45.9B/456B MoE) & 4M   & Hybrid Lightning Attn, ABF & Upsampled long data, add 10\% high quality long QA. \\
%     Qwen 2.5-1M (7B/14B) & 1M   & ABF, DCA, Sparse Attention           & Natural, Fill-in-the-middle, retrieval, and paragraph reorder. \\
%     SmolLM2 (1.7B)       & 8K   & ABF                             & 40\% long, 60\% short \\
%     Gemma 3 (1B-27B)    & 128K & Interleave $83\%$ local, ABF, PI & Undisclosed \\
%     Llama 4 Scout (17B/109B MoE) & 10M  & Interleave local, NoPE  & Undisclosed \\ 
%     Phi-4 (14B)         & 16K  & ABF                             & Upsampled academic, books, code, with 30\% longer data \\
%     MiMo (7B)           & 32K  & ABF three-stage                 & add $10\%$ synthetic data for math, code, and writing. \\ 
%     Qwen 3 (235B)       & 128K & ABF, YaRN, DCA                  & 25\% length in (4k,16k), 75\% length in (16k,32k) \\
%     PANGU PRO MoE (16B/72B MoE) & 32K  & Undisclosed                     & Undisclosed \\
%     MiniCPM-4 (8B)      & 128K & InfLLM v2, LongRoPE, YaRN       & Undisclosed \\
%     % MiniMax-M1          & 1M   & Hybrid Arc, Lightning Attention & 3 stage datasets of 360B token, with 10\% high-quality long QA data \\
%     Hunyuan A13B        & 256K & NTK-aware                       & Undisclosed \\
%     SmolLM3 (3B)        & 128K & ABF, NoPE, YaRN                 & Upsampled math, code, and reasoning \\
%     \hline
%     \end{tabular}
%     \label{tab:long_context_llms}
% \end{table*}

% context, technique, data, evaluation benchmarks, 

% The Nemotron-4 15B~\cite{parmar2024nemotron} is naively pretrained with a context window of 4k. It is continued pretrained on two distinct distributions. The first distribution is the pretraining dataset with a larger sampling weight on the higher-quality sources. The second distribution uses a small number of benchmark-style alignment examples and up-weights the data sources with low performance, to better prepare the model for downstream evaluations. 

% Gemma 1 (7B) – The original Gemma-7B is trained natively at 8K with standard RoPE. 

% The InternLM2~\cite{cai2024internlm2} is an open-source model that explicitly focuses on long-context modeling. It supports context lengths up to $32K$ tokens after pre-training, and was demonstrated to be effective on a 200K-token Needle In A Haystack test. InternLM2’s architecture uses Grouped-Query Attention (GQA) to reduce memory usage per head. The model was first trained on 4K-token sequences, then transitioned to 32K-token pre-training in a second phase, with ~9\% of training steps using mixed long sequences, half of which are longer than $4K$. During this long-context phase, the RoPE base frequency was expanded from $50K$ to $1M$ to stretch the positional embedding to $32k$. High-quality long-context data like concatenated code files from repositories, lengthy books and papers was curated and filtered for this phase. 

% The Phi-3~\cite{abdin2024phi3} model series, supports up to $128K$ context length. The authors proposed a blocksparse attention module. For each attention head, the blocksparse attention enforces different sparsity patterns on KV caches, which ensures all tokens are attended to by different heads under a given choice of sparsity. Specifically, the phi-3-mini-128K uses the LongRope~\cite{ding2024longrope} technique to extend the model from $4K$ to $128K$. 

% The DeepSeek-V2~\cite{liu2024deepseek} is an efficient MoE model, with 236B total parameters and 21B activation parameters per token. It supports a context length of $128K$. The authors proposed a novel MLA architecture, which supports efficient inference under long context, since it compresses KV-cache significantly. During context extension of the model, YaRN~\cite{peng2023yarn} is used to extend the model from $4K$ to $128K$. YaRN is applied to the decoupled shared key which uses RoPE. For YaRN, the authors set the scale $s$ to $40$, $\alpha$ to $1$, $\beta$ to $32$, and the target max length to $160K$. The author also adjusted the length scaling factor to modulate the attention entropy. The factor $\sqrt{t}$ is set to $\sqrt{t}=0.0707 \ln{s} + 1$ to minimize the perplexity. During long context extension, the model is trained for $1000$ steps, with a sequence length of $32K$ and a batch size of $576$. 

% The Nemotron-4 340B~\cite{adler2024nemotron} is a model pretrained with a $4K$ context window. There is no long context extension in the mid-training phase. 

% The Llama-3.1~\cite{Dubey2024TheL3}, the improved version of Llama 3, introduced $128K$ context support. Meta achieved this by training Llama-3 in stages: initially on ~8K sequences, then using a long-context pre-training phase up to 32K, and finally employing positional extrapolation to reach 128K at inference. Indeed, Llama-3.1 continues to use Rotary Positional Embeddings, but with an adjusted base and YaRN~\cite{peng2023yarn} to extend beyond the trained length. The architecture also uses Grouped-Query Attention (GQA) to manage memory and speed with long inputs. Llama-3.1 was trained on a massive 15T-token dataset, so plenty of long texts (books, code, etc.) were included implicitly. The long-context pretraining and careful RoPE scaling resulted in Llama-3.1 nearly perfectly identifying information in 128k-token tests. 

% The Gemma 2~\cite{team2024gemma} models have a context length of $8K$. To speed up, the model alternates between global attention and local attention in every other layer. The global attention layers span a range of $8K$ tokens, while the local attention layers span a range of $4K$ tokens. 

% The Granite 3.0~\cite{granite2024granite} is a set of lightweight foundation models from 400M to 8B active parameters. It is pre-trained naively with a sequence length of $4K$, and there is no long context extension. 

% The Hunyuan-Large~\cite{sun2024hunyuan} is an MoE LLM with 389 billion total parameters and 52 billion activation parameters, and it can support a context length of $256K$ tokens. After the learning rate annealing phase, the model is extended from $4K$ to $32K$, then from $32K$ to $256K$, with $\theta$ set to 1 billion using the ABF~\cite{xiong2023effective} techniques. The dataset used has $25\%$ natural long context data from books and codes, and $75\%$ normal-length pretraining data. In these two long context extension stages, about 10 billion tokens are used. The RULER~\cite{hsieh2024ruler} and LV-Eval are used to evaluate the long context capabilities of the model. 

% The Yi-Lightning~\cite{wake2024yi} is an MoE LLM with a context length of $64K$. The model uses hybrid attention blocks which combine three sliding-window-attention layers with one full attention layer. After a fast learning rate decreasing phase, the model is progressively extended from $8K$ to $16K$, $16K$ to $32K$, $32K$ to $64K$ with the ABF~\cite{xiong2023effective} techniques. 

% The Qwen 2.5~\cite{Yang2024Qwen25TR} model series supports a context length of up to $128K$, and the Qwen2.5-Turbo model series supports a context length of $1M$ tokens. The Qwen2.5 has a two-phase pretraining approach. After the initial $4k$ pretraining stage, the model is extended to $32k$ during the second pretraining stage with the ABF techniques, which adjust the RoPE $\theta$ from $10k$ to $1M$. The Qwen2.5-Turbo has a four-stage pretraining, and the context length is extended to $32k$, $64k$, $128k$, $256k$ in the four stages. The final RoPE $\theta$ is set to $10M$. At each stage, the authors used $40\%$ long context dataset and $60\%$ short dataset. YARN~\cite{peng2023yarn} and DCA are used during inference, achieving a four-fold increase in sequence-length capacity. The Qwen2.5-1M model is evaluated with the RULER~\cite{hsieh2024ruler}, the LV-Eval~\cite{yuan2024lv}, the Longbench-chat~\cite{bai2024longalign}. 

% Granite 3.1 language models extend the context length of Granite 3.0 language models from 4K to 128K using a progressive training strategy by increasing the supported context length in increments while adjusting RoPE theta until the models successfully adapt to the desired length of 128K. This long-context pre-training stage was performed using approximately 500B tokens. 

% The DeepSeek-V3~\cite{DeepSeekAI2024DeepSeekV3TR} used MLA and MoE, and it has improved long-context performance. After the first phase of pretraining, the authors applied YaRN~\cite{peng2023yarn} for context extension and performed two additional training phases. Both phases have 1000 steps to expand the context window of the model from 4k to 32k and then to 128k. It used the same YaRN configuration with DeepSeek-V2, and it is applied to the decoupled shared $\bf{k}_t^R$. The hyper-parameters remain the same for two phrases, with the scale $s=40$, $\alpha=15$, $\beta=32$, and the scaling factor $\sqrt{t}=0.1 \ln s+1$. The first phase is set to 32k with batch size $1920$, the second phase is set to 128k with batch size $480$. The learning rate of both phases is $7.3e-6$, and it is the final learning rate from the pretraining stage. The Deepseek-v3 evaluated its long context performance on various benchmarks, including the DROP, the FRAMES, the LongBench v2, and the Needle In A Haystack(NIAH). 

% The OLMo 2~\cite{olmo20242} is a fully open language model with a $4K$ context length. The model is pretrained with a $4K$ context window, and there is no specific mid-training phase to further extend the context length. 

% The MiniMax-01 model~\cite{li2025minimax} features a context length of 4 million tokens.  It introduced the lightning attention, an I/O aware linear attention. To save kv-cache, it introduced a hybrid architecture, where one transformer block follows seven transformer blocks with lightning attention. The data used in the long context extension phase is generated by upsampling long context data and preserving the distribution in critical domains. The author mixes in 10\% high-quality question answer data during the last 20\% of training cycles in each stage. To ensure a smooth data distribution transition, the author uses linear interpolation of source-specific weights in the data transition phase. The long context ability of this model is evaluated on several benchmarks, including the Multi-Round Needle-In-A-Haystack(MR-NIAH), the RULER~\cite{hsieh2024ruler}, the LongBench-v2 and the Machine Translation from One Book(MTOB). 

% Qwen 2.5-1M~\cite{yang2025qwen2} supports up to 1 million tokens in the context window. During pretraining, the model is exposed to both natural and synthetic long text data. The natural data includes Common Crawl, arXiv, books and code repositories. The synthetic data includes Fill-in-the-middle tasks, keyword-based and position-based retrieval tasks, and paragraph reordering tasks, to expose the model to enough tasks requiring long dependencies. The context length is progressively extended in five stages, from $4K$ to $256k$, along with the Adaptive Base Frequenciy(ABF)~\cite{xiong2023effective}. In the first two stages, the model is extended from $4K$ to $32k$, adjusting the RoPE $\theta$ from $10K$ to $1M$. In the next three stages, the model is expanded to $64K$, $128K$ and $256k$, with the RoPE base frequencies set to $1M$, $5M$, and $10M$, respectively. The Qwen2.5-1M model is evaluated with the RULER~\cite{hsieh2024ruler}, the LV-Eval~\cite{yuan2024lv}, the Longbench-chat~\cite{bai2024longalign}. 

% The SmolLM2~\cite{allal2025smollm2} is a 1.7B small model with a context length of 8k. During the long context extension phase, the RoPE $\theta$ is increased from 10k to 130k. The data used for long context extension consists of 40\% long context documents from DCLM(10\%), FineWeb-Edu(10\%), and the books subset of Dolma (20\%), and the remaining 60\% from the last pretraining stage. The model is evaluated with Needle in the Haystack (NIAH) and the HELMET benchmark.

% The Gemma 3~\cite{team2025gemma} is a multimodal LLM that has a 128k context length. To achieve a long context, the model reduces the kv-cache by using interleaving global and local layers at the ratio of 1:5. The model is pre-trained with 32K sequences. During the context extension phase, the RoPE $\theta$ in the global layers is increased from 10k to 1M, while the $\theta$ in the local layers is unchanged. To further extend the context length, the authors also used the Position Interpolation~\cite{chen2023extending}. % Did not mentioned long context dataset. 
% The models are evaluated with the RULER~\cite{hsieh2024ruler} and MRCR~\cite{vodrahalli2024michelangelo}. 

% The Llama 4~\cite{meta2025llama4} scout supports a context length of up to 10M. Llama 4 Scout uses the iRoPE architecture with interleaved global and local layers. In the global layers, it uses no position embedding (NoPE). The model is pre-trained with 256k context length. Additionally, the authors use inference-time temperature scaling of attention to enhance length generalization. The model is evaluated on the Needle In A Haystack(NIAH) benchmark. 

% The Phi-4~\cite{abdin2024phi} has a context length of 16k. During the mid-training stage, the authors increased the base frequency of RoPE to 250K with ABF~\cite{xiong2023effective}, and continued pretraining the model on datasets with longer context. They further filtered and up-weighted the high-quality non-synthetic datasets, such as academic, books, and code data. The authors also created new synthetic datasets longer than 4k. The final data mix includes 30\% newly curated longer context data, and 70\% of the original pretraining dataset. Evaluations used included the RULER~\cite{hsieh2024ruler} and the HELMET. Specifically, the authors reported the sub-tasks of Recall, RAG, Re-rank, ICL, QA, Summ. 

% The MiMo-7B~\cite{xiaomi2025mimo} uses a three-stage pretraining schema. In the first two stages, the model is pretrained with a context length of 8K. In the third stage, the authors extended the context length of the model from 8K to 32K, and they increased the RoPE $\theta$ from $10,000$ to $64,000$ in the third stage. In the third stage, the authors incorporated $10\%$ synthetic responses for mathematics, code, and creative writing queries in the dataset, to boost the model's capability to solve complex tasks. The model is evaluated with the RULER~\cite{hsieh2024ruler}.

% The Qwen 3~\cite{yang2025qwen3} has a context length of 128k. It has three pretraining stages: the general stage (S1), the reasoning stage (S2), and the Long Context Stage (S3). In the long context stage, the authors extend the context length to 32k with a mixture corpus of 25\% short (4k-16k length) data and 75\% long (16k-32k) data. The RoPE $\theta$ is increased from $10,000$ to $1,000,000$ using the ABF~\cite{xiong2023effective} techniques. After extending the model to 32k, the model is further extended to 128k with YARN~\cite{peng2023yarn} and Dual Chunk Attention techniques. These training-free techniques increased the context length of the model by four times. The model is evaluated with the RULER~\cite{hsieh2024ruler}.

% The PANGU PRO MoE~\cite{tang2025pangu} is a MoE model with 72B total parameters and 16B active parameters. trained with three phases. After the first general pretraining phase, the authors trained the model with a significantly higher proportion of high-quality STEM data, code, and synthetic data in the second reasoning phase, under 32k context length. In the third annealing stage, the authors bridged the gap between the pre-training phase and the alignment phase by including 20\% instruction data and 18\% advanced-level STEM education data, under 32k context length. 

% The MiniCPM-4 8B~\cite{team2025minicpm4} is an efficient small model trained for End Devices. The author proposed InfLLM v2, a trainable sparse attention mechanism which has up to 81\% sparsity. After pretraining, the model underwent a mid-training phase with 1T tokens. The author extended the model from $4K$ to $32K$ with 20 billion tokens using LongRoPE~\cite{ding2024longrope}. Although the model is trained only on $32K$ context length, the model is further extended to $128K$ with YaRN~\cite{peng2023yarn}. The long-context ability of this model is evaluated with RULER-NIAH~\cite{hsieh2024ruler}. 

% % Minimax-M1 reused minimax-01 pretrain checkpoint. 

% The Hunyuan-A13B is an MoE model with 80B total parameters and 13B active parameters for each token. After pretraining with 4k context length, the model went through a fast annealing stage which extended its context length to 8k. In the context extension stage, the model went through two sequential stages to expand its context length to 32k and then to 256k. Both phases use the NTK-aware~\cite{bloc97_ntk_rope_2023} position encoding with the alpha value set to 50 for 32k and 1000 for 256k respectively. 

% The SmolLM3~\cite{bakouch2025smollm3} is a 3B powerful multilingual long context reasoner model with $128K$ context length. The model is pre-trained with five different stages with $11.2T$ tokens. The model implements NoPE, which selectively removes rotary position embeddings from every 4th layer. This approach improves long context performance without affecting short context capabilities.
% %After the first pretraining with $8T$ tokens, the model is further pretrained with $2T$ of high-quality tokens. In the third stage the model is trained with a linear decay learning rate on $1.1T$ of high-quality code, math, instruction, and reasoning data. 
% The fourth stage is long-context extension, the model is progressively extended from $4K$ to $32K$ with $\theta=1.5M$, then from $32K$ to $64K$ with $\theta=5M$, both stages are with upsampled math, code, and reasoning data. During inference, the model is further extended with YARN~\cite{peng2023yarn} to $128K$ context length. The model is evaluated using RULER and HELMET benchmarks. 

\begin{table*}[h]
  \centering
  \caption{Pre-training benchmarks by category.}
  \label{tab:std-benchmarks}
  \begin{tabular}{p{2.2cm}p{11cm}}
    \toprule
    \textbf{Task Category} & \textbf{Benchmarks} \\
    \midrule
    General & MMLU~\cite{hendrycks2021measuringmassivemultitasklanguage} (5-shot),
        MMLU-Pro~\cite{wang2024mmluprorobustchallengingmultitask} (5-shot, CoT),
        ARC ~\cite{clark2018thinksolvedquestionanswering} (0-shot), 
        HellaSwag~\cite{zellers2019hellaswagmachinereallyfinish} (0-shot),
        BBH~\cite{suzgun2022challengingbigbenchtaskschainofthought} (3-shot, CoT), 
        WinoGrande~\cite{sakaguchi2019winograndeadversarialwinogradschema} (0-shot),
        PIQA~\cite{bisk2019piqareasoningphysicalcommonsense} (0-shot) \\[2pt]
    \midrule
    Math \& STEM & GPQA~\cite{rein2023gpqagraduatelevelgoogleproofqa} (5-shot, CoT),
        MATH~\cite{hendrycks2021measuringmathematicalproblemsolving} (4-shot, CoT),
        GSM8K~\cite{cobbe2021gsm8k} (4-shot, CoT) \\[2pt]
    \midrule
    Coding & HumanEval~\cite{chen2021evaluating} (0-shot), 
        MBPP~\cite{austin2021programsynthesislargelanguage} (3-shot),  
        MultiPL-E~\cite{cassano2022multiplescalableextensibleapproach} (0-shot; Python, C++, Java, PHP, TypeScript, C\#, Bash, JavaScript) \\[2pt]
    \midrule
    Multilingual & MGSM~\cite{shi2022languagemodelsmultilingualchainofthought} (8-shot, CoT),
        Flores-101~\cite{goyal2021flores101evaluationbenchmarklowresource} (5-shot) \\[2pt]
    \midrule
    Long context & LongBench~\cite{bai2024longbenchbilingualmultitaskbenchmark}, 
        Ruler~\cite{hsieh2024rulerwhatsrealcontext} \\[2pt]
    \bottomrule
  \end{tabular}
\end{table*}

\subsection{Context Extension in Popular LLMs}
\label{sec:ctx-ext-popular}

Most systems scale context using three levers. Firstly, \emph{frequency remapping} of RoPE via \(h(\theta_j;s)\), such as ABF, PI, NTK, NTP and LongRoPE mentioned in Sec.~\ref{sec:unified-context-ext}. Secondly, optional \emph{temperature scaling} \(\alpha(s)\) at the attention logits, e.g., YaRN~\cite{peng2023yarn}. Thirdly, \emph{architecture assists} that reduce KV cost, e.g., the interleaved global and local attention, GQA, MLA, and MoE. 
In practice, windows are grown progressively \(4\text{k}\!\to\!32\text{k}\!\to\!128\text{k}\) with a mix of long and short tokens and occasional synthetic long-dependency tasks. We summarize the techniques used in many LLMs in Table~\ref{tab:ctx-ext-at-a-glance}.

\medskip

\noindent\textbf{No context extension (native RoPE).} 
The Nemotron-4 15B~\cite{parmar2024nemotron} uses native RoPE at 4K, and subsequent tuning reweights data but keeps the window unchanged. 
The Gemma 1 (7B) employs native RoPE at 8K. 
The Nemotron-4 340B is pretrained at 4K with native RoPE and does not include a long-context continuation phase. 
The Gemma 2~\cite{team2024gemma} fixed 8K with native RoPE, and it interleaves global and local layers roughly 1{:}1 to reduce KV cost.
The Granite 3.0~\cite{granite2024granite} family are trained at 4K with native RoPE. 
The OLMo 2~\cite{olmo20242} is a 4K model using native RoPE. 

\medskip

\noindent\textbf{Adaptive Base Frequency (ABF).}
The InternLM2~\cite{cai2024internlm2} extends to 32K via ABF~\cite{xiong2023effective} (RoPE base $50\mathrm{K}\!\to\!1\mathrm{M}$), GQA, and $\sim9\%$ mixed long sequences. Its training data includes book, patent, paper, and CC long sequences.
The Hunyuan-Large~\cite{sun2024hunyuan} uses staged 4K$\to$32K$\to$256K with ABF and a MoE backbone. About 25\% long natural data is included, and evaluations cover RULER/LV-Eval.
The Yi-Lightning~\cite{wake2024yi} uses progressive ABF from 8K to 64K. It combines sliding-window and full attention with an interleave of around 3{:}1.
The SmolLM2~\cite{allal2025smollm2} reaches 8K by increasing the RoPE base via ABF. It is trained with roughly 40\% long data.
The MiMo-7B~\cite{xiaomi2025mimo} follows three ABF stages to $32\mathrm{K}$ (RoPE $\theta:10\mathrm{K}\!\to\!64\mathrm{K}$). The data mix includes about 10\% synthetic math, code, and creative tasks to elicit longer-range behaviors.
The Gemma 3~\cite{team2025gemma} reaches $128\mathrm{K}$ with ABF on global layers plus position interpolation (PI). Global and local layers are interleaved at approximately $1{:}5$ for memory balance.
The Phi-4~\cite{abdin2024phi} reaches $16\mathrm{K}$ via ABF with a high RoPE base ($\sim\!250\mathrm{K}$). Later training up-weights natural and synthetic $>\!4\mathrm{K}$ data.
The SmolLM3~\cite{bakouch2025smollm3} scales $4\mathrm{K}\!\to\!64\mathrm{K}$ via ABF (with $\theta$ up to $1.5\mathrm{M}\!-\!5\mathrm{M}$) and selective NoPE for every 4th layer. Inference-time YaRN then lifts to $128\mathrm{K}$, focusing on math, code, and reasoning workloads. 

\medskip

\noindent\textbf{YaRN}.
DeepSeek-V2~\cite{liu2024deepseek} uses YaRN to scale $4\mathrm{K}\!\to\!128\mathrm{K}$. And the novel MLA and MoE design provide KV and compute efficiency.
DeepSeek-V3~\cite{DeepSeekAI2024DeepSeekV3TR} applies two YaRN phases with 1000 steps to extend the context length from ($4\mathrm{K}\!\to\!32\mathrm{K}\!\to\!128\mathrm{K}$). The MLA{+}MoE structure provides good efficiency, and it has strong long-context results.
Llama-3.1~\cite{Dubey2024TheL3} trains to $32\mathrm{K}$ then uses YaRN (with an adjusted RoPE base) to reach $128\mathrm{K}$. GQA reduces memory traffic for longer sequences.

\medskip

\noindent\textbf{DCA paired with ABF/YaRN at inference}.
The Qwen-2.5 / 2.5-Turbo~\cite{Yang2024Qwen25TR} uses ABF to extend their context to $128\mathrm{K}$, then uses inference-time DCA{+}YaRN to expand the window to $256\mathrm{K}+$. GQA and Minference-style sparse attention reduce KV cost at inference.
The Qwen-2.5-1M~\cite{yang2025qwen2} uses five ABF stages to extend their context to $256\mathrm{K}$ (with $\theta$ up to $10\mathrm{M}$), then uses DCA{+}YaRN push to $\sim\!1\mathrm{M}$ tokens. The data mix includes $\sim\!40\%$ long and synthetic tasks.
The Qwen-3~\cite{yang2025qwen3} uses ABF to support $32\mathrm{K}$ with a heavy long-data ratio ($\sim\!75\%$), then uses DCA{+}YaRN extend to $128\mathrm{K}$. The recipe retains ABF priors while scaling inference reach.

\medskip

\noindent\textbf{LongRoPE}.
The Phi-3 Mini (128K)~\cite{abdin2024phi3} achieves $128\mathrm{K}$ via LongRoPE~\cite{ding2024longrope} frequency remapping. The block-sparse attention also helps to manage memory and latency.
The MiniCPM-4 8B~\cite{team2025minicpm4} uses LongRoPE to support $32\mathrm{K}$ and uses inference-time YaRN to reach $128\mathrm{K}$. Sparse attention (InfLLM v2) further reduces KV pressure.

\medskip

\noindent\textbf{NTK-aware/PI/NoPE-family}.
The Hunyuan-A13B adopts NTK-aware scaling after a fast $4\mathrm{K}\!\to\!8\mathrm{K}$ stage. The temperature scale is $\alpha\!\approx\!50$ at $32\mathrm{K}$ and $\approx\!1000$ at $256\mathrm{K}$ with a MoE backbone.
The Llama-4 Scout~\cite{meta2025llama4} is pretrained at $256\mathrm{K}$ with iRoPE architecture, which has interleaved attention layers with no position embedding (NoPE). It uses inference-time temperature scaling. The reported inference scaling reaches $\sim\!10\mathrm{M}$ tokens.

\medskip

\noindent\textbf{Architecture-dominant}.
The MiniMax-01~\cite{li2025minimax} targets a $4\mathrm{M}$ context via a $7{:}1$ LightningAttn:full attention interleave. Its training upsamples long data and adds a late-stage QA mix to consolidate extremely long-span utility.
The PANGU PRO MoE~\cite{tang2025pangu} follows a three-phase training under $32\mathrm{K}$ on a MoE backbone. Later stages emphasize STEM and code to strengthen long-range reasoning.

\medskip

\noindent\textit{Takeaway.} Across families, a robust recipe pairs progressive ABF with a non-uniform \(h(\cdot)\) and a mild \(\alpha(s)\) schedule (e.g., YaRN), while architecture assists primarily unlock the memory budget required for long windows rather than replacing frequency remapping.

% Put this block immediately after \label{sec:longcontext} and before Background.

\begin{table*}[h]
  \centering
  % \scriptsize
  \caption{Mid-Training uplift summary.}
  \label{tab:annealing uplift summary}
  \begin{tabular}{p{1.5cm}p{7.6cm}p{7.6cm}}
    \toprule
    \textbf{Model} & \textbf{Quantified Metric Uplift} & \textbf{Qualitative Impact / Purpose of Strategy} \\
    \midrule
    InternLM2     & Achieves substantial performance improvements in coding, reasoning, question answering, examinations, and long-context modeling. & \textbf{[Advanced Reasoning \& Adaptability]} Leveraging enriched datasets for coding, reasoning, QA, and exams, the model outperforms predecessors across 6 dimensions and 30 benchmarks, notably in long-context modeling and subjective evaluation. \\
    \midrule
    MiniCPM4      & Achieves comparable performance to Qwen3 using 22\% of tokens (8T vs 36T). MiniCPM4-0.5B surpasses Llama3.2-1B, Gemma3-1B. MiniCPM4-8B outperforms Gemma3-12B, Phi4-14B. & \textbf{[Scalable Efficiency]} Reduces training costs, improves performance \& efficiency, enables predictable scaling, critical for low-bit LLMs. \\
    \midrule
    DeepSeek-v2   & Achieves 88.9\% pass ratio on MiniF2F-test. Solves 47/658 PutnamBench problems. Pass rate can be improved significantly with CoT reasoning. & \textbf{[Advanced Reasoning \& Adaptability]} Overall training significantly enhances formal theorem proving capabilities. \\
    \midrule
    Zamba         & Base model (Phase 1) achieves Llama2 performance with 1T tokens (Llama2 uses 2T). Overall performance approaches leading ~7B transformers. & \textbf{[Advanced Reasoning \& Adaptability]} Improves adaptation, mitigates forgetting, enables fast adaptation from high-quality data, contributes to data efficiency. \\
    \midrule
    MAP-Neo       & Achieves rapid improvements and stabilization of code metrics during the decay phase, and enhances professional response generation. & \textbf{[Reliable Foundations]} Rectifies tokenizer flaws, enhances code generation, targeted remediation of specific deficiencies. \\
    \midrule
    Nemotron4     & Achieves a better learning from the data and a significant improvement in overall model quality. & \textbf{[Reliable Foundations]} Refines model, effectively integrates new data, contributes to overall polish and robustness. \\
    \midrule
    Llama3        & Annealing improved Llama 3 8B on GSM8k by 24.0\% and MATH by 6.4\%, but a negligible uplift for 405B model. & \textbf{[Scalable Efficiency]} Boosts performance for smaller models, ensures training stability, efficient method to judge value of domain-specific datasets. \\
    \midrule
    Gemma2        & Gemma2-9B shows up to 10\% improvement in some benchmarks compared to previous versions. & \textbf{[Scalable Efficiency]} Significantly improves model quality for smaller models, simulates training beyond available tokens, provides richer gradients. \\
    \midrule
    Hunyuan-Large & Achieves substantial improvements in downstream task performance and generalization ability. & \textbf{[Reliable Foundations]} Optimizes alignment pipeline, front-loads instruction-following, shapes model behavior, improves robustness. \\
    \midrule
    Yi-Lightning  & Achieves model capabilities enhancing and context length extending. & \textbf{[Scalable Efficiency]} Enables efficient training, deployment, inference, and high practical efficacy. Builds capabilities progressively from broad to specialized. \\
    \midrule
    DeepSeek-v3   & Achieves a remarkably stable training without irrecoverable loss spikes or rollbacks. & \textbf{[Scalable Efficiency]} Ensures robust and cost-effective training process, prevents costly failures, critical for large-scale training. \\
    \midrule
    OLMo2         & Mid-training improves 7B by 10.6 pts, 13B by 10.3 pts (avg performance). Specific uplifts on Arc Challenge (+7.2/3.3), MMLU (+3.9/4.1), GSM8K (+43.4/37.8), DROP (+20.1/21.1) for 7B and 13B respectively. & \textbf{[Reliable Foundations]} Dramatically improves performance across diverse benchmarks, enhances robustness, reduces experimentation cost. \\
    \midrule
    Phi-4         & Substantially surpasses teacher model on STEM-focused QA. Meets or exceeds Llama-3.1-405B on reasoning. Scores well above its weight-class on competition math. & \textbf{[Advanced Reasoning \& Adaptability]} Enables emergent reasoning capabilities, improves on teacher models, stabilizes learning during context extension. \\
    \midrule
    PanguProMoE   & Annealing phase can still bring about a large improvement in model performance, stabilize convergence, and reduce catastrophic forgetting. & \textbf{[Reliable Foundations]} Ensures robust generalization and effective specialization, systematically acquires diverse skills, prevents knowledge degradation. \\
    \midrule
    SmolLM3       & Achieves SOTA at 3B scale, and competitive performance with 4B models. Linear Blend retains robust CoT and instruction alignment while recovering top-tier performance on 128 K-token tasks. & \textbf{[Scalable Efficiency]} Maximizes performance from limited parameters, ensures robustness across capabilities, enables long-context robustness. \\
    \bottomrule
  \end{tabular}
\end{table*}

\subsection{Key paradigms and insights at a glance}
\label{sec:longcontext-takeaways}
\noindent{Frequency remapping and length/temperature scaling unify most methods.} Long-context schemes are well modeled by two independent controls:
(i) a \emph{frequency remapping} \(h(\theta_j;s)\) that modifies RoPE’s base angles in each rotary sub-plane (cf.\ Eq.~\eqref{eq:rope-extrapol-h}); and
(ii) an optional \emph{length/temperature scaling} \(\alpha(s)\) applied to queries/keys or softmax (e.g., YaRN in Eq.~\eqref{eq:yarn-alpha}).
PI, NTK-aware, NTK-by-parts, YaRN, and LongRoPE are all specific choices of \(h(\cdot)\) (and sometimes \(\alpha(\cdot)\)).

\noindent\textbf{Trade off between global reach and local fidelity.} Pure interpolation (e.g., PI, Eq.~\eqref{eq:h-pi}) stretches all frequencies equally but can blur local cues at long ranges; NTK-aware~(Eq.~\eqref{eq:h-ntk}) and NTK-by-parts~(Eq.~\eqref{eq:h-ntp}) spare high-frequency bands to preserve short-range structure; YaRN adds a gentle temperature schedule to stabilize optimization at large \(s\).

% \textbf{Piecewise or learned non-uniformity is robust.} Methods that vary scaling across \emph{dimensions} (NTK-by-parts) and/or \emph{positions} (LongRoPE, Eq.~\eqref{eq:h-longrope}) better match the pretrained geometry, yielding stronger length generalization with minimal short-context regression.

\noindent\textbf{Architecture assists matter but are orthogonal.} Interleaved global/local attention, GQA/MLA, and sparse attention primarily reduce KV-cache and memory cost; they \emph{enable} long sequences and pair well with frequency remapping, but do not replace it.

\noindent\textbf{Stage the data and schedule.} Effective recipes expand the window progressively (e.g., \(4\text{k}\!\rightarrow\!32\text{k}\!\rightarrow\!128\text{k}\)), anneal the learning rate, and mix long/short tokens (often 30–60\% long) with synthetic tasks (FIM, retrieval, reordering) to teach cross-document dependencies.

% \textbf{Claimed \(\neq\) effective length.} Report both the advertised window and an \emph{effective} range (e.g., RULER wAvg\(_{\text{inc/dec}}\)) where accuracy remains stable; many models degrade well before the claimed maximum if frequency scaling is mismatched.

\noindent\textbf{Practical default.} For \(L_{\text{target}}\!\le\!128\text{k}\): set ABF to \(\theta\!\approx\!10^6\), 
%use NTK-by-parts \(h(\cdot)\) with \((\alpha,\beta)\!\approx\!(1,32)\), and 
apply YaRN \(\alpha(s)=0.1\ln s+1\). For \(\ge\!1\text{M}\): consider staged ABF (\(\theta\) up to \(10^7\!-\!10^8\)) and dimension/position-aware scaling (e.g., LongRoPE), while protecting short-context performance via short-length recovery.

\section{Evaluation}
\label{sec:evaluation}

\subsection{Standard Benchmarks}
\label{subsec:standard-benchmarks}
At the mid-training stage, model evaluation continues to rely predominantly on widely adopted standard benchmarks rather than task-specific or newly curated datasets. These benchmarks cover a broad spectrum of capabilities, including general knowledge and comprehension, reasoning and problem solving, mathematics and scientific knowledge, coding and software engineering, multilingual understanding, and long-context processing. Such coverage ensures comparability with prior work while highlighting progress across diverse domains. In our study, we report results on \textbf{15 representative benchmarks}, summarized in Table~\ref{tab:std-benchmarks}, which collectively reflect the core skill areas emphasized in large-scale LLM evaluation.

% \begin{table*}[h]
%   \centering
%   \caption{Pre-training benchmarks by category.}
%   \label{tab:std-benchmarks}
%   \begin{tabular}{p{2.2cm}p{11cm}}
%     \toprule
%     \textbf{Task Category} & \textbf{Benchmarks} \\
%     \midrule
%     General & MMLU~\cite{hendrycks2021measuringmassivemultitasklanguage} (5-shot),
%         MMLU-Pro~\cite{wang2024mmluprorobustchallengingmultitask} (5-shot, CoT),
%         ARC ~\cite{clark2018thinksolvedquestionanswering} (0-shot), 
%         HellaSwag~\cite{zellers2019hellaswagmachinereallyfinish} (0-shot),
%         BBH~\cite{suzgun2022challengingbigbenchtaskschainofthought} (3-shot, CoT), 
%         WinoGrande~\cite{sakaguchi2019winograndeadversarialwinogradschema} (0-shot),
%         PIQA~\cite{bisk2019piqareasoningphysicalcommonsense} (0-shot) \\[2pt]
%     \midrule
%     Math \& STEM & GPQA~\cite{rein2023gpqagraduatelevelgoogleproofqa} (5-shot, CoT),
%         MATH~\cite{hendrycks2021measuringmathematicalproblemsolving} (4-shot, CoT),
%         GSM8K~\cite{cobbe2021gsm8k} (4-shot, CoT) \\[2pt]
%     \midrule
%     Coding & HumanEval~\cite{chen2021evaluating} (0-shot), 
%         MBPP~\cite{austin2021programsynthesislargelanguage} (3-shot),  
%         MultiPL-E~\cite{cassano2022multiplescalableextensibleapproach} (0-shot; Python, C++, Java, PHP, TypeScript, C\#, Bash, JavaScript) \\[2pt]
%     \midrule
%     Multilingual & MGSM~\cite{shi2022languagemodelsmultilingualchainofthought} (8-shot, CoT),
%         Flores-101~\cite{goyal2021flores101evaluationbenchmarklowresource} (5-shot) \\[2pt]
%     \midrule
%     Long context & LongBench~\cite{bai2024longbenchbilingualmultitaskbenchmark}, 
%         Ruler~\cite{hsieh2024rulerwhatsrealcontext} \\[2pt]
%     \bottomrule
%   \end{tabular}
% \end{table*}

% \subsubsection{Experimental setup.}
% For each benchmark, we compute scores for various pre-trained models of comparable sizes.
% Where possible, we recompute numbers with our own pipeline for those models.  
% To ensure a fair comparison, we then select the best score between the score that we computed and the reported number for that model with comparable or more conservative settings.  
% The evaluation results of some popular models are shown in Table \ref{tab:pre-trained model evaluation of general tasks} and \ref{tab:pre-trained model evaluation of Math & STEM, Code, Multilingual, and Long Context tasks}.

% \begin{table*}[h]
%   \centering
%   \caption{General tasks evaluation of pre-trained models.}
%   \label{tab:pre-trained model evaluation of general tasks}
%   \begin{tabular}{lcccccccccc}
%     \toprule
%     & \multicolumn{2}{c}{Architecture} & \multicolumn{8}{c}{General Tasks} \\
%     & Total & Act & MMLU & MMLU-Pro & Arc-e & ARC-c & HellaSwag & BBH & WinoGrande & PIQA \\
%     \hline
%     Nemotron4-15B  & 15B  & 15B  & 64.2 &   -  & 80.9 & 55.5 & 82.4 & 58.7 & 78.0 &   -  \\
%     Gemma-7B       & 7B   & 7B   & 64.3 &   -  & 81.5 & 53.2 & 81.2 & 55.1 & 72.3 & 81.2 \\
%     InternLM2-7B   & 7B   & 7B   & 64.0 &   -  &   -  &   -  &   -  & 55.5 &   -  &   -  \\
%     InternLM2-20B  & 20B  & 20B  & 64.0 &   -  &   -  &   -  &   -  & 72.1 &   -  &   -  \\
%     MiniCPM-1.2B   & 1.2B & 1.2B & 49.6 &   -  & 80.9 & 66.8 & 54.7 & 34.7 &   -  &   -  \\
%     MiniCPM-2.4B   & 2.4B & 2.4B & 53.5 &   -  & 85.4 & 68.0 & 68.3 & 36.9 &   -  &   -  \\
%     Phi3-mini      & 3.8B & 3.8B & 68.8 &   -  & 94.6 & 84.9 & 76.7 & 71.7 & 70.8 & 84.2 \\
%     Phi3-small     & 7B   & 7B   & 75.7 &   -  & 97.0 & 90.7 & 77.0 & 79.7 & 81.5 & 86.9 \\
%     Phi3-medium    & 14B  & 14B  & 78.0 &   -  & 91.6 & 91.0 & 82.4 & 81.4 & 81.5 & 87.9 \\
%     Deepseek-v2    & 236B & 21B  & 78.5 &   -  & 97.6 & 92.4 & 84.2 & 78.9 & 84.9 & 83.7 \\
%     Zamba-7B       & 7B   & 7B   & 57.2 &   -  & 74.5 & 46.5 & 80.2 &   -  & 76.4 & 81.4 \\
%     MAP-Neo-7B     & 7B   & 7B   & 54.1 &   -  &   -  & 68.1 & 70.7 &   -  & 59.8 & 76.6 \\
%     Nemotron4-340B & 340B & 340B & 81.1 &   -  &   -  & 94.2 & 90.5 & 85.4 & 89.5 &   -  \\
%     Llama3-8B      & 8B   & 8B   & 66.7 & 37.1 &   -  & 79.7 &   -  & 64.2 & 75.7 & 85.0 \\
%     Llama3-70B     & 70B  & 70B  & 79.5 & 53.8 &   -  & 92.9 &   -  & 81.6 & 83.5 & 83.8 \\
%     Llama3-405B    & 405B & 405B & 85.2 & 61.6 &   -  & 96.1 &   -  & 85.9 & 82.2 & 85.6 \\
%     Gemma2-2B      & 2B   & 2B   & 52.2 &   -  &   -  & 55.7 & 72.9 & 41.9 & 71.3 & 78.4 \\
%     Gemma2-9B      & 9B   & 9B   & 71.3 &   -  &   -  & 68.4 & 81.9 & 68.2 & 80.6 & 81.7 \\
%     Gemma2-27B     & 27B  & 27B  & 75.2 &   -  &   -  & 71.4 & 86.4 & 74.9 & 83.7 & 83.2 \\
%     Hunyuan-Large  & 389B & 52B  & 88.4 & 60.2 &   -  & 95.0 & 86.8 & 86.3 & 88.7 & 88.3 \\
%     Qwen2.5-14B    & 14B  & 14B  & 79.7 & 51.2 &   -  & 67.3 & 84.3 & 78.2 & 81.0 &   -  \\
%     Qwen2.5-72B    & 72B  & 72B  & 86.1 & 58.1 &   -  & 72.4 & 87.6 & 86.3 & 83.9 &   -  \\
%     Qwen2.5-Plus   & 271B & 37B  & 85.4 & 64.0 &   -  & 70.9 & 89.2 & 85.8 & 85.5 &   -  \\
%     DeepSeek-v3    & 671B & 37B  & 87.2 & 59.8 & 87.6 & 75.2 & 88.9 & 87.5 & 84.9 & 84.7 \\
%     OLMo2-7B       & 7B   & 7B   & 63.7 & 31.0 &   -  & 79.8 & 83.8 &   -  & 77.2 &   -  \\
%     OLMo2-13B      & 13B  & 13B  & 67.5 & 35.1 &   -  & 83.5 & 86.4 &   -  & 81.5 &   -  \\
%     MiniMax-01     & 456B & 46B  & 88.5 & 75.7 &   -  &   -  &   -  &   -  &   -  &   -  \\
%     SmolLM2        & 1.7B & 1.7B &   -  & 19.4 & 60.5 &   -  & 68.7 &   -  & 59.4 & 77.6 \\
%     Gemma3-27B     & 27B  & 27B  & 78.7 & 52.9 &   -  &   -  & 84.1 & 80.0 & 71.9 & 83.3 \\
%     Llama4-Scount  & 109B & 17B  & 78.3 & 56.1 &   -  &   -  & 81.9 & 82.4 & 71.4 & 88.9 \\
%     Qwen3-235B     & 235B & 22B  & 87.8 & 68.2 &   -  &   -  &   -  & 88.9 &   -  &   -  \\
%     MiMo-7B        & 7B   & 7B   & 71.2 & 41.9 & 85.2 & 62.3 & 80.0 & 75.2 & 78.0 & 79.4 \\
%     PanguProMoE    & 72B  & 16B  & 87.4 & 63.5 &   -  &   -  & 93.5 & 81.2 & 75.6 & 89.3 \\
%     BitCPM4        & 1B   & 1B   & 59.2 &   -  &   -  &   -  &   -  & 57.6 &   -  &   -  \\
%     SmolLM3        & 3B   & 3B   &   -  & 19.6 & 65.6 &   -  & 76.2 &   -  &   -  &   -  \\
%     \hline
%   \end{tabular}
% \end{table*}

% \begin{table*}[h]
%   \centering
%   \caption{Math \& STEM, Code, Multilingual, and Long Context tasks evaluation of pre-trained models.}
%   \label{tab:pre-trained model evaluation of Math & STEM, Code, Multilingual, and Long Context tasks}
%   \begin{tabular}{lcccccccccc}
%     \toprule
%     & \multicolumn{2}{c}{Architecture} & \multicolumn{3}{c}{Math \& STEM} & \multicolumn{3}{c}{Code} & \multicolumn{2}{c}{Multilingual} \\
%     & Total & Act & GPQA & MATH & GSM8K & HumanEval & MBPP & MultiPL-E & MGSM & FLORES \\
%     \hline
%     Nemotron4-15B  & 15B  & 15B  &   -  &   -  & 46.0 & 31.6 & 40.6 &   -  &   -  &   -  \\
%     Gemma-7B       & 7B   & 7B   &   -  & 24.3 & 46.4 & 32.3 & 44.4 &   -  &   -  &   -  \\
%     InternLM2-7B   & 7B   & 7B   &   -  & 8.9  & 36.0 & 56.1 & 51.8 &   -  &   -  & 5.1  \\
%     InternLM2-20B  & 20B  & 20B  &   -  & 13.6 & 54.3 & 48.8 & 63.0 &   -  &   -  & 6.0  \\
%     MiniCPM-1.2B   & 1.2B & 1.2B &   -  & 10.6 & 31.8 & 44.5 & 32.8 &   -  &   -  &   -  \\
%     MiniCPM-2.4B   & 2.4B & 2.4B &   -  & 10.2 & 53.8 & 50.0 & 47.3 &   -  &   -  &   -  \\
%     Phi3-mini      & 3.8B & 3.8B & 32.8 & 41.3 & 82.5 & 58.5 & 70.0 &   -  & 47.9 &   -  \\
%     Phi3-small     & 7B   & 7B   & 34.3 & 34.6 & 89.6 & 61.0 & 71.7 &   -  & 58.7 &   -  \\
%     Phi3-medium    & 14B  & 14B  &   -  & 53.1 & 91.0 & 62.2 & 75.2 &   -  & 31.8 &   -  \\
%     Deepseek-v2    & 236B & 21B  &   -  & 43.6 & 79.2 & 48.8 & 66.6 &   -  &   -  &   -  \\
%     Zamba-7B       & 7B   & 7B   &   -  &   -  &   -  &   -  &   -  &   -  &   -  &   -  \\
%     MAP-Neo-7B     & 7B   & 7B   &   -  & 20.7 & 53.7 & 23.8 & 34.9 &   -  &   -  &   -  \\
%     Nemotron4-340B & 340B & 340B &   -  &   -  &   -  & 57.3 &   -  &   -  &   -  &   -  \\
%     Llama3-8B      & 8B   & 8B   & 32.8 & 20.3 & 57.2 & 72.6 & 60.8 &   -  & 68.9 &   -  \\
%     Llama3-70B     & 70B  & 70B  & 46.7 & 41.4 & 83.7 & 80.5 & 75.4 &   -  & 86.9 & 38.0 \\
%     Llama3-405B    & 405B & 405B & 51.1 & 53.8 & 89.0 & 89.0 & 82.3 &   -  & 91.6 &   -  \\
%     Gemma2-2B      & 2B   & 2B   &   -  & 16.0 & 24.3 & 20.1 & 30.2 &   -  &   -  &   -  \\
%     Gemma2-9B      & 9B   & 9B   &   -  & 36.6 & 68.6 & 40.2 & 52.4 &   -  &   -  &   -  \\
%     Gemma2-27B     & 27B  & 27B  &   -  & 42.3 & 74.0 & 51.8 & 62.6 &   -  &   -  &   -  \\
%     Hunyuan-Large  & 389B & 52B  &   -  & 69.8 & 92.8 & 71.4 & 72.6 &   -  &   -  &   -  \\
%     Qwen2.5-14B    & 14B  & 14B  & 32.8 & 55.6 & 90.2 & 56.7 & 76.7 & 53.5 &   -  & 36.2 \\
%     Qwen2.5-72B    & 72B  & 72B  & 45.9 & 62.1 & 91.6 & 59.1 & 84.7 & 60.5 & 76.7 & 39.0 \\
%     Qwen2.5-Plus   & 271B & 37B  & 43.9 & 64.4 & 93.0 & 59.1 & 79.7 & 61.0 &   -  & 40.4 \\
%     DeepSeek-v3    & 671B & 37B  & 41.9 & 62.6 & 87.6 & 65.2 & 75.4 & 63.8 & 82.7 &   -  \\
%     OLMo2-7B       & 7B   & 7B   &   -  &   -  & 67.5 &   -  &   -  &   -  &   -  &   -  \\
%     OLMo2-13B      & 13B  & 13B  &   -  &   -  & 75.1 &   -  &   -  &   -  &   -  &   -  \\
%     MiniMax-01     & 456B & 46B  & 54.4 & 77.4 & 94.8 & 86.9 &   -  &   -  &   -  &   -  \\
%     SmolLM2        & 1.7B & 1.7B &   -  & 11.6 & 31.1 & 22.6 &   -  &   -  &   -  &   -  \\
%     Gemma3-27B     & 27B  & 27B  & 26.3 & 51.8 & 81.2 &   -  & 68.4 & 45.0 & 73.7 &   -  \\
%     Llama4-Scount  & 109B & 17B  & 40.4 & 51.7 & 85.4 &   -  & 68.6 & 47.4 & 79.9 &   -  \\
%     Qwen3-235B     & 235B & 22B  & 47.5 & 71.8 & 94.4 &   -  & 81.4 & 65.9 & 83.5 &   -  \\
%     MiMo-7B        & 7B   & 7B   & 25.8 & 37.4 & 75.2 & 51.8 & 69.2 &   -  &   -  &   -  \\
%     PanguProMoE    & 72B  & 16B  &   -  & 55.9 & 86.5 & 63.7 &   -  &   -  &   -  &   -  \\
%     BitCPM4        & 1B   & 1B   &   -  &   -  & 60.8 & 37.2 & 61.5 &   -  &   -  &   -  \\
%     SmolLM3        & 3B   & 3B   &   -  & 46.1 &   -  &   -  &   -  &   -  &   -  & 56.3 \\
%     \bottomrule
%   \end{tabular}
% \end{table*}


% \subsubsection{Significance estimates.}
% Benchmark scores are estimates of a model’s true performance. These estimates have variance because benchmark sets are finite samples drawn from some underlying distribution. We follow \cite{madaan2024quantifyingvarianceevaluationbenchmarks} and report this variance via 95\% confidence intervals (CIs), assuming that benchmark scores are Gaussian-distributed. While this assumption is imperfect (scores are bounded), preliminary bootstrap experiments suggest CIs (for discrete metrics) are a good approximation:
% \begin{equation}
%   \mathrm{CI}(S) \;=\; 1.96 \times
%   \sqrt{\frac{S \,\bigl(1 - S\bigr)}{N}},
%   \label{eq:ci}
% \end{equation}
% where $S$ is the observed benchmark score (e.g.\ accuracy or EM) and $N$ is the sample size of the benchmark.  
% We omit CIs for benchmark scores that are not simple averages.  
% Because subsampling is not the only source of variation, the CIs in Eq.~\eqref{eq:ci} lower-bound the actual variation in the capability estimate.


\subsection{Mid-Training Uplift}
The mid-training strategy has consistently been shown to enhance LLM performance, making it a central component of modern training practice. In this section, we summarize reported gains from a range of prominent models (see Table~\ref{tab:annealing uplift summary}). Our review suggests that the improvements attributed to mid-training fall into several major dimensions: including \textbf{(A)} Advanced Reasoning \& Adaptability, \textbf{(B)} Scalable Efficiency, \textbf{(C)} Reliable Foundations. To facilitate understanding, the table also categorizes these benefits by model, providing a structured comparison of how mid-training contributes across different training regimes.

% \begin{table*}[h]
%   \centering
%   \caption{Mid-Training uplift summary.}
%   \label{tab:annealing uplift summary}
%   \begin{tabular}{p{1.6cm}p{7.6cm}p{8cm}}
%     \toprule
%     \textbf{Model} & \textbf{Quantified Metric Uplift} & \textbf{Qualitative Impact / Purpose of Strategy} \\
%     \midrule
%     InternLM2     & Achieves substantial performance improvements in coding, reasoning, question answering, examinations, and long-context modeling. & \textbf{[Advanced Reasoning \& Adaptability]} Leveraging enriched datasets for coding, reasoning, QA, and exams, the model outperforms predecessors across 6 dimensions and 30 benchmarks, notably in long-context modeling and subjective evaluation. \\
%     \midrule
%     MiniCPM4      & Achieves comparable performance to Qwen3 using 22\% of tokens (8T vs 36T). MiniCPM4-0.5B surpasses Llama3.2-1B, Gemma3-1B. MiniCPM4-8B outperforms Gemma3-12B, Phi4-14B. & \textbf{[Scalable Efficiency]} Reduces training costs, improves performance \& efficiency, enables predictable scaling, critical for low-bit LLMs. \\
%     \midrule
%     DeepSeek-v2   & Achieves 88.9\% pass ratio on MiniF2F-test. Solves 47/658 PutnamBench problems. Pass rate can be improved significantly with CoT reasoning. & \textbf{[Advanced Reasoning \& Adaptability]} Overall training significantly enhances formal theorem proving capabilities. \\
%     \midrule
%     Zamba         & Base model (Phase 1) achieves Llama2 performance with 1T tokens (Llama2 uses 2T). Overall performance approaches leading ~7B transformers. & \textbf{[Advanced Reasoning \& Adaptability]} Improves adaptation, mitigates forgetting, enables fast adaptation from high-quality data, contributes to data efficiency. \\
%     \midrule
%     MAP-Neo       & Achieves rapid improvements and stabilization of code metrics during the decay phase, and enhances professional response generation. & \textbf{[Reliable Foundations]} Rectifies tokenizer flaws, enhances code generation, targeted remediation of specific deficiencies. \\
%     \midrule
%     Nemotron4     & Achieves a better learning from the data and a significant improvement in overall model quality. & \textbf{[Reliable Foundations]} Refines model, effectively integrates new data, contributes to overall polish and robustness. \\
%     \midrule
%     Llama3        & Annealing improved Llama 3 8B on GSM8k by 24.0\% and MATH by 6.4\%, but a negligible uplift for 405B model. & \textbf{[Scalable Efficiency]} Boosts performance for smaller models, ensures training stability, efficient method to judge value of domain-specific datasets. \\
%     \midrule
%     Gemma2        & Gemma2-9B shows up to 10\% improvement in some benchmarks compared to previous versions. & \textbf{[Scalable Efficiency]} Significantly improves model quality for smaller models, simulates training beyond available tokens, provides richer gradients. \\
%     \midrule
%     Hunyuan-Large & Achieves substantial improvements in downstream task performance and generalization ability. & \textbf{[Reliable Foundations]} Optimizes entire alignment pipeline, front-loads instruction-following, shapes model behavior, improves robustness. \\
%     \midrule
%     Yi-Lightning  & Achieves model capabilities enhancing and context length extending. & \textbf{[Scalable Efficiency]} Enables efficient training, deployment, inference, and high practical efficacy. Builds capabilities progressively from broad to specialized. \\
%     \midrule
%     DeepSeek-v3   & Achieves a remarkably stable training without irrecoverable loss spikes or rollbacks. & \textbf{[Scalable Efficiency]} Ensures robust and cost-effective training process, prevents costly failures, critical for large-scale training. \\
%     \midrule
%     OLMo2         & Mid-training improves 7B by 10.6 pts, 13B by 10.3 pts (avg performance). Specific uplifts on Arc Challenge (+7.2/3.3), MMLU (+3.9/4.1), GSM8K (+43.4/37.8), DROP (+20.1/21.1) for 7B and 13B respectively. & \textbf{[Reliable Foundations]} Dramatically improves performance across diverse benchmarks, enhances robustness, reduces experimentation cost. \\
%     \midrule
%     Phi-4         & Substantially surpasses teacher model on STEM-focused QA. Meets or exceeds Llama-3.1-405B on reasoning. Scores well above its weight-class on competition math. & \textbf{[Advanced Reasoning \& Adaptability]} Enables emergent reasoning capabilities, improves on teacher models, stabilizes learning during context extension. \\
%     \midrule
%     PanguProMoE   & Annealing phase can still bring about a large improvement in model performance, stabilize convergence, and reduce catastrophic forgetting. & \textbf{[Reliable Foundations]} Ensures robust generalization and effective specialization, systematically acquires diverse skills, prevents knowledge degradation. \\
%     \midrule
%     SmolLM3       & Achieves SOTA at 3B scale, and competitive performance with 4B models. Linear Blend retains robust chain-of-thought and instruction alignment while recovering top-tier performance on 128 K-token tasks. & \textbf{[Scalable Efficiency]} Maximizes performance from limited parameters, ensures robustness across capabilities, enables long-context robustness. \\
%     \bottomrule
%   \end{tabular}
% \end{table*}

% The following part shows each model is grouped under a broad category (e.g., Efficiency, Reasoning, Alignment), with a specific focus tag that highlights its main innovation or improvement. 
% \begin{itemize}
%   \item \textbf{Efficiency}
%   \begin{itemize}
%     \item MiniCPM4: [Efficiency: Low-Bit Scaling \& Cost Reduction]
%     \item Yi-Lightning: [Efficiency: Practical Training \& Deployment]
%     \item DeepSeek-v3: [Efficiency: Cost-Efficient Large-Scale Training]
%     \item SmolLM3: [Efficiency: Parameter-Efficient Robustness \& Long-Context]
%   \end{itemize}

%   \item \textbf{Capabilities \& Reasoning}
%   \begin{itemize}
%     \item InternLM2: [Capabilities: Broad Capability Uplift \& Long-Context]
%     \item DeepSeek-v2: [Reasoning: Formal Theorem Proving]
%     \item Phi-4: [Reasoning: Emergent Reasoning \& Context Stabilization]
%   \end{itemize}

%   \item \textbf{Adaptation \& Data Handling}
%   \begin{itemize}
%     \item Zamba: [Adaptation: Data Efficiency \& Forgetting Mitigation]
%     \item MAP-Neod: [Infrastructure: Tokenizer Remediation \& Code Generation]
%   \end{itemize}

%   \item \textbf{Refinement \& Robustness}
%   \begin{itemize}
%     \item Nemotron4: [Refinement: Data Integration \& Robustness]
%     \item OLMo2: [Robustness: Benchmark Strength \& Experiment Cost Reduction]
%   \end{itemize}

%   \item \textbf{Scaling \& Small-Model Performance}
%   \begin{itemize}
%     \item Llama3: [Scaling: Small-Model Stability \& Dataset Evaluation]
%     \item Gemma2: [Scaling: Small-Model Quality \& Gradient Enrichment]
%   \end{itemize}

%   \item \textbf{Alignment \& Behavior Shaping}
%   \begin{itemize}
%     \item Hunyuan-Large: [Alignment: Instruction Optimization \& Robustness]
%   \end{itemize}

%   \item \textbf{Architecture: Mixture-of-Experts}
%   \begin{itemize}
%     \item PanguProMoE: [Architecture: MoE Generalization \& Specialization]
%   \end{itemize}
% \end{itemize}


\section{Future Work}
% \label{sec:related}

% Research related to \emph{annealing} in large language model (LLM) training is scattered across several neighboring areas—continual pre-training, domain-adaptive training, learning-rate schedule theory, and long-context extension. 
% %In this section, we synthesize these strands and clarify how our survey differs from prior overviews.

% \subsection{Continual Pre-training and Continual Learning}
% The body of work most closely aligned with what many practitioners now call \emph{mid-training} or \emph{annealing} is surveyed under the umbrella of \emph{continual learning for LLMs}. Wu et al.\ (2024) provide a taxonomy that explicitly names ``continual pre-training'' as the first stage of continual learning, mapping almost one-to-one to industrial mid-training practice. Their survey covers objectives, data curation, catastrophic forgetting, and the design space of parameter-efficient vs.\ full-model updates. A companion, continuously updated GitHub survey tracks roughly 200 papers on domain refresh, ReLoRA, and soft masking strategies. %:contentReference[oaicite:0]{index=0}

% \subsection{Domain-Adaptive Pre-training and General PLM Surveys}
% Earlier but still relevant work on \emph{domain-adaptive pre-training} (DAPT) for BERT/GPT-style models synthesizes principles of data selection, regularization, and continual objectives that have since been scaled to 7--70B LLMs. Guo and Yu (2022) give a thorough taxonomy of DAPT methods, while Chen et al.\ (2024) include a concise chapter on training strategies, curriculum/annealing schedules, and specialized continued pre-training within a broad LLM survey. %:contentReference[oaicite:1]{index=1}

% \subsection{Learning-Rate Annealing and Optimization Schedules}
% Despite its ubiquity, learning-rate decay itself has not been the focus of many surveys. Tissue et al.\ (2024) offer one of the first systematic analyses of cosine/linear annealing for LLMs, proposing predictive scaling laws that link schedules to loss trajectories. %:contentReference[oaicite:2]{index=2} 
% Recent meta-analyses emphasize that ``annealing per se lacks a dedicated survey'' and remains treated as a hyperparameter detail in most recipe papers. %:contentReference[oaicite:3]{index=3} 
% Complementary works (e.g., Defazio et al.\ 2023; Bergsma et al.\ 2025) study optimal linear decay and the empirical success of straight-to-zero schedules, enriching the theoretical toolkit for annealing stages. %:contentReference[oaicite:4]{index=4}

% \subsection{Long-Context Extension Surveys}
% Because many annealing passes coincide with context-length extension, surveys on long-context LLMs are directly relevant. Liu et al.\ (2025) present the most up-to-date LLM-centric synthesis, organizing 250+ papers by data strategies, RoPE remapping variants, retrieval/memory mechanisms, and evaluation. Wang et al.\ (IJCAI'24) catalog practical tricks for extending windows (positional encoding, chunking, hierarchical processing), while Zhao et al.\ (EMNLP'24) give a deep dive into positional-encoding math and length extrapolation protocols. %:contentReference[oaicite:5]{index=5}

% \subsection{Efficient Attention and State-Space Alternatives}
% Historical surveys on ``efficient transformers'' (e.g., Tay et al.\ 2020) provide the foundational taxonomy of local/sparse/low-rank attention and memory mechanisms that long-context methods draw upon. %:contentReference[oaicite:6]{index=6} 
% More recent state-space model (SSM) surveys (Tiezzi 2024; Patro 2024) catalog Mamba, S4, Hyena, and RetNet as transformer alternatives that sidestep quadratic attention, complementing the transformer-centric context-extension literature. %:contentReference[oaicite:7]{index=7}

% \subsection{Overtraining, Scaling Laws, and Stability Analyses}
% A line of work investigates the risks and dynamics of \emph{overtraining}—a phenomenon particularly relevant to late-stage annealing. Springer et al.\ (2025) show that overtrained LMs can be harder to fine-tune, emphasizing the need for carefully budgeted annealing passes. %:contentReference[oaicite:8]{index=8} 
% Broader scaling-law research (Kaplan et al.\ 2020; Hoffmann et al.\ 2022; Li et al.\ 2025) and curvature-based stability analyses (Gilmer et al.\ 2022) provide theoretical baselines for compute-optimality and optimization stability that intersect with learning-rate and data-annealing design. %:contentReference[oaicite:9]{index=9}

% \subsection{Position of This Survey}
% To our knowledge, there is no prior survey that \emph{unifies} (i) high-quality data curation, (ii) learning-rate scheduler theory and practice, and (iii) long-context extension into a \emph{single late-stage training paradigm}. Previous surveys cover one or two axes in isolation. Our work integrates all three, adding a design-oriented perspective (recipes, ablations, and open challenges) tailored for practitioners executing annealing under real-world constraints. %:contentReference[oaicite:10]{index=10}

% \subsection{Open Challenges and Future Directions}
\label{sec:future}
Despite its growing adoption, \emph{mid-training} for LLMs remains more of an engineering art than a principled science. Below we outline key open challenges and sketch promising directions for future research.


\textbf{Dynamic Curriculum Design that adaptively adjust data composition.} 
A promising direction for mid-training is the development of dynamic curriculum strategies that adaptively adjust data composition based on the model’s evolving capabilities. Rather than relying on fixed sampling ratios, future approaches could monitor intermediate performance signals (e.g., loss sharpness, gradient variance, or benchmark outcomes) to guide the gradual transition from broad natural corpora toward specialized reasoning, coding, or multilingual tasks. Such adaptive curriculum would allow mid-training to better align data exposure with the model’s learning trajectory, improving efficiency while mitigating overfitting.

{\textbf{Theoretically grounded schedulers with adaptive mechanisms that co-evolve with scale and data.}}
{Learning rate schedulers remain a critical yet underexplored component in LLM training. Their interaction with other experimental parameters, such as batch size and model size, continues to be an active area of empirical and theoretical research. Future work should aim to unify these disparate empirical findings under a theoretical framework, and develop adaptive or scale-aware schedulers that generalize across model sizes and training regimes.}

\textbf{Long-context progress hinges on two fronts: fixing position-embedding OOD and cutting compute/latency.} Beyond RoPE tweaks, push toward general, learnable position functions and index warps, with robustness curricula and on-the-fly calibration. Beyond sparse attention, attack efficiency via KV-cache compression and quantization, token pruning/merging, hybrid attention-SSM stacks, retrieval-summarized contexts, and length-aware scheduling with compute controllers and MoE specialized by sequence length. We should also come up with better long-range datasets and metrics, and theory that ties spectral generalization and compute–length scaling laws—validated by concrete experiments on learned PEs, hybrid stacks, KV compression, retrieval summaries, and OOD calibration.
% \section{Conclusion}
% \label{sec:conclusion}
% Summarise consensus: small, high-quality, late passes punch above compute weight.
% Highlight emerging trend: “mid-training” as an industry standard pipeline step.
% Stress need for reproducible ablations and open benchmarks.

\section{Conclusion}
\label{sec:conclusion}

% This survey reframed \emph{mid-training} for large language models. We argued that effective mid-training requires \emph{jointly} manipulating three levers: (i) \emph{data distribution refinement} to concentrate scarce, capability-dense tokens where they matter most; (ii) \emph{optimization cooling} via carefully designed learning-rate (and optimizer) schedules that stabilize convergence; and (iii) \emph{context-length extension} strategies that expand an LLM’s usable window without sacrificing stability or perplexity. By treating these components as interdependent rather than isolated tweaks, practitioners can unlock significant capability gains under a modest compute budget.

% We synthesized empirical recipes and theoretical intuitions behind data valuation, curriculum-like token allocation, and schedule design, and we unified recent long-context techniques under a frequency-remapping view of positional encodings. Alongside practical guidelines, we highlighted gaps: the lack of formal objectives for mid-training, limited understanding of why certain recipes generalize, and the absence of standardized, contamination-aware benchmarks to measure incremental gains.

% % \paragraph{Takeaway.} 
% Mid-training is no longer a footnote in training logs; it is a strategic phase where careful engineering meets principled design. We hope this survey provides both a conceptual scaffold and a practical playbook, accelerating the transition of mid-training from ad-hoc recipes to a disciplined methodology for building more capable, stable, and compute-efficient LLMs.

This paper presents a first-of-its-kind survey of LLM mid-training approaches. We introduce a taxonomy that categorizes existing methods into three key domains: data distribution, learning rate scheduling, and long-context extension. We summarize the main insights in each domain, providing a structured reference for researchers and practitioners. We also compile common evaluation benchmarks and reported gains, enabling a comparative view of how mid-training improves model performance. Finally, we identify open challenges and propose future research avenues, positioning mid-training as a central stage for shaping the next generation of large language models.
% \bibliographystyle{IEEEtran}
% \bibliography{main}
\input{main.bbl}
% \clearpage
% \appendix
% \input{CompassLLM/content/8_appendix}
\end{document}
